% Generated mechanically from the frozen spaces-divisors.tex target.
% Do not edit this generated file; edit the bound locale source instead.

% OK, start here.
%

\setcounter{chapter}{70}
\chapter{代数空间上的除子}



\phantomsection
\label{spaces-divisors-section-phantom}


\section{引言}
\label{spaces-divisors-section-introduction}

\noindent
本章研究代数空间上的除子及相关主题。关于代数空间的基本参考文献是
\cite{Kn}。









\section{伴随点与弱伴随点}
\label{spaces-divisors-section-associated}

\noindent
在概形的情形，我们已经引入了两个相互竞争的伴随点概念，即通常的伴随点
（《除子》第 \ref{divisors-section-associated} 节）和弱伴随点
（《除子》第 \ref{divisors-section-weakly-associated} 节）。对于一般代数空间，
伴随点的概念基本无用，我们甚至不打算引入它。若代数空间局部诺特，则允许
用“伴随点”一词代替“弱伴随点”，因为对诺特概形两种概念相同
（《除子》引理 \ref{divisors-lemma-ass-weakly-ass}）。在给出定义之前，
先需要一个引理。

\begin{lemma}
\label{spaces-divisors-lemma-associated}
设 $S$ 为概形，$X$ 为 $S$ 上的代数空间，$\mathcal{F}$ 为拟凝聚
$\mathcal{O}_X$-模，且 $x \in |X|$。下列条件等价：
\begin{enumerate}
\item 对某个平展态射 $f : U \to X$（其中 $U$ 是概形）以及映到 $x$ 的
$u \in U$，点 $u$ 弱伴随于 $f^*\mathcal{F}$；
\item 对每个平展态射 $f : U \to X$（其中 $U$ 是概形）以及映到 $x$ 的
$u \in U$，点 $u$ 弱伴随于 $f^*\mathcal{F}$；
\item $\mathcal{O}_{X, \overline{x}}$ 的极大理想是茎
$\mathcal{F}_{\overline{x}}$ 的弱伴随素理想。
\end{enumerate}
若 $X$ 局部诺特，则这些条件还等价于：
\begin{enumerate}
\item[(4)] 对某个平展态射 $f : U \to X$（其中 $U$ 是概形）以及映到 $x$
的 $u \in U$，点 $u$ 伴随于 $f^*\mathcal{F}$；
\item[(5)] 对每个平展态射 $f : U \to X$（其中 $U$ 是概形）以及映到 $x$
的 $u \in U$，点 $u$ 伴随于 $f^*\mathcal{F}$；
\item[(6)] $\mathcal{O}_{X, \overline{x}}$ 的极大理想是茎
$\mathcal{F}_{\overline{x}}$ 的伴随素理想。
\end{enumerate}
\end{lemma}

\begin{proof}
选择带有点 $u$ 的概形 $U$ 以及把 $u$ 映到 $x$ 的平展态射
$f : U \to X$。把 $\overline{x}$ 提升为 $U$ 上位于 $u$ 上方的几何点。
回忆
$\mathcal{O}_{X, \overline{x}} = \mathcal{O}_{U, u}^{sh}$
，其中严格亨泽尔化是相对于所选的 $\overline{x}$ 的提升而言；参见
《空间的性质》引理
\ref{spaces-properties-lemma-describe-etale-local-ring}。最后，有
$$
\mathcal{F}_{\overline{x}} =
(f^*\mathcal{F})_u \otimes_{\mathcal{O}_{U, u}}
\mathcal{O}_{X, \overline{x}} =
(f^*\mathcal{F})_u \otimes_{\mathcal{O}_{U, u}}
\mathcal{O}_{U, u}^{sh}
$$
其中等式由《空间的性质》引理
\ref{spaces-properties-lemma-stalk-quasi-coherent} 得到。
因此，(1)、(2) 与 (3) 的等价性来自《平坦性的更多内容》引理
\ref{flat-lemma-weakly-associated-henselization}。若 $X$ 局部诺特，
则上述任意 $U$ 都局部诺特，故由《除子》引理
\ref{divisors-lemma-ass-weakly-ass}，(1)（相应地 (2)）等价于 (4)
（相应地 (5)）。另一方面，在局部诺特情形，局部环
$\mathcal{O}_{X, \overline{x}}$ 也诺特（《空间的性质》引理
\ref{spaces-properties-lemma-Noetherian-local-ring-Noetherian}）。
% P09-ERR-0147
因此，由同一引理（或由《代数》引理
\ref{algebra-lemma-ass-weakly-ass}）得到 (3) 与 (6) 的等价性。
\end{proof}

\begin{definition}
\label{spaces-divisors-definition-weakly-associated}
设 $S$ 为概形，$X$ 为 $S$ 上的代数空间，$\mathcal{F}$ 为 $X$ 上的
拟凝聚层，且 $x \in |X|$。
\begin{enumerate}
\item 若引理 \ref{spaces-divisors-lemma-associated} 的等价条件 (1)、(2)、(3) 成立，
则称 $x$ {\it 弱伴随于} $\mathcal{F}$。
\item % P09-ERR-0148
用 $\text{WeakAss}(\mathcal{F})$ 表示 $\mathcal{F}$ 的弱伴随点集合。
\item $X$ 的{\it 弱伴随点}是 $\mathcal{O}_X$ 的弱伴随点。
\end{enumerate}
若 $X$ 局部诺特，则当且仅当 $x$ 弱伴随于 $\mathcal{F}$ 时，称
{\it $x$ 伴随于 $\mathcal{F}$}，并置
$\text{Ass}(\mathcal{F}) = \text{WeakAss}(\mathcal{F})$。最后
（仍假设 $X$ 局部诺特），当且仅当 $x$ 是 $X$ 的弱伴随点时，称
{\it $x$ 是 $X$ 的伴随点}。
\end{definition}

\noindent
现在可以证明一些必备引理。

\begin{lemma}
\label{spaces-divisors-lemma-weakly-ass-support}
设 $S$ 为概形，$X$ 为 $S$ 上的代数空间，$\mathcal{F}$ 为拟凝聚
$\mathcal{O}_X$-模。则
$\text{WeakAss}(\mathcal{F}) \subset \text{Supp}(\mathcal{F})$。
\end{lemma}

\begin{proof}
这由定义立即得到。$X$ 上阿贝尔层的支撑定义见《空间的性质》定义
\ref{spaces-properties-definition-support}。
\end{proof}

\begin{lemma}
\label{spaces-divisors-lemma-ses-weakly-ass}
设 $S$ 为概形，$X$ 为 $S$ 上的代数空间。设
$0 \to \mathcal{F}_1 \to \mathcal{F}_2 \to \mathcal{F}_3 \to 0$
为 $X$ 上拟凝聚层的短正合列。则
$\text{WeakAss}(\mathcal{F}_2) \subset
\text{WeakAss}(\mathcal{F}_1) \cup \text{WeakAss}(\mathcal{F}_3)$
且
$\text{WeakAss}(\mathcal{F}_1) \subset \text{WeakAss}(\mathcal{F}_2)$。
\end{lemma}

\begin{proof}
对每个几何点 $\overline{x} \in X$，茎的序列
$0 \to \mathcal{F}_{1, \overline{x}} \to
\mathcal{F}_{2, \overline{x}} \to
\mathcal{F}_{3, \overline{x}} \to 0$
是 $\mathcal{O}_{X, \overline{x}}$-模的短正合列。因此本引理由
《代数》引理 \ref{algebra-lemma-weakly-ass} 得出。
\end{proof}

\begin{lemma}
\label{spaces-divisors-lemma-weakly-ass-zero}
设 $S$ 为概形，$X$ 为 $S$ 上的代数空间，$\mathcal{F}$ 为拟凝聚
$\mathcal{O}_X$-模。则
$$
\mathcal{F} = (0) \Leftrightarrow \text{WeakAss}(\mathcal{F}) = \emptyset
$$
\end{lemma}

\begin{proof}
取概形 $U$ 以及满平展态射 $f : U \to X$。则 $\mathcal{F}$ 为零当且仅当
$f^*\mathcal{F}$ 为零。因此，本引理由定义以及概形情形的对应引理得出；参见
《除子》引理 \ref{divisors-lemma-weakly-ass-zero}。
\end{proof}

\begin{lemma}
\label{spaces-divisors-lemma-minimal-support-in-weakly-ass}
设 $S$ 为概形，$X$ 为 $S$ 上的代数空间，$\mathcal{F}$ 为拟凝聚
$\mathcal{O}_X$-模，且 $x \in |X|$。若
\begin{enumerate}
\item $x \in \text{Supp}(\mathcal{F})$；
\item $x$ 是 $X$ 的余维 $0$ 点（《空间的性质》定义
\ref{spaces-properties-definition-dimension-local-ring}）。
\end{enumerate}
则 $x \in \text{WeakAss}(\mathcal{F})$。若 $\mathcal{F}$ 是有限型
$\mathcal{O}_X$-模，其概形论支撑为 $Z$（《空间的态射》定义
\ref{spaces-morphisms-definition-scheme-theoretic-support}），且 $x$ 是 $Z$
的余维 $0$ 点，则 $x \in \text{WeakAss}(\mathcal{F})$。
\end{lemma}

\begin{proof}
由于 $x \in \text{Supp}(\mathcal{F})$，茎
$\mathcal{F}_{\overline{x}}$ 非零。因此，由《代数》引理
\ref{algebra-lemma-weakly-ass-zero}，
$\text{WeakAss}(\mathcal{F}_{\overline{x}})$ 非空。另一方面，
$\mathcal{O}_{X, \overline{x}}$ 的谱仅含一个点。因此按定义，$x$ 是
$\mathcal{F}$ 的弱伴随点。最后一个断言成立是因为
% P09-ERR-0149
$\mathcal{O}_{X, \overline{x}} \to \mathcal{O}_{Z, \overline{z}}$
是满射，$\mathcal{O}_{Z, \overline{z}}$ 的谱仅含一个点，并且
$\mathcal{F}_{\overline{x}}$ 是非零的 $\mathcal{O}_{Z, \overline{z}}$-模。
\end{proof}

\begin{lemma}
\label{spaces-divisors-lemma-minimal-support-in-weakly-ass-decent}
设 $S$ 为概形，$X$ 为 $S$ 上的代数空间，$\mathcal{F}$ 为拟凝聚
$\mathcal{O}_X$-模，且 $x \in |X|$。若
\begin{enumerate}
\item $X$ 是良态的（例如拟分离或局部分离）；
\item $x \in \text{Supp}(\mathcal{F})$；
\item $x$ 不是 $\text{Supp}(\mathcal{F})$ 中另一个点的特化。
\end{enumerate}
则 $x \in \text{WeakAss}(\mathcal{F})$。
\end{lemma}

\begin{proof}
（拟分离代数空间是良态的，参见《良态空间》第
\ref{decent-spaces-section-reasonable-decent} 节；局部分离代数空间是良态的，
参见《良态空间》引理 \ref{decent-spaces-lemma-locally-separated-decent}。）
取概形 $U$、点 $u \in U$ 以及把 $u$ 映到 $x$ 的平展态射
$f : U \to X$。由《良态空间》引理
\ref{decent-spaces-lemma-decent-no-specializations-map-to-same-point}，
若 $u' \leadsto u$ 是非平凡特化，则 $f(u') \not = x$。由此可见，
$u \in \text{Supp}(f^*\mathcal{F})$ 不是
$\text{Supp}(f^*\mathcal{F})$ 中另一个点的特化。因此，由
% P09-ERR-0150
《除子》引理 \ref{spaces-divisors-lemma-minimal-support-in-weakly-ass}，
$u \in \text{WeakAss}(f^*\mathcal{F})$。
\end{proof}

\begin{lemma}
\label{spaces-divisors-lemma-finite-ass}
设 $S$ 为概形，$X$ 为 $S$ 上的局部诺特代数空间，$\mathcal{F}$ 为凝聚
$\mathcal{O}_X$-模。则对每个拟紧开子集 $W \subset |X|$，
$\text{Ass}(\mathcal{F}) \cap W$ 是有限集。
\end{lemma}

\begin{proof}
取拟紧概形 $U$ 以及平展态射 $U \to X$，使得 $W$ 是
$|U| \to |X|$ 的像。于是 $U$ 是诺特概形，应用《除子》引理
\ref{divisors-lemma-finite-ass} 即得结论。
\end{proof}

\begin{lemma}
\label{spaces-divisors-lemma-restriction-injective-open-contains-weakly-ass}
设 $S$ 为概形，$X$ 为 $S$ 上的代数空间，$\mathcal{F}$ 为拟凝聚
$\mathcal{O}_X$-模。若 $U \to X$ 是平展态射且
$\text{WeakAss}(\mathcal{F}) \subset \Im(|U| \to |X|)$，则
$\Gamma(X, \mathcal{F}) \to \Gamma(U, \mathcal{F})$ 是单射。
\end{lemma}

\begin{proof}
设 $s \in \Gamma(X, \mathcal{F})$ 是限制到 $U$ 上为零的截面。令
$\mathcal{F}' \subset \mathcal{F}$ 为 $s$ 所定义的映射
$\mathcal{O}_X \to \mathcal{F}$ 的像。则 $\mathcal{F}'|_U = 0$。这蕴含
$\text{WeakAss}(\mathcal{F}') \cap \Im(|U| \to |X|) = \emptyset$
（按弱伴随点的定义）。另一方面，由引理 \ref{spaces-divisors-lemma-ses-weakly-ass}，
$\text{WeakAss}(\mathcal{F}') \subset \text{WeakAss}(\mathcal{F})$
。故 $\text{WeakAss}(\mathcal{F}') = \emptyset$。于是由引理
\ref{spaces-divisors-lemma-weakly-ass-zero}，$\mathcal{F}' = 0$。
\end{proof}

\begin{lemma}
\label{spaces-divisors-lemma-weakass-pushforward}
设 $S$ 为概形，$f : X \to Y$ 为 $S$ 上代数空间间的拟紧拟分离态射，
$\mathcal{F}$ 为拟凝聚 $\mathcal{O}_X$-模。设 $y \in |Y|$ 不在
$|f|$ 的像中。则 $y$ 不弱伴随于 $f_*\mathcal{F}$。
\end{lemma}

\begin{proof}
由《空间的态射》引理 \ref{spaces-morphisms-lemma-pushforward}，
$\mathcal{O}_Y$-模 $f_*\mathcal{F}$ 拟凝聚，故本引理的陈述有意义。
取仿射概形 $V$、点 $v \in V$ 以及把 $v$ 映到 $y$ 的平展态射
$V \to Y$。可将 $f : X \to Y$、$\mathcal{F}$、$y$ 分别替换为
$X \times_Y V \to V$、$\mathcal{F}|_{X \times_Y V}$、$v$。
因此可设 $Y$ 为仿射概形。此时 $X$ 拟紧，故可取仿射概形 $U$ 以及满平展态射
$U \to X$。
% P09-ERR-0151
用 $g : U \to Y$ 表示合成态射。于是
$f_*\mathcal{F} \subset g_*(\mathcal{F}|_U)$。
由引理 \ref{spaces-divisors-lemma-ses-weakly-ass}，可归约到概形的情形，即《除子》引理
\ref{divisors-lemma-weakass-pushforward}。
\end{proof}

\begin{lemma}
\label{spaces-divisors-lemma-check-injective-on-weakass}
设 $S$ 为概形，$X$ 为 $S$ 上的代数空间。设
$\varphi : \mathcal{F} \to \mathcal{G}$ 为拟凝聚 $\mathcal{O}_X$-模间的
映射。假设对每个 $x \in |X|$，下列情形至少有一个发生：
\begin{enumerate}
\item $\mathcal{F}_{\overline{x}} \to \mathcal{G}_{\overline{x}}$
是单射；或
\item $x \not \in \text{WeakAss}(\mathcal{F})$。
\end{enumerate}
则 $\varphi$ 是单射。
\end{lemma}

\begin{proof}
这些假设蕴含 $\text{WeakAss}(\Ker(\varphi)) = \emptyset$，故由引理
\ref{spaces-divisors-lemma-weakly-ass-zero}，$\Ker(\varphi) = 0$。
\end{proof}

\begin{lemma}
\label{spaces-divisors-lemma-weakass-reduced}
设 $S$ 为概形，$X$ 为 $S$ 上的既约代数空间。则
% P09-ERR-0152
$X$ 的弱伴随点恰为 $X$ 的余维 $0$ 点。
\end{lemma}

\begin{proof}
在平展局部工作，本引理由《除子》引理
\ref{divisors-lemma-weakass-reduced} 及《空间的性质》引理
\ref{spaces-properties-lemma-codimension-0-points} 得出。
\end{proof}










\section{态射与弱伴随点}
\label{spaces-divisors-section-morphisms-weakly-associated}

\begin{lemma}
\label{spaces-divisors-lemma-weakly-ass-reverse-functorial}
设 $S$ 为概形，$f : X \to Y$ 为 $S$ 上代数空间间的仿射态射，
$\mathcal{F}$ 为拟凝聚 $\mathcal{O}_X$-模。则有
% P09-ERR-0153
$$
\text{WeakAss}_S(f_*\mathcal{F}) \subset f(\text{WeakAss}_X(\mathcal{F}))
$$
\end{lemma}

\begin{proof}
取概形 $V$ 以及满平展态射 $V \to Y$。置 $U = X \times_Y V$。则
$U \to V$ 是概形间的仿射态射。由弱伴随点的定义，问题归约为概形间的态射
$U \to V$；这种情形已由《除子》引理
\ref{divisors-lemma-weakly-ass-reverse-functorial} 处理。
\end{proof}

\begin{lemma}
\label{spaces-divisors-lemma-ass-functorial-equal}
设 $S$ 为概形，$f : X \to Y$ 为 $S$ 上代数空间间的仿射态射，
$\mathcal{F}$ 为拟凝聚 $\mathcal{O}_X$-模。若 $X$ 局部诺特，则有
$$
\text{WeakAss}_Y(f_*\mathcal{F}) =
f(\text{WeakAss}_X(\mathcal{F}))
$$
\end{lemma}

\begin{proof}
取概形 $V$ 以及满平展态射 $V \to Y$。置 $U = X \times_Y V$。则
$U \to V$ 是概形间的仿射态射，且 $U$ 局部诺特。由弱伴随点的定义，
问题归约为概形间的态射 $U \to V$；这种情形已由《除子》引理
\ref{divisors-lemma-ass-functorial-equal} 处理。
\end{proof}

\begin{lemma}
\label{spaces-divisors-lemma-weakly-associated-finite}
设 $S$ 为概形，$f : X \to Y$ 为 $S$ 上代数空间间的有限态射，
$\mathcal{F}$ 为拟凝聚 $\mathcal{O}_X$-模。则
$\text{WeakAss}(f_*\mathcal{F}) = f(\text{WeakAss}(\mathcal{F}))$。
\end{lemma}

\begin{proof}
取概形 $V$ 以及满平展态射 $V \to Y$。置 $U = X \times_Y V$。则
$U \to V$ 是概形间的有限态射。由弱伴随点的定义，问题归约为概形间的态射
$U \to V$；这种情形已由《除子》引理
\ref{divisors-lemma-weakly-associated-finite} 处理。
\end{proof}

\begin{lemma}
\label{spaces-divisors-lemma-weakly-ass-pullback}
设 $S$ 为概形，$f : X \to Y$ 为 $S$ 上代数空间间的态射，
$\mathcal{G}$ 为拟凝聚 $\mathcal{O}_Y$-模。设 $x \in |X|$ 且
$y = f(x) \in |Y|$。若
\begin{enumerate}
\item % P09-ERR-0154
$y \in \text{WeakAss}_S(\mathcal{G})$；
\item $f$ 在 $x$ 处平坦；且
\item $f$ 的纤维在 $x$ 处的局部环维数为零（《空间的态射》定义
\ref{spaces-morphisms-definition-dimension-fibre}）；
\end{enumerate}
则 $x \in \text{WeakAss}(f^*\mathcal{G})$。
\end{lemma}

\begin{proof}
取概形 $V$、点 $v \in V$ 以及把 $v$ 映到 $y$ 的平展态射 $V \to Y$。
取概形 $U$、点 $u \in U$ 以及平展态射
% P09-ERR-0155
$U \to V \times_Y X$，它把 $v$ 映到位于 $v$ 与 $x$ 上方的一个点。
这是可行的，因为由《空间的性质》引理
% P09-ERR-0156
\ref{spaces-properties-lemma-points-cartesian}，存在
$t \in |V \times_Y X|$ 映到 $(v, y)$。按定义可见，
$\mathcal{O}_{U_v, u}$ 的维数为零。因此 $u$ 是纤维 $U_v$ 的一个一般点。
由弱伴随点的定义，问题归约为概形间的态射 $U \to V$。这种情形已由
《除子》引理 \ref{divisors-lemma-weakly-ass-pullback} 处理。
\end{proof}

\begin{lemma}
\label{spaces-divisors-lemma-weakly-ass-change-fields}
设 $K/k$ 为域扩张，$X$ 为 $k$ 上的代数空间，$\mathcal{F}$ 为拟凝聚
$\mathcal{O}_X$-模。设 $y \in X_K$ 的像为 $x \in X$。若 $y$ 是拉回
$\mathcal{F}_K$ 的弱伴随点，则 $x$ 是 $\mathcal{F}$ 的弱伴随点。
\end{lemma}

\begin{proof}
这是将《除子》引理 \ref{divisors-lemma-weakly-ass-change-fields}
改述为代数空间的语言所得的陈述。略去这种改述的细节。
\end{proof}

\begin{lemma}
\label{spaces-divisors-lemma-finite-flat-weak-assassin-up-down}
设 $S$ 为概形，$f : X \to Y$ 为代数空间间的有限平坦态射，
$\mathcal{G}$ 为拟凝聚 $\mathcal{O}_Y$-模。设 $x \in |X|$ 的像为
$y \in |Y|$。则
% P09-ERR-0157
$$
x \in \text{WeakAss}(g^*\mathcal{G})
\Leftrightarrow
y \in \text{WeakAss}(\mathcal{G})
$$
\end{lemma}

\begin{proof}
由概形的情形（《平坦性的更多内容》引理
\ref{flat-lemma-finite-flat-weak-assassin-up-down}）经平展局部化立即得出。
\end{proof}

\begin{lemma}
\label{spaces-divisors-lemma-etale-weak-assassin-up-down}
设 $S$ 为概形，$f : X \to Y$ 为代数空间间的平展态射，
$\mathcal{G}$ 为拟凝聚 $\mathcal{O}_Y$-模。设 $x \in |X|$ 的像为
$y \in |Y|$。则
$$
x \in \text{WeakAss}(f^*\mathcal{G})
\Leftrightarrow
y \in \text{WeakAss}(\mathcal{G})
$$
\end{lemma}

\begin{proof}
这由弱伴随点的定义立即得到；事实上，概形情形的对应引理
（《平坦性的更多内容》引理
\ref{flat-lemma-etale-weak-assassin-up-down}）正是本定义的基础。
\end{proof}







\section{相对弱伴随点集}
\label{spaces-divisors-section-relative-weak-assassin}

\noindent
为了定义这个对象，需要先证明几个引理。

\begin{lemma}
\label{spaces-divisors-lemma-locally-noetherian-fibre}
设 $S$ 为概形，$f : X \to Y$ 为 $S$ 上代数空间间的态射，且
$y \in |Y|$。下列条件等价：
\begin{enumerate}
\item 对某个概形 $V$、点 $v \in V$ 以及把 $v$ 映到 $y$ 的平展态射
$V \to Y$，代数空间 $X_v$ 局部诺特；
\item 对每个概形 $V$、点 $v \in V$ 以及把 $v$ 映到 $y$ 的平展态射
$V \to Y$，代数空间 $X_v$ 局部诺特；且
\item 存在域 $k$ 以及表示 $y$ 的态射 $\Spec(k) \to Y$，使得 $X_k$
局部诺特。
\end{enumerate}
若存在域 $k_0$ 以及表示 $y$ 的单态射 $\Spec(k_0) \to Y$，则上述条件还
等价于：
\begin{enumerate}
\item[(4)] 代数空间 $X_{k_0}$ 局部诺特。
\end{enumerate}
\end{lemma}

\begin{proof}
注意 $X_v = v \times_Y X = \Spec(\kappa(v)) \times_Y X$。因此蕴含
(2) $\Rightarrow$ (1) $\Rightarrow$ (3) 是显然的。假设从某个域的谱出发的
态射 $\Spec(k) \to Y$ 使得 $X_k$ 局部诺特。设 $V \to Y$ 是从概形
$V$ 出发的平展态射，并且
% P09-ERR-0158
$v \in V$ 是映到 $y$ 的点。则概形 $v \times_Y \Spec(k)$ 非空。取一点
$w \in v \times_Y \Spec(k)$。考虑态射
$$
X_v \longleftarrow X_w \longrightarrow X_k
$$
由于 $V \to Y$ 平展，且可将 $w$ 视为 $V \times_Y \Spec(k)$ 的一个点，
故 $\kappa(w)/k$ 是有限可分域扩张（《态射》引理
\ref{morphisms-lemma-etale-over-field}）。因此，作为
$w \to \Spec(k)$ 的基变换，$X_w \to X_k$ 是有限平展态射。故
$X_w$ 局部诺特（《空间的态射》引理
\ref{spaces-morphisms-lemma-locally-finite-type-locally-noetherian}）。
态射 $X_w \to X_v$ 是满、仿射、平坦的，因为它是满、仿射、平坦态射
$w \to v$ 的基变换。于是 $X_w$ 局部诺特蕴含 $X_v$ 局部诺特。为看出
这一点，取满平展态射 $U \to X$，再利用 $U_w \to U_v$ 满、仿射且平坦。
在概形 $U_v$ 上仿射局部地工作，由《代数》引理
% P09-ERR-0159
\ref{algebra-lemma-descent-Noetherian} 可得 $U_w$ 局部诺特。

\medskip\noindent
最后，只需在有表示 $y$ 的单态射 $\Spec(k_0) \to Y$ 时证明 (3) 蕴含
(4)。此时 $\Spec(k) \to Y$ 分解为
$\Spec(k) \to \Spec(k_0) \to Y$。上述论证于是表明，$X_k$ 局部诺特
% P09-ERR-0160
蕴含 $X_{k_0}$ 局部诺特。
\end{proof}

\begin{definition}
\label{spaces-divisors-definition-locally-Noetherian-fibre}
设 $S$ 为概形，$f : X \to Y$ 为 $S$ 上代数空间间的态射，且
$y \in |Y|$。若引理 \ref{spaces-divisors-lemma-locally-noetherian-fibre} 的等价条件
(1)、(2)、(3) 成立，则称 {\it $f$ 在 $y$ 上的纤维局部诺特}。若这对每个
$y \in |Y|$ 都成立，则称 {\it $f$ 的各纤维局部诺特}。
\end{definition}

\noindent
当然，保证纤维局部诺特的通常方法是假设该态射局部有限型。

\begin{lemma}
\label{spaces-divisors-lemma-locally-finite-type-locally-Noetherian-fibres}
设 $S$ 为概形，$f : X \to Y$ 为 $S$ 上代数空间间的态射。若 $f$
局部有限型，则 $f$ 的各纤维局部诺特。
\end{lemma}

\begin{proof}
这由《空间的态射》引理
\ref{spaces-morphisms-lemma-locally-finite-type-locally-noetherian}
以及域的谱为诺特概形这一事实得出。
\end{proof}

\begin{lemma}
\label{spaces-divisors-lemma-relative-assassin}
设 $S$ 为概形，$f : X \to Y$ 为 $S$ 上代数空间间的态射。设
$x \in |X|$ 且 $y = f(x) \in |Y|$，并设 $\mathcal{F}$ 为拟凝聚
$\mathcal{O}_X$-模。考虑交换图
$$
\xymatrix{
X \ar[d] & X \times_Y V \ar[d] \ar[l] & X_v \ar[d] \ar[l] \\
Y & V \ar[l] & v \ar[l]
}
\quad
\xymatrix{
X \ar[d] & U \ar[d] \ar[l] & U_v \ar[d] \ar[l] \\
Y & V \ar[l] & v \ar[l]
}
\quad
\xymatrix{
x \ar@{|->}[d] &
x' \ar@{|->}[d] \ar@{|->}[l] &
u \ar@{|->}[ld] \ar@{|->}[l] \\
y &
v \ar@{|->}[l]
}
$$
其中 $V$ 和 $U$ 是概形，$V \to Y$ 与 $U \to X \times_Y V$ 平展，而
$v \in V$、$x' \in |X_v|$、$u \in U$ 是按最后一个图相关联的点。
% P09-ERR-0161
用 $\mathcal{F}|_{X_v}$ 和 $\mathcal{F}|_{U_v}$ 表示 $\mathcal{F}$
的拉回。下列条件等价：
\begin{enumerate}
\item 对某组如上的 $V, v, x'$，点 $x'$ 是
$\mathcal{F}|_{X_v}$ 的弱伴随点；
\item 对每组如上的 $V \to Y, v, x'$，点 $x'$ 是
$\mathcal{F}|_{X_v}$ 的弱伴随点；
\item 对某组如上的 $U, V, u, v$，点 $u$ 是
$\mathcal{F}|_{U_v}$ 的弱伴随点；
\item 对每组如上的 $U, V, u, v$，点 $u$ 是
$\mathcal{F}|_{U_v}$ 的弱伴随点；
\item 对某个域 $k$、某个表示 $y$ 的态射 $\Spec(k) \to Y$，以及某个映到
$x$ 的 $t \in |X_k|$，点 $t$ 是 $\mathcal{F}|_{X_k}$ 的弱伴随点。
\end{enumerate}
若存在域 $k_0$ 以及表示 $y$ 的单态射 $\Spec(k_0) \to Y$，则上述条件还
等价于：
\begin{enumerate}
\item[(6)] $x_0$ 是 $\mathcal{F}|_{X_{k_0}}$ 的弱伴随点，其中
$x_0 \in |X_{k_0}|$ 是映到 $x$ 的唯一点。
\end{enumerate}
若 $f$ 在 $y$ 上的纤维局部诺特，则在条件 (1)、(2)、(3)、(4)、(6) 中，
可以用“伴随”代替“弱伴随”。
\end{lemma}

\begin{proof}
注意，给定引理中如上的 $V, v, x'$，可以找到
$U \to X \times_Y V$ 和映到 $x'$ 的 $u \in U$，于是态射
$U_v \to X_v$ 平展。因此，(1) 与 (3) 等价以及 (2) 与 (4) 等价都是
显然的；其中每个条件都蕴含 (5)。我们将证明 (5) 蕴含 (2)。设给定
$V, v, x'$，以及
% P09-ERR-0162
$\Spec(k) \to X$ 和 $t \in |X_k|$，使得点 $t$ 是
$\mathcal{F}|_{X_k}$ 的弱伴随点。可以取一点
$w \in v \times_Y \Spec(k)$。于是得到态射
$$
X_v \longleftarrow X_w \longrightarrow X_k
$$
由于 $V \to Y$ 平展，且可将 $w$ 视为 $V \times_Y \Spec(k)$ 的一点，
故 $\kappa(w)/k$ 是有限可分域扩张（《态射》引理
\ref{morphisms-lemma-etale-over-field}）。因此，作为
$w \to \Spec(k)$ 的基变换，$X_w \to X_k$ 是有限平展态射。于是由引理
\ref{spaces-divisors-lemma-etale-weak-assassin-up-down}，$X_w$ 中位于 $t$ 上方的任一点
$x''$ 都是 $\mathcal{F}|_{X_w}$ 的弱伴随点。可以选取映到 $x'$ 的
$x''$（《空间的性质》引理
\ref{spaces-properties-lemma-points-cartesian}）。于是引理
\ref{spaces-divisors-lemma-weakly-ass-change-fields} 蕴含 $x'$ 是
$\mathcal{F}|_{X_v}$ 的弱伴随点。

\medskip\noindent
为完成证明，只需说明：若给定 (6) 中的 $\Spec(k_0) \to Y$，则等价条件
(1)--(5) 蕴含 (6)。此时，(5) 中的态射 $\Spec(k) \to Y$ 唯一分解为
$\Spec(k) \to \Spec(k_0) \to Y$。于是 $x_0$ 是 $t$ 在态射
$X_k \to X_{k_0}$ 下的像。因此，上述同一引理表明 (6) 成立。
\end{proof}

\begin{definition}
\label{spaces-divisors-definition-relative-weak-assassin}
设 $S$ 为概形，$f : X \to Y$ 为 $S$ 上代数空间间的态射，
$\mathcal{F}$ 为拟凝聚 $\mathcal{O}_X$-模。{\it $\mathcal{F}$ 在 $X$
中相对于 $Y$ 的相对弱伴随点集}是集合
$\text{WeakAss}_{X/Y}(\mathcal{F}) \subset |X|$，它由满足引理
\ref{spaces-divisors-lemma-relative-assassin} 中等价条件的那些 $x \in |X|$ 组成。
若 $f$ 的各纤维局部诺特（定义
\ref{spaces-divisors-definition-locally-Noetherian-fibre}），则使用记号
$\text{Ass}_{X/Y}(\mathcal{F})$。
\end{definition}

\noindent
用这个记号，可以表述若干已对概形证明的结果。

\begin{lemma}
\label{spaces-divisors-lemma-bourbaki}
设 $S$ 为概形，$f : X \to Y$ 为 $S$ 上代数空间间的态射，
$\mathcal{F}$ 为拟凝聚 $\mathcal{O}_X$-模，而 $\mathcal{G}$ 为拟凝聚
$\mathcal{O}_Y$-模。假设
\begin{enumerate}
\item $\mathcal{F}$ 在 $Y$ 上平坦；
\item $X$ 和 $Y$ 局部诺特；且
\item $f$ 的各纤维局部诺特。
\end{enumerate}
则
$$
\text{Ass}_X(\mathcal{F} \otimes_{\mathcal{O}_X} f^*\mathcal{G}) =
\{x \in \text{Ass}_{X/Y}(\mathcal{F})\text{ 使得 }
f(x) \in \text{Ass}_Y(\mathcal{G}) \}
$$
\end{lemma}

\begin{proof}
经平展局部化，这立即由概形情形的结果得出；参见《除子》引理
\ref{divisors-lemma-bourbaki}。概形情形的结果更一般，仅仅是因为我们没有
为非诺特代数空间定义伴随点（因此，为了使本结果甚至能够表述，必须假设
$X$ 以及 $X \to Y$ 的各纤维局部诺特）。
\end{proof}

\begin{lemma}
\label{spaces-divisors-lemma-base-change-relative-assassin}
设 $S$ 为概形。设
$$
\xymatrix{
X' \ar[d]_{f'} \ar[r]_{g'} & X \ar[d]^f \\
Y' \ar[r]^g & Y
}
$$
为 $S$ 上代数空间的笛卡尔图。设 $\mathcal{F}$ 为拟凝聚
$\mathcal{O}_X$-模，并置 $\mathcal{F}' = (g')^*\mathcal{F}$。若 $f$
局部有限型，则
\begin{enumerate}
\item $x' \in \text{Ass}_{X'/Y'}(\mathcal{F}')
\Rightarrow g'(x') \in \text{Ass}_{X/Y}(\mathcal{F})$；
\item 若 $x \in \text{Ass}_{X/Y}(\mathcal{F})$，则给定满足
$f(x) = g(y')$ 的 $y' \in |Y'|$，存在
$x' \in \text{Ass}_{X'/Y'}(\mathcal{F}')$，使得
$g'(x') = x$ 且 $f'(x') = y'$。
\end{enumerate}
\end{lemma}

\begin{proof}
经平展局部化，这由概形的情形得出。我们完整写出细节。取概形 $V$ 以及满平展
态射 $V \to Y$。取概形 $U$ 以及满平展态射
$U \to V \times_Y X$。取概形 $V'$ 以及满平展态射
$V' \to V \times_Y Y'$。则 $U' = V' \times_V U$ 是概形，并且态射
$U' \to X'$ 满且平展。

\medskip\noindent
(1) 的证明。取映到 $x'$ 的 $u' \in U'$。
% P09-ERR-0163
用 $v' \in V'$ 表示 $u'$ 的像。按定义，
$x' \in \text{Ass}_{X'/Y'}(\mathcal{F}')$ 等价于
$u' \in \text{Ass}(\mathcal{F}|_{U'_{v'}})$（因为 $U'_{v'}$ 局部诺特，
所以写 $\text{Ass}$ 而不是 $\text{WeakAss}$ 是有意义的）。应用《除子》
引理 \ref{divisors-lemma-base-change-relative-assassin} 可见，$u'$ 的像
$u \in U$ 属于 $\text{Ass}(\mathcal{F}|_{U_v})$，其中
$v \in V$ 是 $u$ 的像。这又意味着
$g'(x') \in \text{Ass}_{X/Y}(\mathcal{F})$。

\medskip\noindent
(2) 的证明。取映到 $x$ 的 $u \in U$。
% P09-ERR-0164
用 $v \in V$ 表示 $u$ 的像。按定义，
$x \in \text{Ass}_{X/Y}(\mathcal{F})$ 等价于
$u \in \text{Ass}(\mathcal{F}|_{U_v})$。取一点 $v' \in V'$，它映到
$y' \in |Y'|$ 及 $v \in V$（这由《空间的性质》引理
\ref{spaces-properties-lemma-points-cartesian} 可行）。令
$t \in \Spec(\kappa(v') \otimes_{\kappa(v)} \kappa(u))$
为某个不可约分支的一般点，并令 $u' \in U'$ 为 $t$ 的像。应用《除子》
引理 \ref{divisors-lemma-base-change-relative-assassin} 可见
$u' \in \text{Ass}(\mathcal{F}'|_{U'_{v'}})$。这又意味着
$x' \in \text{Ass}_{X'/Y'}(\mathcal{F}')$，其中 $x' \in |X'|$
是 $u'$ 的像。
\end{proof}

\begin{lemma}
\label{spaces-divisors-lemma-base-change-relative-assassin-quasi-finite}
采用引理 \ref{spaces-divisors-lemma-base-change-relative-assassin} 的记号与假设。假设 $g$
局部拟有限，或更一般地，假设对每个 $y' \in |Y'|$，$y'/g(y')$ 的超越次数
为 $0$。则 $\text{Ass}_{X'/Y'}(\mathcal{F}')$ 是
$\text{Ass}_{X/Y}(\mathcal{F})$ 的逆像。
\end{lemma}

\begin{proof}
点在其像上的超越次数定义见《空间的态射》定义
\ref{spaces-morphisms-definition-dimension-fibre}。设 $x' \in |X'|$
的像为 $x \in |X|$。取概形 $V$ 以及满平展态射 $V \to Y$。取概形
$U$ 以及满平展态射 $U \to V \times_Y X$。取概形 $V'$ 以及满平展态射
$V' \to V \times_Y Y'$。则 $U' = V' \times_V U$ 是概形，并且态射
$U' \to X'$ 满且平展。取映到 $x$ 的 $u \in U$。
% P09-ERR-0165
用 $v \in V$ 表示 $u$ 的像。按定义，
$x \in \text{Ass}_{X/Y}(\mathcal{F})$ 等价于
$u \in \text{Ass}(\mathcal{F}|_{U_v})$。取一点 $u' \in U'$，它映到
$x' \in |X'|$ 及 $u \in U$（这由《空间的性质》引理
\ref{spaces-properties-lemma-points-cartesian} 可行）。令 $v' \in V'$
为 $u'$ 的像。按定义，$x' \in \text{Ass}_{X'/Y'}(\mathcal{F}')$
等价于 $u' \in \text{Ass}(\mathcal{F}'|_{U'_{v'}})$。现在，将《除子》
注记 \ref{divisors-remark-base-change-relative-assassin} 中的讨论应用于
$u' \in \Spec(\kappa(v') \otimes_{\kappa(v)} \kappa(u))$，即得本引理。
\end{proof}

\begin{lemma}
\label{spaces-divisors-lemma-relative-weak-assassin-finite}
设 $S$ 为概形，$f : X \to Y$ 为 $S$ 上代数空间间的态射，
$i : Z \to X$ 为有限态射，而 $\mathcal{G}$ 为拟凝聚
$\mathcal{O}_Z$-模。则
$\text{WeakAss}_{X/Y}(i_*\mathcal{G}) =
i(\text{WeakAss}_{Z/Y}(\mathcal{G}))$。
\end{lemma}

\begin{proof}
经平展局部化，这由概形情形（《除子》引理
\ref{divisors-lemma-relative-weak-assassin-finite}）得出。略去细节。
\end{proof}

\begin{lemma}
\label{spaces-divisors-lemma-relative-assassin-constructible}
设 $Y$ 为概形，$X$ 为 $Y$ 上有限表示的代数空间，$\mathcal{F}$ 为有限表示的
拟凝聚 $\mathcal{O}_X$-模。设 $U \subset X$ 是开子空间，并且
$U \to Y$ 拟紧。则集合
$$
E = \{y \in Y \mid \text{Ass}_{X_y}(\mathcal{F}_y) \subset |U_y|\}
$$
在 $Y$ 中局部可构造。
\end{lemma}

\begin{proof}
注意，由于 $Y$ 是概形，取纤维
$X_y = \Spec(\kappa(y)) \times_Y X$ 是有意义的。（此外，按我们的定义，
集合 $\text{Ass}_{X_y}(\mathcal{F}_y)$ 恰为
$\text{Ass}_{X/Y}(\mathcal{F}) \to Y$ 在 $y$ 上的纤维，但这里不需要
这一点。）问题在 $Y$ 上是局部的；事实上，只需在 $Y$ 仿射时证明 $E$
可构造。在这种情形，$X$ 拟紧。取仿射概形 $W$ 以及满平展态射
$\varphi : W \to X$。则对所有 $y \in Y$，
$\text{Ass}_{X_y}(\mathcal{F}_y)$ 是
$\text{Ass}_{W_y}(\varphi^*\mathcal{F}_y)$ 的像。因此，本引理由开集
$\varphi^{-1}(U) \subset W$ 与态射 $W \to Y$ 的概形情形得出。概形情形
就是《态射的更多内容》引理
\ref{more-morphisms-lemma-relative-assassin-constructible}。
\end{proof}









\section{Fitting 理想}
\label{spaces-divisors-section-fitting-ideals}

\noindent
本节继续《除子》第 \ref{divisors-section-fitting-ideals} 节中的讨论。
设 $S$ 为概形，$X$ 为 $S$ 上的代数空间，而 $\mathcal{F}$ 为有限型拟凝聚
$\mathcal{O}_X$-模。在这种情形，可以构造 Fitting 理想
$$
0 = \text{Fit}_{-1}(\mathcal{F}) \subset \text{Fit}_0(\mathcal{F}) \subset
\text{Fit}_1(\mathcal{F}) \subset \ldots \subset \mathcal{O}_X
$$
，即由下述性质刻画的拟凝聚理想层序列：
% P09-ERR-0166
对每个平展于 $X$ 的仿射概形 $U = \Spec(A)$，若 $\mathcal{F}|_U$ 对应于
$A$-模 $M$，则 $\text{Fit}_i(\mathcal{F})|_U$ 对应于理想
$\text{Fit}_i(M) \subset A$。这是良定义的拟凝聚理想层，因为若
$A \to B$ 是平展环同态，则由《代数的更多内容》引理
\ref{more-algebra-lemma-fitting-ideal-basics} 的 (3)，$M \otimes_A B$
作为 $B$-模的第 $i$ 个 Fitting 理想等于 $\text{Fit}_i(M) B$。更准确地说
（或许），拟凝聚理想层
% P09-ERR-0167
$\text{Fit}_0(\mathcal{O}_X)$ 的存在性例如可由《空间的性质》引理
\ref{spaces-properties-lemma-characterize-quasi-coherent-small-etale}
中对拟凝聚层的描述，以及《除子》引理
\ref{divisors-lemma-base-change-fitting-ideal} 给出的拉回性质得到。

\medskip\noindent
以这种方式构造 Fitting 理想的优点是，可以立即看出 Fitting 理想的形成与
平展局部化交换，因而 Fitting 理想的许多性质立即归约为概形情形的对应性质。
我们将经常使用《空间的性质》第
\ref{spaces-properties-section-properties-modules} 节中的讨论，在概形上的
拟凝聚层与代数空间上的拟凝聚层的性质之间作转化。

\begin{lemma}
\label{spaces-divisors-lemma-base-change-fitting-ideal}
设 $S$ 为概形，$f : X \to Y$ 为 $S$ 上代数空间间的态射，
$\mathcal{F}$ 为有限型拟凝聚 $\mathcal{O}_Y$-模。则
$f^{-1}\text{Fit}_i(\mathcal{F}) \cdot \mathcal{O}_X =
\text{Fit}_i(f^*\mathcal{F})$。
\end{lemma}

\begin{proof}
经平展局部化，归约为《除子》引理
\ref{divisors-lemma-base-change-fitting-ideal}。
\end{proof}

\begin{lemma}
\label{spaces-divisors-lemma-fitting-ideal-of-finitely-presented}
设 $S$ 为概形，$X$ 为 $S$ 上的代数空间，$\mathcal{F}$ 为有限表示
$\mathcal{O}_X$-模。则 $\text{Fit}_r(\mathcal{F})$ 是有限型拟凝聚理想。
\end{lemma}

\begin{proof}
经平展局部化，归约为《除子》引理
\ref{divisors-lemma-fitting-ideal-of-finitely-presented}。
\end{proof}

\begin{lemma}
\label{spaces-divisors-lemma-on-subscheme-cut-out-by-Fit-0}
设 $S$ 为概形，$X$ 为 $S$ 上的代数空间，而 $\mathcal{F}$ 为有限型拟凝聚
$\mathcal{O}_X$-模。设 $Z_0 \subset X$ 为由
$\text{Fit}_0(\mathcal{F})$ 切出的闭子空间，$Z \subset X$ 为
$\mathcal{F}$ 的概形论支撑。则
\begin{enumerate}
\item 作为闭子空间，$Z \subset Z_0 \subset X$；
\item 作为 $|X|$ 的闭子集，
$|Z| = |Z_0| = \text{Supp}(\mathcal{F})$；
\item 存在有限型拟凝聚 $\mathcal{O}_{Z_0}$-模 $\mathcal{G}_0$，使得
$$
(Z_0 \to X)_*\mathcal{G}_0 = \mathcal{F}.
$$
\end{enumerate}
\end{lemma}

\begin{proof}
回忆 $Z$ 的形成与平展局部化交换；参见《空间的态射》定义
\ref{spaces-morphisms-definition-scheme-theoretic-support}（其中用
《空间的态射》引理
\ref{spaces-morphisms-lemma-scheme-theoretic-support} 定义 $Z$）。因此
(1) 和 (2) 由概形情形得出；参见《除子》引理
\ref{divisors-lemma-on-subscheme-cut-out-by-Fit-0}。为得到 (3) 中的
$\mathcal{G}_0$，可使用《空间的态射》引理
\ref{spaces-morphisms-lemma-scheme-theoretic-support} 中 $Z$ 上的
$\mathcal{G}$，并置 $\mathcal{G}_0 = (Z \to Z_0)_*\mathcal{G}$。
\end{proof}

\begin{lemma}
\label{spaces-divisors-lemma-fitting-ideal-generate-locally}
设 $S$ 为概形，$X$ 为 $S$ 上的代数空间，$\mathcal{F}$ 为有限型拟凝聚
$\mathcal{O}_X$-模，且 $x \in |X|$。则 $\mathcal{F}$ 在 $x$ 的某个平展
邻域中可由 $r$ 个元素生成，当且仅当
$\text{Fit}_r(\mathcal{F})_{\overline{x}} = \mathcal{O}_{X, \overline{x}}$。
\end{lemma}

\begin{proof}
经平展局部化，归约为《除子》引理
\ref{divisors-lemma-fitting-ideal-generate-locally}（还要使用《空间的性质》
第 \ref{spaces-properties-section-stalks-structure-sheaf} 节中对局部环的
描述，以及局部环的严格亨泽尔化是忠实平坦的这一事实，以看出严格亨泽尔化
上的等式等价于局部环上的等式）。
\end{proof}

\begin{lemma}
\label{spaces-divisors-lemma-fitting-ideal-finite-locally-free}
设 $S$ 为概形，$X$ 为 $S$ 上的代数空间，$\mathcal{F}$ 为有限型拟凝聚
$\mathcal{O}_X$-模。设 $r \geq 0$。下列条件等价：
\begin{enumerate}
\item $\mathcal{F}$ 是秩为 $r$ 的有限局部自由模；
\item $\text{Fit}_{r - 1}(\mathcal{F}) = 0$ 且
$\text{Fit}_r(\mathcal{F}) = \mathcal{O}_X$；且
\item 对 $k < r$ 有 $\text{Fit}_k(\mathcal{F}) = 0$，而对
$k \geq r$ 有 $\text{Fit}_k(\mathcal{F}) = \mathcal{O}_X$。
\end{enumerate}
\end{lemma}

\begin{proof}
经平展局部化，归约为《除子》引理
\ref{divisors-lemma-fitting-ideal-finite-locally-free}。
\end{proof}

\begin{lemma}
\label{spaces-divisors-lemma-locally-free-rank-r-pullback}
设 $S$ 为概形，$X$ 为 $S$ 上的代数空间，$\mathcal{F}$ 为有限型拟凝聚
$\mathcal{O}_X$-模。由 $\mathcal{F}$ 的 Fitting 理想定义的闭子空间
% P09-ERR-0168
$$
X = Z_{-1} \supset Z_0 \supset Z_1 \supset Z_2 \ldots
$$
具有下列性质：
\begin{enumerate}
\item 交集 $\bigcap Z_r$ 为空。
\item 由下述规则定义的函子
$(\Sch/X)^{opp} \to \textit{Sets}$
$$
T \longmapsto
\left\{
\begin{matrix}
\{*\} & \text{若 }\mathcal{F}_T\text{ 局部可由 }
\leq r\text{ 个截面生成} \\
\emptyset & \text{否则}
\end{matrix}
\right.
$$
可由开子空间 $X \setminus Z_r$ 表示。
\item 由下述规则定义的函子
$F_r : (\Sch/X)^{opp} \to \textit{Sets}$
% P09-ERR-0169
$$
T \longmapsto
\left\{
\begin{matrix}
\{*\} & \text{若 }\mathcal{F}_T\text{ 是局部自由且秩为 }r\\
\emptyset & \text{否则}
\end{matrix}
\right.
$$
可由 $X$ 的局部闭子空间 $Z_{r - 1} \setminus Z_r$ 表示。
\end{enumerate}
若 $\mathcal{F}$ 有限表示，则
$Z_r \to X$、$X \setminus Z_r \to X$ 及
$Z_{r - 1} \setminus Z_r \to X$ 都是有限表示态射。
\end{lemma}

\begin{proof}
经平展局部化，归约为《除子》引理
\ref{divisors-lemma-locally-free-rank-r-pullback}。
\end{proof}

\begin{lemma}
\label{spaces-divisors-lemma-finite-presentation-module}
设 $S$ 为概形，$X$ 为 $S$ 上的代数空间，$\mathcal{F}$ 为有限表示
$\mathcal{O}_X$-模。设
% P09-ERR-0170
$X = Z_{-1} \subset Z_0 \subset Z_1 \subset \ldots$
如引理 \ref{spaces-divisors-lemma-locally-free-rank-r-pullback} 中所述。置
$X_r = Z_{r - 1} \setminus Z_r$。则
$X' = \coprod_{r \geq 0} X_r$ 表示函子
% P09-ERR-0171
% P09-ERR-0172
$$
F_{flat} : \Sch/X \longrightarrow \textit{Sets},\quad\quad
T \longmapsto
\left\{
\begin{matrix}
\{*\} & \text{若 }\mathcal{F}_T\text{ 平坦于 }T\\
\emptyset & \text{否则}
\end{matrix}
\right.
$$
。此外，$\mathcal{F}|_{X_r}$ 是秩为 $r$ 的局部自由模，且态射
$X_r \to X$ 与 $X' \to X$ 都是有限表示态射。
\end{lemma}

\begin{proof}
经平展局部化，归约为《除子》引理
\ref{divisors-lemma-finite-presentation-module}。
\end{proof}














\section{有效 Cartier 除子}
\label{spaces-divisors-section-effective-Cartier-divisors}

\noindent
% P09-ERR-0173
出于某种原因，先于其他一切概念定义有效 Cartier 除子的概念似乎较为方便。注意，
在《空间的态射》第 \ref{spaces-morphisms-section-closed-immersions} 节中，
我们讨论了闭子空间与拟凝聚理想层之间的对应。此外，在《空间的性质》第
\ref{spaces-properties-section-properties-modules} 节中，我们讨论了拟凝聚模的
性质，特别是“局部由 $1$ 个元素生成”这一性质。这些参考结果表明，下面的定义
与概形情形的定义相容。

\begin{definition}
\label{spaces-divisors-definition-effective-Cartier-divisor}
设 $S$ 为概形，$X$ 为 $S$ 上的代数空间。
\begin{enumerate}
\item $X$ 的一个{\it 局部主闭子空间}是指其理想层局部由 $1$ 个元素生成的
闭子空间。
\item $X$ 上的一个{\it 有效 Cartier 除子}是指闭子空间 $D \subset X$，
使得理想层 $\mathcal{I}_D \subset \mathcal{O}_X$ 是可逆
$\mathcal{O}_X$-模。
\end{enumerate}
\end{definition}

\noindent
因此，有效 Cartier 除子是局部主闭子空间，但逆命题并非总成立。有效
Cartier 除子是在最强意义下具有纯余维 $1$ 的闭子空间；也就是说，它们局部由
一个非零因子切出。特别地，它们无处稠密。

\begin{lemma}
\label{spaces-divisors-lemma-characterize-effective-Cartier-divisor}
设 $S$ 为概形，$X$ 为 $S$ 上的代数空间，并设 $D \subset X$ 为闭子空间。
下列条件等价：
\begin{enumerate}
\item 子空间 $D$ 是 $X$ 上的有效 Cartier 除子。
\item 存在某个概形 $U$ 及满平展态射 $U \to X$，使得逆像
$D \times_X U$ 是 $U$ 上的有效 Cartier 除子。
\item 对任意概形 $U$ 及任意平展态射 $U \to X$，逆像
$D \times_X U$ 都是 $U$ 上的有效 Cartier 除子。
\item 对每个 $x \in |D|$，存在带点代数空间的平展态射
$(U, u) \to (X, x)$，使得 $U = \Spec(A)$ 且
$D \times_X U = \Spec(A/(f))$，其中 $f \in A$ 不是零因子。
\end{enumerate}
\end{lemma}

\begin{proof}
(1)--(3) 的等价性由定义 \ref{spaces-divisors-definition-effective-Cartier-divisor} 及其前面的
参考结果立即得到。假设 (1) 成立，并令 $x \in |D|$。选取概形 $W$ 及满平展
态射 $W \to X$。选取映到 $x$ 的点 $w \in D \times_X W$。由 (3)，
$D \times_X W$ 是 $W$ 上的有效 Cartier 除子。于是，按照《除子》引理
\ref{divisors-lemma-characterize-effective-Cartier-divisor}，在 $W$ 中选取 $w$
的一个仿射开邻域，即可得到仿射平展邻域 $U$。
% P09-ERR-0174

\medskip\noindent
假设 (4) 成立。由《除子》引理
\ref{divisors-lemma-characterize-effective-Cartier-divisor} 可知，
$\mathcal{I}_D|_U$ 可逆。所有这样的 $U$ 与 $X \setminus D$ 合在一起给出
$X$ 的一个平展覆盖，故 $\mathcal{I}_D$ 是可逆 $\mathcal{O}_X$-模。
\end{proof}

\begin{lemma}
\label{spaces-divisors-lemma-complement-locally-principal-closed-subscheme}
设 $S$ 为概形，$X$ 为 $S$ 上的代数空间。设 $Z \subset X$ 为局部主闭子空间，
并令 $U = X \setminus Z$。则 $U \to X$ 是仿射态射。
\end{lemma}

\begin{proof}
% P09-ERR-0175
该问题在 $X$ 上是平展局部的，见《空间的态射》引理
\ref{spaces-morphisms-lemma-affine-local} 和引理
\ref{spaces-divisors-lemma-characterize-effective-Cartier-divisor}。因此，本结论归约为概形情形，
即《除子》引理
\ref{divisors-lemma-complement-locally-principal-closed-subscheme}。
\end{proof}

\begin{lemma}
\label{spaces-divisors-lemma-complement-effective-Cartier-divisor}
设 $S$ 为概形，$X$ 为 $S$ 上的代数空间。设 $D \subset X$ 为有效 Cartier
除子，并令 $U = X \setminus D$。则 $U \to X$ 是仿射态射，且 $U$ 在 $X$
中是概形论稠密的。
\end{lemma}

\begin{proof}
仿射性由引理 \ref{spaces-divisors-lemma-complement-locally-principal-closed-subscheme} 给出。
由《空间的态射》定义
\ref{spaces-morphisms-definition-scheme-theoretically-dense}，稠密性问题在 $X$
上是平展局部的。因此，本结论归约为概形情形，即《除子》引理
\ref{divisors-lemma-complement-effective-Cartier-divisor}。
\end{proof}

\begin{lemma}
\label{spaces-divisors-lemma-effective-Cartier-makes-dimension-drop}
设 $S$ 为概形，$X$ 为 $S$ 上的代数空间。设 $D \subset X$ 为有效 Cartier
除子，并令 $x \in |D|$。若 $\dim_x(X) < \infty$，则
$\dim_x(D) < \dim_x(X)$。
\end{lemma}

\begin{proof}
% P09-ERR-0176
有效 Cartier 除子的定义以及代数空间在一点处的维数之定义（《空间的性质》
定义 \ref{spaces-properties-definition-dimension-at-point}）都是平展局部的。
因此，本引理归约为概形情形，即《除子》引理
\ref{divisors-lemma-effective-Cartier-makes-dimension-drop}。
\end{proof}

\begin{definition}
\label{spaces-divisors-definition-sum-effective-Cartier-divisors}
设 $S$ 为概形，$X$ 为 $S$ 上的代数空间。给定 $X$ 上的有效 Cartier 除子
$D_1$、$D_2$，我们令 $D = D_1 + D_2$ 为 $X$ 的闭子空间，它对应于拟凝聚
理想层
% P09-ERR-0177
$\mathcal{I}_{D_1}\mathcal{I}_{D_2} \subset \mathcal{O}_S$。
称它为{\it 有效 Cartier 除子 $D_1$ 与 $D_2$ 的和}。
\end{definition}

\noindent
显然，给定 $X$ 上有限多个有效 Cartier 除子 $D_i$ 及非负整数 $n_i$，我们可以
定义和 $\sum n_iD_i$。

\begin{lemma}
\label{spaces-divisors-lemma-sum-effective-Cartier-divisors}
两个有效 Cartier 除子的和仍是有效 Cartier 除子。
\end{lemma}

\begin{proof}
略。经平展局部化，这归约为如下简单的代数事实：若 $f_1, f_2 \in A$ 是环
$A$ 中的非零因子，则 $f_1f_2 \in A$ 也是非零因子。
\end{proof}

\begin{lemma}
\label{spaces-divisors-lemma-sum-closed-subschemes-effective-Cartier}
设 $S$ 为概形，$X$ 为 $S$ 上的代数空间。设 $Z, Y$ 是 $X$ 的两个闭子空间，
其理想层分别为 $\mathcal{I}$ 和 $\mathcal{J}$。若
$\mathcal{I}\mathcal{J}$ 定义了有效 Cartier 除子 $D \subset X$，则 $Z$ 与
$Y$ 都是有效 Cartier 除子，并且 $D = Z + Y$。
\end{lemma}

\begin{proof}
由引理 \ref{spaces-divisors-lemma-characterize-effective-Cartier-divisor}，这归约为概形情形，
即《除子》引理 \ref{divisors-lemma-sum-closed-subschemes-effective-Cartier}。
\end{proof}

\noindent
回忆一下，我们已在《空间的态射》定义
\ref{spaces-morphisms-definition-inverse-image-closed-subspace} 中定义了闭子空间
在任意代数空间态射下的逆像。

\begin{lemma}
\label{spaces-divisors-lemma-pullback-locally-principal}
设 $S$ 为概形，$f : X' \to X$ 为 $S$ 上代数空间的态射，并设
$Z \subset X$ 为局部主闭子空间。则逆像 $f^{-1}(Z)$ 是 $X'$ 的局部主闭
子空间。
\end{lemma}

\begin{proof}
略。
\end{proof}

\begin{definition}
\label{spaces-divisors-definition-pullback-effective-Cartier-divisor}
设 $S$ 为概形，$f : X' \to X$ 为 $S$ 上代数空间的态射，并设 $D \subset X$
为有效 Cartier 除子。若闭子空间 $f^{-1}(D) \subset X'$ 是有效 Cartier
除子，则称{\it $D$ 沿 $f$ 的拉回有定义}。在这种情形下，将它记作 $f^*D$
或 $f^{-1}(D)$，并称之为{\it 该有效 Cartier 除子的拉回}。
\end{definition}

\noindent
在实际应用中，$f^{-1}(D)$ 是有效 Cartier 除子这一条件经常成立。

\begin{lemma}
\label{spaces-divisors-lemma-pullback-effective-Cartier-defined}
设 $S$ 为概形，$f : X \to Y$ 为 $S$ 上代数空间的态射，并设 $D \subset Y$
为有效 Cartier 除子。在下列每一种情形中，$D$ 沿 $f$ 的拉回都有定义：
\begin{enumerate}
\item 对 $X$ 的任意弱伴随点 $x$，均有 $f(x) \not \in |D|$；
\item $f$ 是平坦的；以及
% P09-ERR-0178
\item 按需在此补充。
\end{enumerate}
\end{lemma}

\begin{proof}
经平展局部化，本引理归约为概形情形，见《除子》引理
\ref{divisors-lemma-pullback-effective-Cartier-defined}。
\end{proof}

\begin{lemma}
\label{spaces-divisors-lemma-pullback-effective-Cartier-divisors-additive}
设 $S$ 为概形，$f : X' \to X$ 为 $S$ 上代数空间的态射，并设 $D_1$、
$D_2$ 为 $X$ 上的有效 Cartier 除子。若 $D_1$ 与 $D_2$ 的拉回都有定义，
则 $D = D_1 + D_2$ 的拉回也有定义，并且
$f^*D = f^*D_1 + f^*D_2$。
\end{lemma}

\begin{proof}
略。
\end{proof}




\section{有效 Cartier 除子与可逆层}
\label{spaces-divisors-section-effective-Cartier-invertible}

\noindent
由于有效 Cartier 除子具有可逆理想层（定义
\ref{spaces-divisors-definition-effective-Cartier-divisor}），下面的定义是有意义的。

\begin{definition}
\label{spaces-divisors-definition-invertible-sheaf-effective-Cartier-divisor}
设 $S$ 为概形，$X$ 为 $S$ 上的代数空间，并设 $D \subset X$ 为有效 Cartier
除子，其理想层为 $\mathcal{I}_D$。
\begin{enumerate}
\item {\it 与 $D$ 相伴的可逆层 $\mathcal{O}_X(D)$} 定义为
$$
\mathcal{O}_X(D) =
\SheafHom_{\mathcal{O}_X}(\mathcal{I}_D, \mathcal{O}_X) =
\mathcal{I}_D^{\otimes -1}.
$$
\item 典范截面通常记作 $1$ 或 $1_D$，它是 $\mathcal{O}_X(D)$ 的整体截面，
对应于包含映射 $\mathcal{I}_D \to \mathcal{O}_X$。
\item 记
$\mathcal{O}_X(-D) = \mathcal{O}_X(D)^{\otimes -1} = \mathcal{I}_D$.
\item 给定另一个有效 Cartier 除子 $D' \subset X$，定义
$\mathcal{O}_X(D - D') =
\mathcal{O}_X(D) \otimes_{\mathcal{O}_X} \mathcal{O}_X(-D')$.
\end{enumerate}
\end{definition}

\noindent
作几点说明。下文将看到，赋值 $D \mapsto \mathcal{O}_X(D)$ 把有效 Cartier
除子的加法（定义 \ref{spaces-divisors-definition-sum-effective-Cartier-divisors}）变为 $X$ 的
Picard 群中的加法（引理
\ref{spaces-divisors-lemma-invertible-sheaf-sum-effective-Cartier-divisors}）。然而，上述定义中
的表达式 $D - D'$ 本身没有任何几何意义。更精确地，可以把 $X$ 上所有有效
Cartier 除子构成的集合看成交换幺半群 $\text{EffCart}(X)$，其零元为空的有效
Cartier 除子。于是，赋值 $(D, D') \mapsto \mathcal{O}_X(D - D')$ 定义了群同态
$$
\text{EffCart}(X)^{gp} \longrightarrow \Pic(X)
$$
，其中左端是 $\text{EffCart}(X)$ 的群化。换言之，在写
$\mathcal{O}_X(D - D')$ 时，可以把 $D - D'$ 看成
$\text{EffCart}(X)^{gp}$ 的一个元素。

\begin{lemma}
\label{spaces-divisors-lemma-conormal-effective-Cartier-divisor}
设 $S$ 为概形，$X$ 为 $S$ 上的代数空间，并设 $D \subset X$ 为有效 Cartier
除子。则余法层满足
% P09-ERR-0179
$\mathcal{C}_{D/X} = \mathcal{I}_D|D =
\mathcal{O}_X(D)^{\otimes -1}|_D$。
\end{lemma}

\begin{proof}
略。
\end{proof}

\begin{lemma}
\label{spaces-divisors-lemma-invertible-sheaf-sum-effective-Cartier-divisors}
设 $S$ 为概形，$X$ 为 $S$ 上的代数空间。设 $D_1$、$D_2$ 为 $X$ 上的有效
Cartier 除子，并令 $D = D_1 + D_2$。则存在唯一的同构
$$
\mathcal{O}_X(D_1) \otimes_{\mathcal{O}_X} \mathcal{O}_X(D_2)
\longrightarrow
\mathcal{O}_X(D)
$$
，它把 $1_{D_1} \otimes 1_{D_2}$ 映到 $1_D$。
\end{lemma}

\begin{proof}
略。
\end{proof}

\begin{definition}
\label{spaces-divisors-definition-regular-section}
设 $S$ 为概形，$X$ 为 $S$ 上的代数空间，并设 $\mathcal{L}$ 为 $X$ 上的可逆
层。若映射 $\mathcal{O}_X \to \mathcal{L}$，$f \mapsto fs$ 是单射，则称整体
截面 $s \in \Gamma(X, \mathcal{L})$ 为{\it 正则截面}。
\end{definition}

\begin{lemma}
\label{spaces-divisors-lemma-regular-section-structure-sheaf}
设 $S$ 为概形，$X$ 为 $S$ 上的代数空间，并设
$f \in \Gamma(X, \mathcal{O}_X)$。下列条件等价：
\begin{enumerate}
\item $f$ 是正则截面；
\item 对任意 $x \in X$，像 $f \in \mathcal{O}_{X, \overline{x}}$ 都不是
零因子；
\item 对任意平展于 $X$ 的仿射概形 $U = \Spec(A)$，限制 $f|_U$ 都是 $A$
中的非零因子；以及
\item 存在概形 $U$ 及满平展态射 $U \to X$，使得 $f|_U$ 是
$\mathcal{O}_U$ 的正则截面。
\end{enumerate}
\end{lemma}

\begin{proof}
略。
\end{proof}

\noindent
注意，可逆 $\mathcal{O}_X$-模 $\mathcal{L}$ 的整体截面 $s$ 可以看作
$\mathcal{O}_X$-模映射 $s : \mathcal{O}_X \to \mathcal{L}$。因此，其对偶是
映射 $s : \mathcal{L}^{\otimes -1} \to \mathcal{O}_X$。（关于对偶可逆层，
见《位点上的模》引理 \ref{sites-modules-lemma-constructions-invertible}。）

\begin{definition}
\label{spaces-divisors-definition-zero-scheme-s}
设 $S$ 为概形，$X$ 为 $S$ 上的代数空间。设 $\mathcal{L}$ 为可逆层，并设
$s \in \Gamma(X, \mathcal{L})$。$s$ 的{\it 零点概形}是闭子空间
$Z(s) \subset X$，它由拟凝聚理想层 $\mathcal{I} \subset \mathcal{O}_X$
定义；该理想层是映射
$s : \mathcal{L}^{\otimes -1} \to \mathcal{O}_X$ 的像。
\end{definition}

\begin{lemma}
\label{spaces-divisors-lemma-zero-scheme}
设 $S$ 为概形，$X$ 为 $S$ 上的代数空间。设 $\mathcal{L}$ 为可逆
$\mathcal{O}_X$-模，并设 $s \in \Gamma(X, \mathcal{L})$。
\begin{enumerate}
\item 考虑满足下述条件的闭浸入 $i : Z \to X$：
% P09-ERR-0180
$i^*s \in \Gamma(Z, i^*\mathcal{L}))$ 为零；并按包含关系对这些闭浸入排序。
零点概形 $Z(s)$ 是这个偏序集的最大元。
\item 对任意 $S$ 上代数空间的态射 $f : Y \to X$，$f^*s = 0$ 在
$\Gamma(Y, f^*\mathcal{L})$ 中成立，当且仅当 $f$ 经由 $Z(s)$ 分解。
\item 零点概形 $Z(s)$ 是 $X$ 的局部主闭子空间。
\item 零点概形 $Z(s)$ 是 $X$ 上的有效 Cartier 除子，当且仅当 $s$ 是
$\mathcal{L}$ 的正则截面。
\end{enumerate}
\end{lemma}

\begin{proof}
略。
\end{proof}

\begin{lemma}
\label{spaces-divisors-lemma-characterize-OD}
设 $S$ 为概形，$X$ 为 $S$ 上的代数空间。
\begin{enumerate}
\item 若 $D \subset X$ 是有效 Cartier 除子，则 $\mathcal{O}_X(D)$ 的典范
截面 $1_D$ 是正则的。
\item 反之，若 $s$ 是可逆层 $\mathcal{L}$ 的正则截面，则存在唯一的有效
Cartier 除子 $D = Z(s) \subset X$ 以及唯一的同构
$\mathcal{O}_X(D) \to \mathcal{L}$，该同构把 $1_D$ 映到 $s$。
\end{enumerate}
构造
$D \mapsto (\mathcal{O}_X(D), 1_D)$ 与 $(\mathcal{L}, s) \mapsto Z(s)$
给出下列集合之间的互逆映射：
$$
\left\{
\begin{matrix}
\text{有效 Cartier 除子，其所在空间为 }X
\end{matrix}
\right\}
\leftrightarrow
\left\{
\begin{matrix}
\text{二元组 }(\mathcal{L}, s)\text{，其中所含对象为一个可逆}\\
\mathcal{O}_X\text{-模和一个正则整体截面}
\end{matrix}
\right\}
$$
\end{lemma}

\begin{proof}
略。
\end{proof}





\section{诺特空间上的有效 Cartier 除子}
\label{spaces-divisors-section-Noetherian-effective-Cartier}

\noindent
在局部诺特情形下，关于有效 Cartier 除子和正则截面的大部分讨论都会有所简化。

\begin{lemma}
\label{spaces-divisors-lemma-effective-Cartier-divisor-Sk}
设 $S$ 为概形，$X$ 为 $S$ 上的局部诺特代数空间。设 $D \subset X$ 为有效
Cartier 除子。若 $X$ 满足 $(S_k)$，则 $D$ 满足 $(S_{k - 1})$。
\end{lemma}

\begin{proof}
由我们对代数空间的性质 $(S_k)$ 所作的定义（《空间的性质》第
\ref{spaces-properties-section-types-properties} 节）以及引理
\ref{spaces-divisors-lemma-characterize-effective-Cartier-divisor}，本结论归约为概形情形
（《除子》引理 \ref{divisors-lemma-effective-Cartier-divisor-Sk}）。
\end{proof}

\begin{lemma}
\label{spaces-divisors-lemma-normal-effective-Cartier-divisor-S1}
设 $S$ 为概形，$X$ 为 $S$ 上的局部诺特正规代数空间。设 $D \subset X$ 为
有效 Cartier 除子。则 $D$ 满足 $(S_1)$。
\end{lemma}

\begin{proof}
由我们对代数空间的正规性所作的定义（《空间的性质》第
\ref{spaces-properties-section-types-properties} 节）以及引理
\ref{spaces-divisors-lemma-characterize-effective-Cartier-divisor}，本结论归约为概形情形
（《除子》引理 \ref{divisors-lemma-normal-effective-Cartier-divisor-S1}）。
\end{proof}

\noindent
下面的引理有时可用于构造有效 Cartier 除子。

\begin{lemma}
\label{spaces-divisors-lemma-complement-open-affine-effective-cartier-divisor}
设 $S$ 为概形，$X$ 为 $S$ 上的正则诺特分离代数空间。设 $U \subset X$ 为
稠密仿射开子空间。则存在有效 Cartier 除子 $D \subset X$，使得
$U = X \setminus D$。
\end{lemma}

\begin{proof}
我们断言，$X \setminus U$ 上的约化诱导代数空间结构 $D$（《空间的性质》
定义 \ref{spaces-properties-definition-reduced-induced-space}）就是所求的有效
Cartier 除子。$D$ 的构造与平展局部化可交换，见《空间的性质》引理
\ref{spaces-properties-lemma-reduced-closed-subspace} 的证明。取满平展态射
$X' \to X$，其中 $X'$ 为仿射概形。由于 $X$ 分离，
$U' = X' \times_X U$ 是仿射的。由于 $|X'| \to |X|$ 是开映射，$U'$ 在 $X'$
中稠密。又由于 $D' = X' \times_X D$ 是 $X' \setminus U'$ 上诱导的约化概形
结构，由《除子》引理
\ref{divisors-lemma-complement-open-affine-effective-cartier-divisor} 及其证明，
$D'$ 是有效 Cartier 除子。这正是所要证明的。
\end{proof}

\begin{lemma}
\label{spaces-divisors-lemma-Noetherian-regular-separated-pic-effective-Cartier}
设 $S$ 为概形，$X$ 为 $S$ 上的正则诺特分离代数空间。则每个可逆
$\mathcal{O}_X$-模都同构于
$$
\mathcal{O}_X(D - D') =
\mathcal{O}_X(D) \otimes_{\mathcal{O}_X} \mathcal{O}_X(D')^{\otimes -1}
$$
，其中 $D, D'$ 是 $X$ 上的某些有效 Cartier 除子。
\end{lemma}

\begin{proof}
设 $\mathcal{L}$ 为可逆 $\mathcal{O}_X$-模。选取稠密仿射开子空间
$U \subset X$，使得 $\mathcal{L}|_U$ 平凡。这是可能的，因为 $X$ 有一个为
概形的稠密开子空间，见《空间的性质》命题
\ref{spaces-properties-proposition-locally-quasi-separated-open-dense-scheme}。
% P09-ERR-0181
将平凡化记作 $s : \mathcal{O}_U \to \mathcal{L}|_U$。$U$ 的补集是一个有效
Cartier 除子 $D$。我们断言，对某个 $n > 0$，映射 $s$ 唯一延拓为映射
$$
s : \mathcal{O}_X(-nD) \longrightarrow \mathcal{L}
$$
。该断言蕴含本引理，因为它表明
$\mathcal{L} \otimes_{\mathcal{O}_X} \mathcal{O}_X(nD)$ 有一个正则整体截面，
故由引理 \ref{spaces-divisors-lemma-characterize-OD}，它同构于某个有效 Cartier 除子 $D'$
所对应的 $\mathcal{O}_X(D')$。为证明该断言，可以在平展局部工作。因此可设
$X$ 为仿射诺特概形。由于
$\mathcal{O}_X(-nD) = \mathcal{I}^n$，其中
$\mathcal{I} = \mathcal{O}_X(-D)$ 是 $D$ 在 $X$ 中的理想层，此情形由
《概形的上同调》引理 \ref{coherent-lemma-homs-over-open} 得到。
\end{proof}

\noindent
下面的引理其实应当放在另一节中。

\begin{lemma}
\label{spaces-divisors-lemma-smooth-over-valuation-ring-effective-Cartier}
设 $R$ 为赋值环，其分式域为 $K$。设 $X$ 为 $R$ 上的代数空间，且
$X \to \Spec(R)$ 光滑。对每个有效 Cartier 除子 $D \subset X_K$，都存在
有效 Cartier 除子 $D' \subset X$，使得 $D'_K = D$。
\end{lemma}

\begin{proof}
令 $D' \subset X$ 为 $D \to X_K \to X$ 的概形论像。由于该态射拟紧，$D'$
的形成与平坦基变换可交换，见《空间的态射》引理
\ref{spaces-morphisms-lemma-flat-base-change-scheme-theoretic-image}。特别地，
$D'_K = D$。因此可以假设 $X$ 仿射，记 $X = \Spec(A)$。于是
$X_K = \Spec(A \otimes_R K)$，且 $D$ 对应于理想
$I \subset A \otimes_R K$。我们必须证明 $J = I \cap A$ 在 $X$ 中切出一个
有效 Cartier 除子。首先注意，$A/J$ 在 $R$ 上平坦（因为它是无挠 $R$-模，
见《代数进阶》引理
\ref{more-algebra-lemma-valuation-ring-torsion-free-flat}），因而由《代数进阶》
引理
\ref{more-algebra-lemma-flat-finite-type-valuation-ring-finite-presentation}
及《代数》引理 \ref{algebra-lemma-extension}，$J$ 是有限生成的。因此，只需证明
对每个素理想 $\mathfrak q \subset A$，$J_\mathfrak q \subset A_\mathfrak q$
均由单个元素生成。令 $\mathfrak p = R \cap \mathfrak q$。则
$R_\mathfrak p$ 是赋值环（《代数》引理
\ref{algebra-lemma-make-valuation-rings}）。另由《代数》引理
\ref{algebra-lemma-characterize-smooth-over-field}，
$A_\mathfrak q/\mathfrak p A_\mathfrak q$ 是正则环。因此，可以应用《代数
进阶》引理 \ref{more-algebra-lemma-picard-group-generic-fibre-regular} 得到
$I(A_\mathfrak q \otimes_R K)$ 由单个元素
$f \in A_\mathfrak p \otimes_R K$ 生成。清去分母后，可以假设
$f \in A_\mathfrak q$。令 $\mathfrak c \subset R_\mathfrak p$ 为 $f$ 的内容
理想（见《代数进阶》定义 \ref{more-algebra-definition-content-ideal} 及《平坦性
进阶》引理 \ref{flat-lemma-flat-finite-type-local-valuation-ring-has-content}）。
由于 $R_\mathfrak p$ 是赋值环，且 $\mathfrak c$ 有限生成（《代数进阶》引理
\ref{more-algebra-lemma-content-finitely-generated}），存在
$\pi \in R_\mathfrak p$ 使得 $\mathfrak c = (\pi)$（《代数》引理
\ref{algebra-lemma-characterize-valuation-ring}）。
% P09-ERR-0182
把 $f$ 替换为 $\pi^{-1}f$ 后，有 $f \in A_\mathfrak q$ 且
$f \not \in \mathfrak pA_\mathfrak q$。
% P09-ERR-0183
断言：$I_\mathfrak q = (f)$；这将完成证明。为证明该断言，注意
$f \in I_\mathfrak q$。因而有满射
$A_\mathfrak q/(f) \to A_\mathfrak q/I_\mathfrak q$，它在 $R$ 上与 $K$ 张量
后成为同构。因此，只要 $A_\mathfrak q/(f)$ 在 $R_\mathfrak p$ 上平坦，证明
就完成了。这由《代数》引理 \ref{algebra-lemma-grothendieck-general} 及我们对
$f$ 的选择得到。
\end{proof}









\section{相对有效 Cartier 除子}
\label{spaces-divisors-section-effective-Cartier-morphisms}

\noindent
下面的引理表明，在基上平坦的有效 Cartier 除子实际上是基上的一个“有效
Cartier 除子族”。例如，它在任意纤维上的限制都是有效 Cartier 除子。

\begin{lemma}
\label{spaces-divisors-lemma-relative-Cartier}
设 $S$ 为概形，$f : X \to Y$ 为 $S$ 上代数空间的态射，并设
$D \subset X$ 为闭子空间。假设
\begin{enumerate}
\item $D$ 是有效 Cartier 除子；并且
\item $D \to Y$ 是平坦态射。
\end{enumerate}
% P09-ERR-0184
则对每个概形态射 $g : Y' \to Y$，拉回 $(g')^{-1}D$ 都是
$X' = Y' \times_Y X$ 上的有效 Cartier 除子，其中
$g' : X' \to X$ 是投影。
\end{lemma}

\begin{proof}
由引理 \ref{spaces-divisors-lemma-characterize-effective-Cartier-divisor}，成为有效 Cartier
除子这一性质是平展局部的。因此，本引理立即归约为概形情形，即《除子》引理
\ref{divisors-lemma-relative-Cartier}。
% P09-ERR-0185
\end{proof}

\noindent
这个引理促成了下面的定义。

\begin{definition}
\label{spaces-divisors-definition-relative-effective-Cartier-divisor}
设 $S$ 为概形，$f : X \to Y$ 为 $S$ 上代数空间的态射。$X/Y$ 上的一个
{\it 相对有效 Cartier 除子}是指有效 Cartier 除子 $D \subset X$，使得
$D \to Y$ 是代数空间的平坦态射。
\end{definition}








\section{亚纯函数与截面}
\label{spaces-divisors-section-meromorphic-functions}

\noindent
本节对应于《除子》第 \ref{divisors-section-meromorphic-functions} 节。请注意：
% P09-ERR-0186
在代数空间情形下，处理这些材料甚至比在概形情形下更容易出错！

\medskip\noindent
设 $S$ 为概形，$X$ 为 $S$ 上的代数空间。对任意平展于 $X$ 的概形 $U$，
我们已经定义了集合
$\mathcal{S}(U) \subset \mathcal{O}_X(U)$
，它由 $U$ 上 $\mathcal{O}_X$ 的正则截面组成，见定义
\ref{spaces-divisors-definition-regular-section}。正则截面限制到平展的 $V/U$ 上仍是正则的。
因此，$\mathcal{S} : U \mapsto \mathcal{S}(U)$ 是 $\mathcal{O}_X$ 的一个
（集合）子层。若要标明它对 $X$ 的依赖性，有时记
$\mathcal{S} = \mathcal{S}_X$。此外，对每个 $U$，$\mathcal{S}(U)$ 都是环
$\mathcal{O}_X(U)$ 的乘法子集。因此，可以考虑 $X_\etale$ 上的环预层
$$
U \longmapsto \mathcal{S}(U)^{-1} \mathcal{O}_X(U),
$$
及其层化，见《位点上的模》第
\ref{sites-modules-section-localizing-sheaves-rings} 节。

\begin{definition}
\label{spaces-divisors-definition-sheaf-meromorphic-functions}
设 $S$ 为概形，$X$ 为 $S$ 上的代数空间。$X$ 上的{\it 亚纯函数层}是与上面
所示预层相伴的 $X_\etale$ 上的层 {\it $\mathcal{K}_X$}。$X$ 上的一个
{\it 亚纯函数}是 $\mathcal{K}_X$ 的整体截面。
\end{definition}

\noindent
由于每个 $\mathcal{S}(U)$ 的每个元素都是 $\mathcal{O}_X(U)$ 中的非零因子，
可见环层的自然映射 $\mathcal{O}_X \to \mathcal{K}_X$ 是单射。此外，由层化
与取茎的相容性，对 $X$ 的任意几何点 $\overline{x}$，有
$$
\mathcal{K}_{X, \overline{x}} =
\mathcal{S}_{\overline{x}}^{-1}\mathcal{O}_{X, \overline{x}}
$$
。集合 $\mathcal{S}_{\overline{x}}$ 是 $\mathcal{O}_{X, \overline{x}}$ 中
非零因子集合的子集，但一般并不等于后者。

\begin{lemma}
\label{spaces-divisors-lemma-describe-S}
设 $S$ 为概形，$X$ 为 $S$ 上的代数空间。若 $U$ 是平展于 $X$ 的仿射概形，
则集合 $\mathcal{S}_X(U)$ 正是 $\mathcal{O}_X(U)$ 中非零因子的集合。
\end{lemma}

\begin{proof}
由引理 \ref{spaces-divisors-lemma-regular-section-structure-sheaf} 得到。
\end{proof}

\noindent
接着，设 $\mathcal{F}$ 为 $X_\etale$ 上的 $\mathcal{O}_X$-模层。考虑预层
$U \mapsto \mathcal{S}(U)^{-1}\mathcal{F}(U)$。它的层化是
$\mathcal{F} \otimes_{\mathcal{O}_X} \mathcal{K}_X$，见《位点上的模》引理
\ref{sites-modules-lemma-simple-invert-module}。

\begin{definition}
\label{spaces-divisors-definition-meromorphic-section}
设 $S$ 为概形，$X$ 为 $S$ 上的代数空间，并设 $\mathcal{F}$ 为 $X_\etale$
上的 $\mathcal{O}_X$-模层。
\begin{enumerate}
% P09-ERR-0187
\item 用 $\mathcal{K}_X(\mathcal{F})$ 表示如下 $\mathcal{K}_X$-模层：它是预层
$U \mapsto \mathcal{S}(U)^{-1}\mathcal{F}(U)$ 的层化。等价地，
$\mathcal{K}_X(\mathcal{F}) =
\mathcal{F} \otimes_{\mathcal{O}_X} \mathcal{K}_X$（见上文）。
\item $\mathcal{F}$ 的一个{\it 亚纯截面}是
$\mathcal{K}_X(\mathcal{F})$ 的整体截面。
\end{enumerate}
\end{definition}

\noindent
特别地，对 $X$ 的任意几何点 $\overline{x}$，有
$$
\mathcal{K}_X(\mathcal{F})_{\overline{x}}
=
\mathcal{F}_{\overline{x}}
\otimes_{\mathcal{O}_{X, \overline{x}}} \mathcal{K}_{X, \overline{x}}
=
\mathcal{S}_{\overline{x}}^{-1}\mathcal{F}_{\overline{x}}
$$
。不过必须谨慎，因为如上所述，$\mathcal{S}_{\overline{x}}$ 未必是平展局部环
$\mathcal{O}_{X, \overline{x}}$ 中非零因子的集合。亚纯截面层一般不是拟凝聚
模，但确实与拟凝聚模共有一些性质。

\begin{lemma}
\label{spaces-divisors-lemma-meromorphic-quasi-coherent}
设 $S$ 为概形，$X$ 为 $S$ 上的代数空间。假设：
\begin{enumerate}
\item[(a)] $X$ 的每个弱伴随点都是余维 $0$ 的点；并且
\item[(b)] $X$ 满足《空间的态射》引理
\ref{spaces-morphisms-lemma-prepare-normalization} 中的等价条件。
\end{enumerate}
则：
\begin{enumerate}
\item $\mathcal{K}_X$ 是拟凝聚 $\mathcal{O}_X$-代数层；
\item 对仿射的 $U \in X_\etale$，$\mathcal{K}_X(U)$ 是
$\mathcal{O}_X(U)$ 的全分式环；
% P09-ERR-0188
\item 对几何点 $\overline{x}$，集合 $\mathcal{S}_{\overline{x}}$ 是
$\mathcal{O}_{X, \overline{x}}$ 中非零因子的集合；并且
\item 对几何点 $\overline{x}$，环 $\mathcal{K}_{X, \overline{x}}$ 是
$\mathcal{O}_{X, \overline{x}}$ 的全分式环。
\end{enumerate}
\end{lemma}

\begin{proof}
% P09-ERR-0189
由引理 \ref{spaces-divisors-lemma-regular-section-structure-sheaf} 可知，对仿射的
$U \in X_\etale$，$\mathcal{S}_X(U) \subset \mathcal{O}_X(U)$ 是
$\mathcal{O}_X(U)$ 中非零因子的集合。因此，预层
$\mathcal{S}^{-1}\mathcal{O}_X$ 在 $X_{affine, \etale}$ 上等于
$$
U \longmapsto Q(\mathcal{O}_X(U))
$$
，记号见《代数》例 \ref{algebra-example-localize-at-prime}。注意，$X$ 的余维
$0$ 点对应于 $U$ 的一般点，见《空间的性质》引理
\ref{spaces-properties-lemma-codimension-0-points}。因此，若
$U = \Spec(A)$，则环 $A$ 只有有限多个极小素理想，并且 $A$ 的每个弱伴随
素理想都是极小的。对 $A$ 的任意平展扩张，同样的结论成立（因为这种扩张的
谱是平展于 $X$ 的仿射概形，因而可在上一句中充当 $A$）。为证明我们的预层
是层且拟凝聚，只需证明，当 $A \to B$ 是平展环同态时，
$$
Q(A) \otimes_A B \longrightarrow Q(B)
$$
是同构，见《空间的性质》引理
\ref{spaces-properties-lemma-characterize-quasi-coherent-small-etale}。（为定义上述
箭头，注意 $A \to B$ 平坦，故它把非零因子映为非零因子。）由《代数》引理
\ref{algebra-lemma-total-ring-fractions-no-embedded-points} 及
\ref{algebra-lemma-weakly-ass-zero-divisors}
% P09-ERR-0190
，有
$$
Q(A) = \prod\nolimits_{\mathfrak p \subset A\text{ 为极小素理想}} A_\mathfrak p
\quad\text{且}\quad
Q(B) = \prod\nolimits_{\mathfrak q \subset B\text{ 为极小素理想}} B_\mathfrak q
$$
。由于 $A \to B$ 平展，$B$ 的极小素理想恰好是位于 $A$ 的极小素理想之上的
$B$ 的素理想（例如由《代数进阶》引理
\ref{more-algebra-lemma-dimension-etale-extension}）。由《代数》引理
\ref{algebra-lemma-local-dimension-zero-henselian}、
\ref{algebra-lemma-characterize-henselian} (13) 及
\ref{algebra-lemma-mop-up}，$A_\mathfrak p \otimes_A B$ 是有限多个局部环的
直积，这些局部环都是 $A_\mathfrak p$ 上的有限平展代数。这显然蕴含
$A_\mathfrak p \otimes_A B =
\prod_{\mathfrak q\text{ 位于 }\mathfrak p} B_\mathfrak q$，正如所需。

\medskip\noindent
至此，我们已经知道 (1) 和 (2) 成立。证明 (3)。设
$s \in \mathcal{O}_{X, \overline{x}}$ 为非零因子。可以找到平展邻域
$(U, \overline{u}) \to (X, \overline{x})$ 以及映到 $s$ 的
$f \in \mathcal{O}_X(U)$。令 $u \in U$ 为 $\overline{u}$ 所确定的点。由于
$\mathcal{O}_{U, u} \to \mathcal{O}_{X, \overline{x}}$ 忠实平坦（后者是前者
的严格亨泽尔化），可见 $f$ 在 $\mathcal{O}_{U, u}$ 中的像是非零因子。由
《除子》引理 \ref{divisors-lemma-meromorphic-weakass-finite}，缩小 $U$ 后可使
$f$ 成为非零因子，因而成为 $\mathcal{S}_X(U)$ 的截面。(4) 通过由 (3) 计算
茎得到。
\end{proof}

\begin{lemma}
\label{spaces-divisors-lemma-meromorphic-quasi-coherent-agree}
设 $S$ 为概形，$X$ 为 $S$ 上的代数空间。假设：
\begin{enumerate}
\item[(a)] $X$ 的每个弱伴随点都是余维 $0$ 的点；
\item[(b)] $X$ 满足《空间的态射》引理
\ref{spaces-morphisms-lemma-prepare-normalization} 中的等价条件；
\item[(c)] $X$ 可由概形 $X_0$ 表示（这是一个别扭但临时使用的记号）。
\end{enumerate}
则亚纯函数层 $\mathcal{K}_X$ 是与拟凝聚亚纯函数层
$\mathcal{K}_{X_0}$ 相伴的拟凝聚 $\mathcal{O}_X$-代数层。
\end{lemma}

\begin{proof}
关于 $\QCoh(\mathcal{O}_X)$ 与 $\QCoh(\mathcal{O}_{X_0})$ 之间的等价，见
《空间的性质》第 \ref{spaces-properties-section-quasi-coherent} 节。本引理成立，
因为由引理 \ref{spaces-divisors-lemma-meromorphic-quasi-coherent} 及《除子》引理
\ref{divisors-lemma-meromorphic-weakass-finite}，$\mathcal{K}_X$ 与
$\mathcal{K}_{X_0}$ 都拟凝聚，并且在 $X$ 与 $X_0$ 的相应仿射开集上取值相同。
\end{proof}

\begin{definition}
\label{spaces-divisors-definition-pullback-meromorphic-sections}
设 $S$ 为概形，$f : X \to Y$ 为 $S$ 上代数空间的态射。若对每个交换图
$$
\xymatrix{
U \ar[r] \ar[d] & X \ar[d] \\
V \ar[r] & Y
}
$$
，其中 $U \in X_\etale$ 且 $V \in Y_\etale$，以及任意截面
$s \in \mathcal{S}_Y(V)$，拉回 $f^\sharp(s) \in \mathcal{O}_X(U)$ 都属于
$\mathcal{S}_X(U)$，则称{\it 亚纯函数沿 $f$ 的拉回有定义}。
\end{definition}

\noindent
在这种情形下，存在诱导映射
$f^\sharp : f_{small}^{-1}\mathcal{K}_Y \to \mathcal{K}_X$；换言之，得到带环
拓扑斯态射的交换图
$$
\xymatrix{
(\Sh(X_\etale), \mathcal{K}_X) \ar[r] \ar[d]^{f_{small}} &
(\Sh(X_\etale), \mathcal{O}_X) \ar[d]^{f_{small}} \\
(\Sh(Y_\etale), \mathcal{K}_Y) \ar[r] &
(\Sh(Y_\etale), \mathcal{O}_Y)
}
$$
对截面 $s \in \Gamma(Y, \mathcal{K}_Y)$，有时记
$f^*(s) = f^\sharp(s)$。

\begin{lemma}
\label{spaces-divisors-lemma-pullback-meromorphic-sections-defined}
% P09-ERR-0191
设 $S$ 为概形，$f : X \to Y$ 为 $S$ 上代数空间的态射。在下列每一种情形中，
亚纯截面的拉回都有定义：
\begin{enumerate}
\item $X$ 的弱伴随点被映到 $Y$ 的余维 $0$ 点；
\item $f$ 是平坦的；
% P09-ERR-0192
\item 按需在此补充。
\end{enumerate}
\end{lemma}

\begin{proof}
经平展局部化，这转化为概形情形，见《除子》引理
\ref{divisors-lemma-pullback-meromorphic-sections-defined}。完成这一转化时，使用
引理 \ref{spaces-divisors-lemma-regular-section-structure-sheaf}（正则截面的描述）、定义
\ref{spaces-divisors-definition-weakly-associated}（弱伴随点的定义）以及《空间的性质》引理
\ref{spaces-properties-lemma-codimension-0-points}（余维 $0$ 点的描述）。
\end{proof}

\begin{lemma}
\label{spaces-divisors-lemma-compute-meromorphic}
设 $S$ 为概形，$X$ 为 $S$ 上的代数空间。假设：
\begin{enumerate}
\item[(a)] $X$ 的每个弱伴随点都是余维 $0$ 的点；
\item[(b)] $X$ 满足《空间的态射》引理
\ref{spaces-morphisms-lemma-prepare-normalization} 中的等价条件；
\item[(c)] $X$ 的每个余维 $0$ 点都可由单态射 $\Spec(k) \to X$ 表示。
\end{enumerate}
令 $X^0 \subset |X|$ 为 $X$ 的余维 $0$ 点所成的集合。则有
$$
\mathcal{K}_X =
\bigoplus\nolimits_{\eta \in X^0} j_{\eta, *}\mathcal{O}_{X, \eta} =
\prod\nolimits_{\eta \in X^0} j_{\eta, *}\mathcal{O}_{X, \eta}
$$
，其中 $j_\eta : \Spec(\mathcal{O}_{X, \eta}) \to X$ 是《概形》第
\ref{schemes-section-points} 节中的典范映射；这是有意义的，因为 $X^0$ 包含在
$X$ 的概形轨迹中。类似地，对每个拟凝聚 $\mathcal{O}_X$-模 $\mathcal{F}$，
得到公式
$$
\mathcal{K}_X(\mathcal{F}) =
\bigoplus\nolimits_{\eta \in X^0} j_{\eta, *}\mathcal{F}_\eta =
\prod\nolimits_{\eta \in X^0} j_{\eta, *}\mathcal{F}_\eta
$$
，它描述 $\mathcal{F}$ 的亚纯截面层。最后，$X$ 的有理函数环就是 $X$ 上的
亚纯函数环；写成公式即 $R(X) = \Gamma(X, \mathcal{K}_X)$。
\end{lemma}

\begin{proof}
由《良态空间》引理 \ref{decent-spaces-lemma-get-reasonable} 及第
\ref{decent-spaces-section-reasonable-decent} 节，$X$ 是良态的
\footnote{反之，若 $X$ 良态，则条件 (c) 自动成立。}。因此，
$X^0 \subset |X|$ 是不可约分支的一般点所成的集合（《良态空间》引理
\ref{decent-spaces-lemma-decent-generic-points}），且由 (b)，$X^0$ 在 $|X|$
中局部有限。于是 $X^0$ 包含在 $|X|$ 的每个稠密开子集中。特别地，$X^0$
包含在概形轨迹中（《良态空间》定理
\ref{decent-spaces-theorem-decent-open-dense-scheme}）。因此，局部环
$\mathcal{O}_{X, \eta}$ 与态射 $j_\eta$ 都有定义。

\medskip\noindent
注意，模层的局部有限直和等于其直积。这个事实连同 $X^0$ 在 $|X|$ 中局部
有限，解释了陈述中直和与直积之间的等式。又由于
$\mathcal{K}_X(\mathcal{F}) =
\mathcal{F} \otimes_{\mathcal{O}_X} \mathcal{K}_X$，第二个等式由第一个等式
得到。

\medskip\noindent
令
$j : Y = \coprod\nolimits_{\eta \in X^0} \Spec(\mathcal{O}_{X, \eta}) \to X$
% P09-ERR-0193
为诸态射 $j_\eta$ 的乘积。我们要证明
$\mathcal{K}_X = j_*\mathcal{O}_Y$。注意，$\mathcal{K}_Y = \mathcal{O}_Y$，
因为 $Y$ 是若干维数为 $0$ 的局部环之谱的不交并：在零维局部环中，任意非零
因子都是单位。接着，由引理
\ref{spaces-divisors-lemma-pullback-meromorphic-sections-defined}，亚纯函数沿 $j$ 的拉回有定义。
这给出映射
$$
\mathcal{K}_X \longrightarrow j_*\mathcal{O}_Y.
$$
设仿射的 $U \in X_\etale$。由引理 \ref{spaces-divisors-lemma-meromorphic-quasi-coherent}，
% P09-ERR-0194
左端的取值是 $\mathcal{O}_X(U)$ 的全分式环。另一方面，右端等于 $U$ 在各余维
$0$ 点（即 $U$ 的一般点）处的局部环之直积。这两个环相等（正如我们已在引理
\ref{spaces-divisors-lemma-meromorphic-quasi-coherent} 的证明中看到的），这是由《代数》引理
\ref{algebra-lemma-total-ring-fractions-no-embedded-points} 及
\ref{algebra-lemma-weakly-ass-zero-divisors} 得到的。因此，我们的映射是同构。

\medskip\noindent
最后，必须证明 $R(X) = \Gamma(X, \mathcal{K}_X)$。把概形情形（《除子》引理
\ref{divisors-lemma-meromorphic-weakass-finite}）应用于概形轨迹
$X' \subset X$ 即得。确切地说，由于它是 $X$ 的稠密开子空间（见上文），
按照定义，$X$ 的有理函数环与 $X'$ 的有理函数环相同。当把 $X'$ 看作概形时，
$R(X')$ 当然与其有理函数环一致。另一方面，由上面对 $\mathcal{K}_X$ 的描述，
以及上面所见的 $X^0 \subset |X'|$ 包含在每个稠密开集中的事实，可知
$\Gamma(X, \mathcal{K}_X) = \Gamma(X', \mathcal{K}_{X'})$。最后使用引理
\ref{spaces-divisors-lemma-meromorphic-quasi-coherent-agree} 中记录的相容性。
\end{proof}

\begin{definition}
\label{spaces-divisors-definition-regular-meromorphic-section}
设 $S$ 为概形，$X$ 为 $S$ 上的代数空间，并设 $\mathcal{L}$ 为可逆
$\mathcal{O}_X$-模。若诱导映射
$\mathcal{K}_X \to \mathcal{K}_X(\mathcal{L})$ 是单射，则称 $\mathcal{L}$ 的
亚纯截面 $s$ 是{\it 正则的}。
\end{definition}

\noindent
下面详细说明（正则）亚纯截面何时可以拉回。

\begin{lemma}
\label{spaces-divisors-lemma-meromorphic-sections-pullback}
设 $S$ 为概形，$f : X \to Y$ 为 $S$ 上代数空间的态射。假设亚纯函数沿 $f$
的拉回有定义（见定义 \ref{spaces-divisors-definition-pullback-meromorphic-sections}）。
\begin{enumerate}
\item 设 $\mathcal{F}$ 为 $\mathcal{O}_Y$-模层。存在亚纯截面的典范拉回映射
$f^* : \Gamma(Y, \mathcal{K}_Y(\mathcal{F})) \to
\Gamma(X, \mathcal{K}_X(f^*\mathcal{F}))$
，其定义在 $\mathcal{F}$ 的亚纯截面上。
% P09-ERR-0195
\item 设 $\mathcal{L}$ 为可逆 $\mathcal{O}_X$-模。$\mathcal{L}$ 的正则亚纯
截面 $s$ 拉回为 $f^*\mathcal{L}$ 的正则亚纯截面 $f^*s$。
\end{enumerate}
\end{lemma}

\begin{proof}
略。
\end{proof}

\begin{lemma}
\label{spaces-divisors-lemma-regular-meromorphic-section-exists}
设 $S$ 为概形，$X$ 为 $S$ 上的代数空间，且满足引理
\ref{spaces-divisors-lemma-compute-meromorphic} 的 (a)、(b)、(c)。则每个可逆
$\mathcal{O}_X$-模 $\mathcal{L}$ 都有正则亚纯截面。
\end{lemma}

\begin{proof}
使用引理 \ref{spaces-divisors-lemma-compute-meromorphic} 的记号。对所有
% P09-ERR-0196
$\eta \in X^0$，$\mathcal{L}$ 的茎 $\mathcal{L}_\eta$ 都有定义，并且它是秩
$1$ 的自由 $\mathcal{O}_{X, \eta}$-模。对每个 $\eta \in X^0$ 选取生成元
$s_\eta \in \mathcal{L}_\eta$。由引理 \ref{spaces-divisors-lemma-compute-meromorphic} 中对
$\mathcal{K}_X$ 与 $\mathcal{K}_X(\mathcal{L})$ 的描述，立即可知
$s = \prod s_\eta$ 是 $\mathcal{L}$ 的正则亚纯截面。
\end{proof}












\section{相对 Proj}
\label{spaces-divisors-section-relative-proj}

\noindent
本节在代数空间的框架下重新讨论相对 Proj 的构造。本节内容对应于概形情形下
《构造》第 \ref{constructions-section-relative-proj} 节及《除子》第
\ref{divisors-section-relative-proj} 节的内容。

\begin{situation}
\label{spaces-divisors-situation-relative-proj}
这里，$S$ 是概形，$X$ 是 $S$ 上的代数空间，而 $\mathcal{A}$ 是拟凝聚分次
$\mathcal{O}_X$-代数。
\end{situation}

\noindent
在情形 \ref{spaces-divisors-situation-relative-proj} 中，我们将定义函子
$F : (\Sch/S)_{fppf}^{opp} \to \textit{Sets}$，它最终将是一个代数空间。
我们作适当修改后沿用《构造》第
\ref{constructions-section-relative-proj} 节的步骤。首先，给定 $S$ 上的概形
$T$，定义一个{\it $T$ 上的四元组}为系统
$(d, f : T \to X, \mathcal{L}, \psi)$
\begin{enumerate}
\item $d \geq 1$ 是整数；
\item $f : T \to X$ 是 $S$ 上的态射；
\item $\mathcal{L}$ 是可逆 $\mathcal{O}_T$-模；并且
\item
$\psi : f^*\mathcal{A}^{(d)} \to \bigoplus_{n \geq 0}\mathcal{L}^{\otimes n}$
是分次 $\mathcal{O}_T$-代数同态，且
$f^*\mathcal{A}_d \to \mathcal{L}$ 是满的。
\end{enumerate}
称两个四元组 $(d, f, \mathcal{L}, \psi)$ 与
$(d', f', \mathcal{L}', \psi')$ {\it 等价}\footnote{此定义受到《构造》引理
\ref{constructions-lemma-equivalent-relative} 的启发。采用这一定义的优点是，
它显然定义了一个等价关系。}，当且仅当 $f = f'$，并且对某个正整数
$m = ad = a'd'$，存在同构
$\beta : \mathcal{L}^{\otimes a} \to (\mathcal{L}')^{\otimes a'}$
，使得 $\beta \circ \psi|_{f^*\mathcal{A}^{(m)}}$ 与
$\psi'|_{f^*\mathcal{A}^{(m)}}$ 作为如下分次环同态相同：
% P09-ERR-0197
$f^*\mathcal{A}^{(m)} \to \bigoplus_{n \geq 0} (\mathcal{L}')^{\otimes mn}$。
给定四元组 $(d, f, \mathcal{L}, \psi)$ 与态射 $h : T' \to T$，其拉回为
$(d, f \circ h, h^*\mathcal{L}, h^*\psi)$。拉回保持此等价关系。最后，对 $S$
上的{\it 拟紧}概形 $T$，令
$$
F(T) = \text{四元组等价类的集合，其基底为 }T
$$
；对 $S$ 上任意概形 $T$，令
$$
F(T)
=
\lim_{V \subset T\text{ 的拟紧开子集}} F(V).
$$
换言之，$F(T)$ 的元素 $\xi$ 对应于一族相容的元素选择
$\xi_V \in F(V)$，其中 $V$ 遍历 $T$ 的拟紧开子集。因此，我们定义了函子
\begin{equation}
\label{spaces-divisors-equation-proj}
% P09-ERR-0198
F : \Sch^{opp} \longrightarrow \textit{Sets}
\end{equation}
存在函子态射 $F \to X$，它把四元组 $(d, f, \mathcal{L}, \psi)$ 映到 $f$。

\begin{lemma}
\label{spaces-divisors-lemma-relative-proj}
% P09-ERR-0199
在情形 \ref{spaces-divisors-situation-relative-proj} 中，上述函子 $F$ 是代数空间。对任意态射
$g : Z \to X$，其中 $Z$ 是概形，存在与进一步基变换相容的典范同构
$\underline{\text{Proj}}_Z(g^*\mathcal{A}) = Z \times_X F$
。
\end{lemma}

\begin{proof}
只需证明第二个断言，见《空间》引理
\ref{spaces-lemma-representable-over-space}。设 $g : Z \to X$ 为态射，其中 $Z$
是概形。令 $F'$ 为与分次拟凝聚 $\mathcal{O}_Z$-代数
$g^*\mathcal{A}$ 相伴的四元组函子。于是存在典范同构
$F' = Z \times_X F$，它把 $F'$ 的四元组
$(d, f : T \to Z, \mathcal{L}, \psi)$ 映到
$(d, g \circ f, \mathcal{L}, \psi)$（略去细节，见《构造》引理
\ref{constructions-lemma-proj-base-change} 的证明）。由《构造》诸引理
\ref{constructions-lemma-equivalent-relative}、
\ref{constructions-lemma-relative-proj}、
% P09-ERR-0200
\ref{constructions-lemma-glueing-gives-functor-proj} 以及定义
\ref{constructions-definition-relative-proj}
可知，$F'$ 由 $\underline{\text{Proj}}_Z(g^*\mathcal{A})$ 表示。
\end{proof}

\noindent
上述引理说明以下定义是有意义的。

\begin{definition}
\label{spaces-divisors-definition-relative-proj}
设 $S$ 为概形，$X$ 为 $S$ 上的代数空间，并设 $\mathcal{A}$ 为拟凝聚分次
$\mathcal{O}_X$-代数层。$\mathcal{A}$ 在 $X$ 上的{\it 相对齐次谱}，或
$\mathcal{A}$ 在 $X$ 上的{\it 齐次谱}，或 $\mathcal{A}$ 在 $X$ 上的
{\it 相对 Proj}，是引理 \ref{spaces-divisors-lemma-relative-proj} 中 $X$ 上的代数空间 $F$。
将其记作 $\pi : \underline{\text{Proj}}_X(\mathcal{A}) \to X$。
\end{definition}

\noindent
特别地，由构造，相对 Proj 的结构态射是可表的。也可以通过黏合来理解相对
Proj。设 $\varphi : U \to X$ 为满平展态射，其中 $U$ 是概形。令
$R = U \times_X U$，其投影态射为 $s, t : R  \to U$。由引理
\ref{spaces-divisors-lemma-relative-proj}，存在 $U$ 上的典范同构
$$
\gamma : 
\underline{\text{Proj}}_U(\varphi^*\mathcal{A})
\longrightarrow
\underline{\text{Proj}}_X(\mathcal{A}) \times_X U
$$
。令 $\alpha : t^*\varphi^*\mathcal{A} \to s^*\varphi^*\mathcal{A}$ 为
《空间的性质》命题
\ref{spaces-properties-proposition-quasi-coherent} 中的典范同构。则图
$$
\xymatrix{
&
\underline{\text{Proj}}_U(\varphi^*\mathcal{A}) \times_{U, s} R
\ar@{=}[r] &
\underline{\text{Proj}}_R(s^*\varphi^*\mathcal{A})
\ar[dd]_{\text{诱导自 }\alpha} \\
\underline{\text{Proj}}_X(\mathcal{A}) \times_X R
\ar[ru]_{s^*\gamma} \ar[rd]^{t^*\gamma} \\
&
\underline{\text{Proj}}_U(\varphi^*\mathcal{A}) \times_{U, t} R
\ar@{=}[r] &
\underline{\text{Proj}}_R(t^*\varphi^*\mathcal{A})
}
$$
交换（等号来自《构造》引理
\ref{constructions-lemma-relative-proj-base-change}）。因此，若用
$\mathcal{A}_U$、$\mathcal{A}_R$ 表示 $\mathcal{A}$ 到 $U$、$R$ 的拉回，
% P09-ERR-0201
则 $P = \underline{\text{Proj}}_X(\mathcal{A})$ 有由概形
$P_U = \underline{\text{Proj}}_U(\mathcal{A}_U)$ 给出的平展覆盖，并且
$P_U \times_P P_U$ 等于
$P_R = \underline{\text{Proj}}_R(\mathcal{A}_R)$。利用这些说明，可以按通常
方式使用平展局部化，把概形情形中相对 Proj 的结果转移到代数空间情形。

\begin{lemma}
\label{spaces-divisors-lemma-twists-of-structure-sheaf}
在情形 \ref{spaces-divisors-situation-relative-proj} 中，相对 Proj 带有拟凝聚
$\mathbf{Z}$-分次代数层
$\bigoplus_{n \in \mathbf{Z}}
\mathcal{O}_{\underline{\text{Proj}}_X(\mathcal{A})}(n)$
及分次代数的典范同态
$$
\psi :
\pi^*\mathcal{A}
\longrightarrow
\bigoplus\nolimits_{n \geq 0}
\mathcal{O}_{\underline{\text{Proj}}_X(\mathcal{A})}(n)
$$
；它们基变换到 $X$ 上任意概形后，与《构造》引理
\ref{constructions-lemma-glue-relative-proj-twists} 相符。
\end{lemma}

\begin{proof}
如定义 \ref{spaces-divisors-definition-relative-proj} 后的讨论，选取概形 $U$ 及满平展态射
$U \to X$，令 $R = U \times_X U$，其投影为 $s, t : R \to U$，并令
$\mathcal{A}_U = \mathcal{A}|_U$、$\mathcal{A}_R = \mathcal{A}|_R$，以及
$\pi : P = \underline{\text{Proj}}_X(\mathcal{A}) \to X$、
$\pi_U : P_U = \underline{\text{Proj}}_U(\mathcal{A}_U)$ 以及
% P09-ERR-0202
$\pi_R : P_R = \underline{\text{Proj}}_U(\mathcal{A}_R)$。
% P09-ERR-0203
由《构造》引理 \ref{constructions-lemma-glue-relative-proj-twists}，有拟凝聚
$\mathbf{Z}$-分次 $\mathcal{O}_{P_U}$-代数层
$\bigoplus_{n \in \mathbf{Z}} \mathcal{O}_{P_U}(n)$
以及典范映射
$\psi_U : \pi_U^*\mathcal{A}_U \to \bigoplus_{n \geq 0} \mathcal{O}_{P_U}(n)$
；对 $P_R$ 亦然。由《构造》引理
\ref{constructions-lemma-relative-proj-base-change}，经任一投影
$P_R \to P_U$ 拉回 $\mathcal{O}_{P_U}(n)$ 与 $\psi_U$，分别得到
$\mathcal{O}_{P_R}(n)$ 与 $\psi_R$。由《空间的性质》命题
\ref{spaces-properties-proposition-quasi-coherent}
得到 $\mathcal{O}_{P}(n)$ 与 $\psi$。略去其与拉回到 $X$ 上任意概形相容的验证。
\end{proof}

\noindent
构造出相对 Proj 后，下面讨论它的一些基本性质。

\begin{lemma}
\label{spaces-divisors-lemma-relative-proj-base-change}
设 $S$ 为概形，$g : X' \to X$ 为 $S$ 上代数空间的态射，并设
$\mathcal{A}$ 为拟凝聚分次 $\mathcal{O}_X$-代数层。则存在典范同构
$$
r :
\underline{\text{Proj}}_{X'}(g^*\mathcal{A})
\longrightarrow
X' \times_X \underline{\text{Proj}}_X(\mathcal{A})
$$
以及相应的同构
$$
\theta :
r^*\text{pr}_2^*\left(\bigoplus\nolimits_{d \in \mathbf{Z}}
\mathcal{O}_{\underline{\text{Proj}}_X(\mathcal{A})}(d)\right)
\longrightarrow
\bigoplus\nolimits_{d \in \mathbf{Z}}
\mathcal{O}_{\underline{\text{Proj}}_{X'}(g^*\mathcal{A})}(d)
$$
，它是 $\mathbf{Z}$-分次
$\mathcal{O}_{\underline{\text{Proj}}_{X'}(g^*\mathcal{A})}$-代数的同构。
\end{lemma}

\begin{proof}
令 $F$ 为函子 (\ref{spaces-divisors-equation-proj})，并令 $F'$ 为在 $X'$ 上用
$g^*\mathcal{A}$ 定义的相应函子。我们断言存在函子的典范同构
$r : F' \to X' \times_X F$（当然，$r$ 正是本引理中的同构）。只需对 $S$ 上的
拟紧概形 $T$ 构造双射
$r : F'(T) \to X'(T) \times_{X(T)} F(T)$。首先，若
$\xi = (d', f', \mathcal{L}', \psi')$ 是 $F'$ 的一个 $T$ 上四元组，则可令
$r(\xi) = (f', (d', g \circ f', \mathcal{L}', \psi'))$。这是有意义的，因为
% P09-ERR-0204
$(g \circ f')^*\mathcal{A}^{(d)} = (f')^*(g^*\mathcal{A})^{(d)}$。
逆映射把配对 $(f', (d, f, \mathcal{L}, \psi))$ 映到四元组
$(d, f', \mathcal{L}, \psi)$。略去最后一个断言的证明（提示：通过平展局部化
约化到概形情形，并应用《构造》引理
\ref{constructions-lemma-relative-proj-base-change}）。
\end{proof}

\begin{lemma}
\label{spaces-divisors-lemma-relative-proj-separated}
在情形 \ref{spaces-divisors-situation-relative-proj} 中，态射
$\pi : \underline{\text{Proj}}_X(\mathcal{A}) \to X$
是分离的。
\end{lemma}

\begin{proof}
由《空间的态射》引理 \ref{spaces-morphisms-lemma-separated-local} 及相对 Proj
的构造，这归结为概形情形，即《构造》引理
\ref{constructions-lemma-relative-proj-separated}。
\end{proof}

\begin{lemma}
\label{spaces-divisors-lemma-relative-proj-quasi-compact}
在情形 \ref{spaces-divisors-situation-relative-proj} 中，若下列条件之一成立：
\begin{enumerate}
\item $\mathcal{A}$ 作为 $\mathcal{A}_0$-代数层是有限型的；
\item $\mathcal{A}$ 作为 $\mathcal{A}_0$-代数由 $\mathcal{A}_1$ 生成，且
$\mathcal{A}_1$ 是有限型 $\mathcal{A}_0$-模；
\item 存在有限型拟凝聚 $\mathcal{A}_0$-子模
$\mathcal{F} \subset \mathcal{A}_{+}$，使得
$\mathcal{A}_{+}/\mathcal{F}\mathcal{A}$ 是
$\mathcal{A}/\mathcal{F}\mathcal{A}$ 的局部幂零理想层；
\end{enumerate}
则 $\pi : \underline{\text{Proj}}_X(\mathcal{A}) \to X$ 是拟紧的。
\end{lemma}

\begin{proof}
由《空间的态射》引理
\ref{spaces-morphisms-lemma-quasi-compact-local} 及相对 Proj 的构造，这归结为
概形情形，即《除子》引理
\ref{divisors-lemma-relative-proj-quasi-compact}。
\end{proof}

\begin{lemma}
\label{spaces-divisors-lemma-relative-proj-finite-type}
在情形 \ref{spaces-divisors-situation-relative-proj} 中，若 $\mathcal{A}$ 作为
$\mathcal{O}_X$-代数层是有限型的，则
$\pi : \underline{\text{Proj}}_X(\mathcal{A}) \to X$ 是有限型的。
\end{lemma}

\begin{proof}
由《空间的态射》引理
\ref{spaces-morphisms-lemma-finite-type-local} 及相对 Proj 的构造，这归结为
概形情形，即《除子》引理
\ref{divisors-lemma-relative-proj-finite-type}。
\end{proof}

\begin{lemma}
\label{spaces-divisors-lemma-relative-proj-universally-closed}
在情形 \ref{spaces-divisors-situation-relative-proj} 中，若
$\mathcal{O}_X \to \mathcal{A}_0$
是整代数映射\footnote{换言之，$\mathcal{O}_X$ 在 $\mathcal{A}_0$ 中的整闭包
（见《空间的态射》定义
\ref{spaces-morphisms-definition-integral-closure}）等于
$\mathcal{A}_0$。}，且 $\mathcal{A}$ 作为 $\mathcal{A}_0$-代数是有限型的，则
$\pi : \underline{\text{Proj}}_X(\mathcal{A}) \to X$ 是普遍闭的。
\end{lemma}

\begin{proof}
由《空间的态射》引理
\ref{spaces-morphisms-lemma-universally-closed-local} 及相对 Proj 的构造，这归结为
概形情形，即《除子》引理
\ref{divisors-lemma-relative-proj-universally-closed}。
\end{proof}

\begin{lemma}
\label{spaces-divisors-lemma-relative-proj-proper}
在情形 \ref{spaces-divisors-situation-relative-proj} 中，下列条件等价：
\begin{enumerate}
\item $\mathcal{A}_0$ 是有限型 $\mathcal{O}_X$-模，且 $\mathcal{A}$ 作为
$\mathcal{A}_0$-代数是有限型的；
\item $\mathcal{A}_0$ 是有限型 $\mathcal{O}_X$-模，且 $\mathcal{A}$ 作为
$\mathcal{O}_X$-代数是有限型的。
\end{enumerate}
若这些条件成立，则
$\pi : \underline{\text{Proj}}_X(\mathcal{A}) \to X$
是固有的。
\end{lemma}

\begin{proof}
由《空间的态射》引理
\ref{spaces-morphisms-lemma-proper-local} 及相对 Proj 的构造，这归结为概形
情形，即《除子》引理
% P09-ERR-0205
\ref{divisors-lemma-relative-proj-universally-closed}。
\end{proof}

\begin{lemma}
\label{spaces-divisors-lemma-relative-proj-generated-in-degree-1}
设 $S$ 为概形，$X$ 为 $S$ 上的代数空间。设 $\mathcal{A}$ 为拟凝聚分次
$\mathcal{O}_X$-模层，
% P09-ERR-0206
并由 $\mathcal{A}_1$ 作为 $\mathcal{A}_0$-代数生成。令
$P = \underline{\text{Proj}}_X(\mathcal{A})$，则：
\begin{enumerate}
\item $P$ 表示函子 $F_1$；该函子把 $S$ 上的 $T$ 映到三元组
$(f, \mathcal{L}, \psi)$ 的同构类所成的集合，其中 $f : T \to X$ 是 $S$ 上
的态射，$\mathcal{L}$ 是可逆 $\mathcal{O}_T$-模，而
$\psi : f^*\mathcal{A} \to \bigoplus_{n \geq 0} \mathcal{L}^{\otimes n}$
是分次 $\mathcal{O}_T$-代数映射，并诱导满射
$f^*\mathcal{A}_1 \to \mathcal{L}$；
\item 典范映射 $\pi^*\mathcal{A}_1 \to \mathcal{O}_P(1)$ 是满的；并且
\item 每个 $\mathcal{O}_P(n)$ 都可逆，且乘法映射诱导同构
$\mathcal{O}_P(n) \otimes_{\mathcal{O}_P} \mathcal{O}_P(m) =
\mathcal{O}_P(n + m)$。
\end{enumerate}
\end{lemma}

\begin{proof}
略。概形情形见《构造》引理
\ref{constructions-lemma-apply-relative}。
\end{proof}







\section{相对 Proj 的函子性}
\label{spaces-divisors-section-functoriality-relative-proj}

\noindent
本节对应于《构造》第
\ref{constructions-section-functoriality-relative-proj} 节。

\begin{lemma}
\label{spaces-divisors-lemma-morphism-relative-proj}
设 $S$ 为概形，$X$ 为 $S$ 上的代数空间。设
$\psi : \mathcal{A} \to \mathcal{B}$ 为拟凝聚分次
$\mathcal{O}_X$-代数的映射。令
$P = \underline{\text{Proj}}_X(\mathcal{A}) \to X$ 以及
$Q = \underline{\text{Proj}}_X(\mathcal{B}) \to X$。
存在典范开子空间 $U(\psi) \subset Q$ 以及 $X$ 上代数空间的典范态射
$$
r_\psi :
U(\psi)
\longrightarrow
P
$$
和 $\mathbf{Z}$-分次 $\mathcal{O}_{U(\psi)}$-代数的映射
$$
\theta = \theta_\psi :
r_\psi^*\left(
\bigoplus\nolimits_{d \in \mathbf{Z}} \mathcal{O}_P(d)
\right)
\longrightarrow
\bigoplus\nolimits_{d \in \mathbf{Z}} \mathcal{O}_{U(\psi)}(d).
$$
三元组 $(U(\psi), r_\psi, \theta)$ 由下述性质刻画：对任意平展于 $X$ 的概形
$W$，三元组
$$
(U(\psi) \times_X W,\quad
r_\psi|_{U(\psi) \times_X W} : U(\psi) \times_X W \to  P \times_X W,\quad
\theta|_{U(\psi) \times_X W})
$$
等于《构造》引理 \ref{constructions-lemma-morphism-relative-proj} 中与
$\psi : \mathcal{A}|_W \to \mathcal{B}|_W$ 相伴的三元组。
\end{lemma}

\begin{proof}
本引理由平展局部化及概形情形得到，见定义
\ref{spaces-divisors-definition-relative-proj} 后的讨论。略去细节。
\end{proof}

\begin{lemma}
\label{spaces-divisors-lemma-morphism-relative-proj-transitive}
设 $S$ 为概形，$X$ 为 $S$ 上的代数空间。设 $\mathcal{A}$、$\mathcal{B}$、
$\mathcal{C}$ 为拟凝聚分次 $\mathcal{O}_X$-代数。令
$P = \underline{\text{Proj}}_X(\mathcal{A})$、
$Q = \underline{\text{Proj}}_X(\mathcal{B})$ 及
$R = \underline{\text{Proj}}_X(\mathcal{C})$。设
$\varphi : \mathcal{A} \to \mathcal{B}$、
$\psi : \mathcal{B} \to \mathcal{C}$ 为分次 $\mathcal{O}_X$-代数映射。则有
% P09-ERR-0207
$$
U(\psi \circ \varphi) = r_\varphi^{-1}(U(\psi))
\quad
\text{且}
\quad
r_{\psi \circ \varphi}
=
r_\varphi \circ r_\psi|_{U(\psi \circ \varphi)}.
$$
此外有
$$
\theta_\psi \circ r_\psi^*\theta_\varphi
=
\theta_{\psi \circ \varphi}
$$
，其中采用显然的记号。
\end{lemma}

\begin{proof}
略。
\end{proof}

\begin{lemma}
\label{spaces-divisors-lemma-surjective-graded-rings-map-relative-proj}
% P09-ERR-0209
沿用上述引理 \ref{spaces-divisors-lemma-morphism-relative-proj} 的假设与记号。假设当
$d \gg 0$ 时，$\mathcal{A}_d \to \mathcal{B}_d$ 是满的。则：
\begin{enumerate}
\item $U(\psi) = Q$；
% P09-ERR-0208
\item $r_\psi : Q \to R$ 是闭浸入；并且
\item 映射 $\theta : r_\psi^*\mathcal{O}_P(n) \to \mathcal{O}_Q(n)$
是满的，但一般不是同构（即使
$\mathcal{A} \to \mathcal{B}$ 是满的）。
\end{enumerate}
\end{lemma}

\begin{proof}
通过平展局部化，这由概形情形（《构造》引理
\ref{constructions-lemma-surjective-graded-rings-map-relative-proj}）得到。
\end{proof}

\begin{lemma}
\label{spaces-divisors-lemma-eventual-iso-graded-rings-map-relative-proj}
% P09-ERR-0209
沿用上述引理 \ref{spaces-divisors-lemma-morphism-relative-proj} 的假设与记号。假设对所有
$d \gg 0$，$\mathcal{A}_d \to \mathcal{B}_d$ 都是同构。则：
\begin{enumerate}
\item $U(\psi) = Q$；
\item $r_\psi : Q \to P$ 是同构；并且
\item 映射 $\theta : r_\psi^*\mathcal{O}_P(n) \to \mathcal{O}_Q(n)$
都是同构。
\end{enumerate}
\end{lemma}

\begin{proof}
通过平展局部化，这由概形情形（《构造》引理
\ref{constructions-lemma-eventual-iso-graded-rings-map-relative-proj}）得到。
\end{proof}

\begin{lemma}
\label{spaces-divisors-lemma-surjective-generated-degree-1-map-relative-proj}
% P09-ERR-0209
沿用上述引理 \ref{spaces-divisors-lemma-morphism-relative-proj} 的假设与记号。假设当
$d \gg 0$ 时，$\mathcal{A}_d \to \mathcal{B}_d$ 是满的，并且
$\mathcal{A}$ 在 $\mathcal{A}_0$ 上由 $\mathcal{A}_1$ 生成。则：
\begin{enumerate}
\item $U(\psi) = Q$；
\item $r_\psi : Q \to P$ 是闭浸入；并且
\item 映射 $\theta : r_\psi^*\mathcal{O}_P(n) \to \mathcal{O}_Q(n)$
都是同构。
\end{enumerate}
\end{lemma}

\begin{proof}
通过平展局部化，这由概形情形（《构造》引理
\ref{constructions-lemma-surjective-generated-degree-1-map-relative-proj}）得到。
\end{proof}










\section{可逆层与到相对 Proj 的态射}
\label{spaces-divisors-section-invertible-relative-proj}

\noindent
看来我们可能会在某处用到下面这个引理。
情形如下：
\begin{enumerate}
\item 设 $S$ 为概形，$Y$ 为 $S$ 上的代数空间。
\item 设 $\mathcal{A}$ 为拟凝聚分次 $\mathcal{O}_Y$-代数。
\item 以 $\pi : \underline{\text{Proj}}_Y(\mathcal{A}) \to Y$ 表示
$\mathcal{A}$ 在 $Y$ 上的相对 Proj。
\item 设 $f : X \to Y$ 为 $S$ 上代数空间的态射。
\item 设 $\mathcal{L}$ 为可逆 $\mathcal{O}_X$-模。
\item 设
$\psi : f^*\mathcal{A} \to \bigoplus_{d \geq 0} \mathcal{L}^{\otimes d}$
为分次 $\mathcal{O}_X$-代数的同态。
\end{enumerate}
给定这些数据，令 $U(\psi) \subset X$ 为满足下式的开子空间：
$$
|U(\psi)| = \bigcup\nolimits_{d \geq 1}
\{\text{满足下述条件的轨迹：}f^*\mathcal{A}_d \to \mathcal{L}^{\otimes d}
\text{ 为满射}\}
$$
$U(\psi) \subset X$ 的构造与沿任意态射 $X' \to X$ 的拉回可交换。

\begin{lemma}
\label{spaces-divisors-lemma-invertible-map-into-relative-proj}
% P09-ERR-0210
采用上述假设与记号。态射 $\psi$ 诱导一个 $Y$ 上代数空间的典范态射
$$
r_{\mathcal{L}, \psi} :
U(\psi) \longrightarrow \underline{\text{Proj}}_Y(\mathcal{A})
$$
，并伴随一个分次 $\mathcal{O}_{U(\psi)}$-代数映射
$$
\theta :
r_{\mathcal{L}, \psi}^*\left(
\bigoplus\nolimits_{d \geq 0}
\mathcal{O}_{\underline{\text{Proj}}_Y(\mathcal{A})}(d)
\right)
\longrightarrow
\bigoplus\nolimits_{d \geq 0} \mathcal{L}^{\otimes d}|_{U(\psi)}
$$
，它由下列性质刻画：
\begin{enumerate}
\item 对平展态射 $V \to Y$ 和 $d \geq 0$，下图
% P09-ERR-0211
$$
\xymatrix{
\mathcal{A}_d(V) \ar[d]_{\psi} \ar[r]_{\psi} &
\Gamma(V \times_Y X, \mathcal{L}^{\otimes d}) \ar[d]^{restrict} \\
\Gamma(V \times_Y \underline{\text{Proj}}_Y(\mathcal{A}),
\mathcal{O}_{\underline{\text{Proj}}_Y(\mathcal{A})}(d)) \ar[r]^-\theta &
\Gamma(V \times_Y U(\psi), \mathcal{L}^{\otimes d})
}
$$
可交换。
\item 对任意 $d \geq 1$ 以及任意态射 $W \to X$，其中 $W$ 为概形，若
$\psi|_W : f^*\mathcal{A}_d|_W \to \mathcal{L}^{\otimes d}|_W$ 为满射，
则 (a) $W \to X$ 分解经过 $U(\psi)$，且 (b) $W \to U(\psi)$ 与
$r_{\mathcal{L}, \psi}$ 的复合等同于态射
$W \to \underline{\text{Proj}}_Y(\mathcal{A})$；后者因
$\underline{\text{Proj}}_Y(\mathcal{A})$ 的构造而存在，见定义
\ref{spaces-divisors-definition-relative-proj}。
\item 考虑可交换图
$$
\xymatrix{
X' \ar[r]_{g'} \ar[d]_{f'} & X \ar[d]^f \\
Y' \ar[r]^g & Y
}
$$
，其中 $X'$ 与 $Y'$ 为概形。置 $\mathcal{A}' = g^*\mathcal{A}$ 及
$\mathcal{L}' = (g')^*\mathcal{L}$，并以
% P09-ERR-0212
$\psi' : (f')^*\mathcal{A} \to \bigoplus_{d \geq 0} (\mathcal{L}')^{\otimes d}$
表示 $\psi$ 的拉回。令 $U(\psi')$、
% P09-ERR-0213
$r_{\psi', \mathcal{L}'}$ 及 $\theta'$ 为
% P09-ERR-0214
《构造》引理 \ref{spaces-divisors-lemma-invertible-map-into-relative-proj} 中所构造的开子空间、
态射及同态。那么 $U(\psi') = (g')^{-1}(U(\psi))$，且在引理
\ref{spaces-divisors-lemma-relative-proj-base-change} 的同构
$\underline{\text{Proj}}_{Y'}(\mathcal{A}') =
Y' \times_Y \underline{\text{Proj}}_Y(\mathcal{A})$
下，$r_{\psi', \mathcal{L}'}$ 等同于 $r_{\psi, \mathcal{L}}$ 的基变换。
此外，$\theta'$ 是 $\theta$ 的拉回。
\end{enumerate}
\end{lemma}

\begin{proof}
略。提示：首先注意，对 $X$ 上的拟紧概形 $W$，下列条件等价：
\begin{enumerate}
\item $W \to X$ 分解经过 $U(\psi)$；以及
% P09-ERR-0215
\item 存在 $d$，使得
$\psi|_W : f^*\mathcal{A}_d|_W \to \mathcal{L}^{\otimes d}|_W$
为满射。
\end{enumerate}
这就把 $U(\psi)$ 描述为基范畴 $(\Sch/S)_{fppf}$ 上 $X$ 的子函子。
对这样的 $W$ 与 $d$，考虑四元组
$(d, W \to Y, \mathcal{L}|_W, \psi^{(d)}|_W)$。
由 $\underline{\text{Proj}}_Y(\mathcal{A})$ 的定义，得到态射
$W \to \underline{\text{Proj}}_Y(\mathcal{A})$。根据四元组的等价概念可见，
此态射与 $d$ 的选取无关。这显然定义了函子之间的变换
$r_{\psi, \mathcal{L}} : U(\psi) \to \underline{\text{Proj}}_Y(\mathcal{A})$，
即代数空间的态射。按构造，此态射满足 (2)。由于《构造》引理
\ref{constructions-lemma-invertible-map-into-relative-proj} 中构造的态射满足同一
性质，故 (3) 成立。

\medskip\noindent
为构造 $\theta$ 并验证引理中相容性 (1)，在 $Y$ 与 $X$ 上平展局部地工作，
论证同定义 \ref{spaces-divisors-definition-relative-proj} 后的讨论。
\end{proof}








\section{相对丰沛层}
\label{spaces-divisors-section-relatively-ample}

\noindent
本节是《态射》第 \ref{morphisms-section-relatively-ample} 节在代数空间情形下的
对应版本。我们如下定义相对丰沛可逆层。

\begin{definition}
\label{spaces-divisors-definition-relatively-ample}
设 $S$ 为概形，$f : X \to Y$ 为 $S$ 上代数空间的态射，
$\mathcal{L}$ 为可逆 $\mathcal{O}_X$-模。若 $f : X \to Y$ 可表，且对每个
从概形 $Z$ 出发的态射 $Z \to Y$，$\mathcal{L}$ 到
$X_Z = Z \times_Y X$ 上的拉回 $\mathcal{L}_Z$ 按照《态射》定义
\ref{morphisms-definition-relatively-ample} 在 $X_Z/Z$ 上丰沛，则称
$\mathcal{L}$ {\it 相对丰沛}，或称为 {\it $f$-相对丰沛}、
{\it 在 $X/Y$ 上丰沛}、或 {\it $f$-丰沛}。
\end{definition}

\noindent
我们几乎总会把有关相对丰沛可逆层的问题归约到概形情形。因此，本节主要包含
一些基本检验。

\begin{lemma}
\label{spaces-divisors-lemma-relatively-ample-sanity-check}
设 $S$ 为概形，$f : X \to Y$ 为 $S$ 上代数空间的态射，
$\mathcal{L}$ 为可逆 $\mathcal{O}_X$-模。假设 $Y$ 是概形。下列条件等价：
\begin{enumerate}
\item $\mathcal{L}$ 在定义 \ref{spaces-divisors-definition-relatively-ample} 的意义下在
$X/Y$ 上丰沛；以及
\item $X$ 是概形，且 $\mathcal{L}$ 在《态射》定义
\ref{morphisms-definition-relatively-ample} 的意义下在 $X/Y$ 上丰沛。
\end{enumerate}
\end{lemma}

\begin{proof}
这由定义及《态射》引理 \ref{morphisms-lemma-ample-base-change} 得出（该引理说明，
对概形而言，相对丰沛性在基变换下保持）。
\end{proof}

\begin{lemma}
\label{spaces-divisors-lemma-ample-base-change}
设 $S$ 为概形，$f : X \to Y$ 为 $S$ 上代数空间的态射，
$\mathcal{L}$ 为可逆 $\mathcal{O}_X$-模。设 $Y' \to Y$ 为 $S$ 上代数空间的
态射，$f' : X' \to Y'$ 为 $f$ 的基变换，并以 $\mathcal{L}'$ 表示
$\mathcal{L}$ 到 $X'$ 上的拉回。若 $\mathcal{L}$ 是 $f$-丰沛的，则
$\mathcal{L}'$ 是 $f'$-丰沛的。
\end{lemma}

\begin{proof}
这立即由定义得出！（提示：使用基变换的传递性。）
\end{proof}

\begin{lemma}
\label{spaces-divisors-lemma-relatively-ample-properties}
设 $S$ 为概形，$f : X \to Y$ 为 $S$ 上代数空间的态射。若存在一个
$f$-丰沛可逆层，则 $f$ 可表、拟紧且分离。
\end{lemma}

\begin{proof}
这由定义及《态射》引理 \ref{morphisms-lemma-relatively-ample-separated} 显然可得。
（若对此有疑问，可参看《代数空间》引理
\ref{spaces-lemma-representable-transformations-property-implication} 的原理。）
\end{proof}

\begin{lemma}
\label{spaces-divisors-lemma-descend-relatively-ample}
设 $V \to U$ 为仿射概形之间的满平展态射，$X$ 为 $U$ 上的代数空间，
$\mathcal{L}$ 为可逆 $\mathcal{O}_X$-模。令 $Y = V \times_U X$，并令
$\mathcal{N}$ 为 $\mathcal{L}$ 到 $Y$ 上的拉回。下列条件等价：
\begin{enumerate}
\item $\mathcal{L}$ 在 $X/U$ 上丰沛；以及
\item $\mathcal{N}$ 在 $Y/V$ 上丰沛。
\end{enumerate}
\end{lemma}

\begin{proof}
蕴含关系 (1) $\Rightarrow$ (2) 由引理 \ref{spaces-divisors-lemma-ample-base-change} 得出。
假设 (2)。这蕴含 $Y \to V$ 拟紧且分离（引理
\ref{spaces-divisors-lemma-relatively-ample-properties}），并且 $Y$ 是概形。因此，态射
$f : X \to U$ 拟紧且分离（《空间的态射》引理
\ref{spaces-morphisms-lemma-quasi-compact-local} 及
\ref{spaces-morphisms-lemma-separated-local}）。置
$\mathcal{A} = \bigoplus_{d \geq 0} f_*\mathcal{L}^{\otimes d}$。
这是拟凝聚分次 $\mathcal{O}_U$-代数层（《空间的态射》引理
\ref{spaces-morphisms-lemma-pushforward}）。由伴随性得到映射
$\psi : f^*\mathcal{A} \to \bigoplus_{d \geq 0} \mathcal{L}^{\otimes d}$。
应用引理 \ref{spaces-divisors-lemma-invertible-map-into-relative-proj}，得到开子空间
$U(\psi) \subset X$ 及态射
$$
r_{\mathcal{L}, \psi} : U(\psi) \to \underline{\text{Proj}}_U(\mathcal{A})
$$
% P09-ERR-0216
由于 $h : V \to U$ 平展，故
$\mathcal{A}|_V = (Y \to V)_*(\bigoplus_{d \geq 0} \mathcal{N}^{\otimes d})$，
见《空间的性质》引理
\ref{spaces-properties-lemma-pushforward-etale-base-change-modules}。因此，
$\psi$ 到 $Y$ 上的拉回 $\psi'$ 是《态射》引理
\ref{morphisms-lemma-characterize-relatively-ample} 第 (5) 部分中情形
$(Y \to V, \mathcal{N})$ 的伴随映射。由于 $\mathcal{N}$ 在 $Y/V$ 上丰沛，
由刚引用的引理可得 $U(\psi') = Y$，且 $r_{\mathcal{N}, \psi'}$ 是开浸入。
引理 \ref{spaces-divisors-lemma-invertible-map-into-relative-proj} 告诉我们，
$r_{\mathcal{L}, \psi}$ 的构造与基变换可交换，故 $U(\psi) = X$，并且有可交换图
$$
\xymatrix{
Y \ar[r]_-{r'} \ar[d] &
\underline{\text{Proj}}_V(\mathcal{A}|_V) \ar[d] \ar[r] &
V \ar[d] \\
X \ar[r]^-r &
\underline{\text{Proj}}_U(\mathcal{A}) \ar[r] &
U
}
$$
，其中两个方块都是纤维积。由《空间的态射》引理
\ref{spaces-morphisms-lemma-closed-immersion-local} 可知 $r$ 是开浸入。因此，
$X$ 是概形。于是可应用《态射》引理
\ref{morphisms-lemma-characterize-relatively-ample} 第 (5) 部分，得出
$\mathcal{L}$ 在 $X/U$ 上丰沛。
\end{proof}

\begin{lemma}
\label{spaces-divisors-lemma-relatively-ample-local}
设 $S$ 为概形，$f : X \to Y$ 为 $S$ 上代数空间的态射，
$\mathcal{L}$ 为可逆 $\mathcal{O}_X$-模。下列条件等价：
\begin{enumerate}
\item $\mathcal{L}$ 在 $X/Y$ 上丰沛；
\item 对每个概形 $Z$ 及每个态射 $Z \to Y$，代数空间
$X_Z = Z \times_Y X$ 都是概形，且拉回 $\mathcal{L}_Z$ 在 $X_Z/Z$ 上丰沛；
\item 对每个仿射概形 $Z$ 及每个态射 $Z \to Y$，代数空间
$X_Z = Z \times_Y X$ 都是概形，且拉回 $\mathcal{L}_Z$ 在 $X_Z/Z$ 上丰沛；
\item 存在概形 $V$ 及满平展态射 $V \to Y$，使得代数空间
$X_V = V \times_Y X$ 是概形，且拉回 $\mathcal{L}_V$ 在 $X_V/V$ 上丰沛。
\end{enumerate}
\end{lemma}

\begin{proof}
按定义，(1) 与 (2) 等价。蕴含关系 (2) $\Rightarrow$ (3) 是直接的。若 (3)
成立且 $Z \to Y$ 如 (2) 所述，则可见
% P09-ERR-0217
$X_Z \to Z$ 在 $Z$ 上仿射局部地可表。因此，例如由《空间的性质》引理
\ref{spaces-properties-lemma-subscheme}，$X_Z$ 是概形。于是
$\mathcal{L}_Z$ 在 $X_Z/Z$ 上丰沛，因为这在 $Z$ 上是局部的，并且可以使用
《态射》引理 \ref{morphisms-lemma-characterize-relatively-ample}。故 (1)、(2)、
(3) 等价。显然，这些条件蕴含 (4)。

\medskip\noindent
假设 (4)。设 $Z \to Y$ 为态射，其中 $Z$ 仿射。则
$U = V \times_Y Z \to Z$ 是满平展态射，且 $\mathcal{L}_Z$ 沿
$X_U \to X_Z$ 的拉回在 $X_U/U$ 上相对丰沛。
% P09-ERR-0218
当然，我们可以用一个仿射开集替换 $U$。于是由引理
\ref{spaces-divisors-lemma-descend-relatively-ample}，$\mathcal{L}_Z$ 在 $X_Z/Z$ 上丰沛。
故 (4) $\Rightarrow$ (3)，证明完毕。
\end{proof}

\begin{lemma}
\label{spaces-divisors-lemma-quasi-affine-O-ample}
设 $S$ 为概形，$f : X \to Y$ 为 $S$ 上代数空间的态射。则 $f$ 拟仿射，
当且仅当 $\mathcal{O}_X$ 是 $f$-相对丰沛的。
\end{lemma}

\begin{proof}
这由概形情形得出，见《态射》引理
\ref{morphisms-lemma-quasi-affine-O-ample}。
\end{proof}



\section{相对丰沛性与上同调}
\label{spaces-divisors-section-ample-and-proper}

\noindent
本节包含一些与《概形的上同调》第
\ref{coherent-section-applications-formal-functions} 节及第
\ref{coherent-section-ample-cohomology} 节中的结果相关的结论。

\medskip\noindent
下面的引理只是我们所能做的一例。

\begin{lemma}
\label{spaces-divisors-lemma-vanshing-gives-ample}
% P09-ERR-0219
设 $R$ 为诺特环，$X$ 为 $R$ 上的代数空间，并且结构态射
$f : X \to \Spec(R)$ 固有。设 $\mathcal{L}$ 为可逆 $\mathcal{O}_X$-模。
下列条件等价：
\begin{enumerate}
\item $\mathcal{L}$ 在 $X/R$ 上丰沛（定义
\ref{spaces-divisors-definition-relatively-ample}）；
\item 对每个凝聚 $\mathcal{O}_X$-模 $\mathcal{F}$，存在 $n_0 \geq 0$，使得
$H^p(X, \mathcal{F} \otimes \mathcal{L}^{\otimes n}) = 0$
对所有 $n \geq n_0$ 及 $p > 0$ 成立。
\end{enumerate}
\end{lemma}

\begin{proof}
蕴含关系 (1) $\Rightarrow$ (2) 由《概形的上同调》引理
\ref{coherent-lemma-coherent-proper-ample} 得出，因为假设 (1) 蕴含 $X$ 是概形。
蕴含关系 (2) $\Rightarrow$ (1) 即《空间的上同调》引理
\ref{spaces-cohomology-lemma-Noetherian-h1-zero-invertible}。
\end{proof}

\begin{lemma}
\label{spaces-divisors-lemma-ample-on-fibre}
设 $Y$ 为诺特概形，$X$ 为 $Y$ 上的代数空间，并且结构态射
$f : X \to Y$ 固有。设 $\mathcal{L}$ 为可逆 $\mathcal{O}_X$-模，
$\mathcal{F}$ 为凝聚 $\mathcal{O}_X$-模。设 $y \in Y$ 为一点，使得 $X_y$
是概形，且 $\mathcal{L}_y$ 在 $X_y$ 上丰沛。则存在 $d_0$，使得对所有
$d \geq d_0$ 都有
$$
R^pf_*(\mathcal{F} \otimes_{\mathcal{O}_X} \mathcal{L}^{\otimes d})_y = 0
\text{ 对 }p > 0
$$
，并且映射
$$
f_*(\mathcal{F} \otimes_{\mathcal{O}_X} \mathcal{L}^{\otimes d})_y
\longrightarrow
H^0(X_y, \mathcal{F}_y \otimes_{\mathcal{O}_{X_y}} \mathcal{L}_y^{\otimes d})
$$
是满射。
\end{lemma}

\begin{proof}
注意，$\mathcal{O}_{Y, y}$ 是诺特局部环。考虑典范态射
$c : \Spec(\mathcal{O}_{Y, y}) \to Y$，见《概形》公式
(\ref{schemes-equation-canonical-morphism})。由于它诱导局部环的等同，故为平坦
态射。暂以 $f' : X' \to \Spec(\mathcal{O}_{Y, y})$ 表示 $f$ 到该局部环上的
基变换。由《空间的上同调》引理
\ref{spaces-cohomology-lemma-flat-base-change-cohomology} 可见
$c^*R^pf_*\mathcal{F} = R^pf'_*\mathcal{F}'$。此外，纤维 $X_y$ 与 $X'_y$
被等同。因此，可以假设 $Y = \Spec(A)$ 是
诺特局部环 $(A, \mathfrak m, \kappa)$ 的谱，且 $y \in Y$ 对应于
$\mathfrak m$。在此情形下，对所有 $p \geq 0$ 都有
$R^pf_*(\mathcal{F} \otimes_{\mathcal{O}_X} \mathcal{L}^{\otimes d})_y =
H^p(X, \mathcal{F} \otimes_{\mathcal{O}_X} \mathcal{L}^{\otimes d})$
。以 $f_y : X_y \to \Spec(\kappa)$ 表示投影。

\medskip\noindent
令 $B = \text{Gr}_\mathfrak m(A) =
\bigoplus_{n \geq 0} \mathfrak m^n/\mathfrak m^{n + 1}$。考虑拟凝聚分次
$\mathcal{O}_{X_y}$-代数层 $\mathcal{B} = f_y^*\widetilde{B}$。我们将沿用
《空间的上同调》第 \ref{spaces-cohomology-section-theorem-formal-functions} 节的
记号，其中以 $\mathfrak m$ 代替 $I$。由于 $X_y$ 是由
$\mathfrak m\mathcal{O}_X$ 截出的 $X$ 的闭子空间，可以把
$\mathfrak m^n\mathcal{F}/\mathfrak m^{n + 1}\mathcal{F}$ 看作凝聚
$\mathcal{O}_{X_y}$-模，见《空间的上同调》引理
\ref{spaces-cohomology-lemma-i-star-equivalence}。于是
$\bigoplus_{n \geq 0} \mathfrak m^n\mathcal{F}/\mathfrak m^{n + 1}\mathcal{F}$
是有限型拟凝聚分次 $\mathcal{B}$-模，因为它在 $\mathcal{B}$ 上由 0 次部分生成，
% P09-ERR-0220
且其 0 次部分为
$\mathcal{F}_y = \mathcal{F}/\mathfrak m \mathcal{F}$
，这是凝聚 $\mathcal{O}_{X_y}$-模。因此，由《概形的上同调》引理
\ref{coherent-lemma-graded-finiteness} 第 (2) 部分，存在 $d_0$，使得
$$
H^p(X_y, \mathfrak m^n \mathcal{F}/ \mathfrak m^{n + 1}\mathcal{F}
\otimes_{\mathcal{O}_{X_y}} \mathcal{L}_y^{\otimes d}) = 0
$$
对所有 $p > 0$、$d \geq d_0$ 及 $n \geq 0$ 成立。由《空间的上同调》引理
\ref{spaces-cohomology-lemma-relative-affine-cohomology}，这等价于
$
H^p(X, \mathfrak m^n \mathcal{F}/ \mathfrak m^{n + 1}\mathcal{F}
\otimes_{\mathcal{O}_X} \mathcal{L}^{\otimes d}) = 0
$
对所有 $p > 0$、$d \geq d_0$ 及 $n \geq 0$ 成立。

\medskip\noindent
考虑凝聚 $\mathcal{O}_X$-模的短正合列
$$
0 \to \mathfrak m^n\mathcal{F}/\mathfrak m^{n + 1} \mathcal{F}
\to \mathcal{F}/\mathfrak m^{n + 1} \mathcal{F}
\to \mathcal{F}/\mathfrak m^n \mathcal{F} \to 0
$$
。与 $\mathcal{L}^{\otimes d}$ 作张量积是正合函子，因而得到短正合列
$$
0 \to
\mathfrak m^n\mathcal{F}/\mathfrak m^{n + 1} \mathcal{F}
\otimes_{\mathcal{O}_X} \mathcal{L}^{\otimes d}
\to \mathcal{F}/\mathfrak m^{n + 1} \mathcal{F}
\otimes_{\mathcal{O}_X} \mathcal{L}^{\otimes d}
\to \mathcal{F}/\mathfrak m^n \mathcal{F}
\otimes_{\mathcal{O}_X} \mathcal{L}^{\otimes d} \to 0
$$
。利用上同调长正合列及上述消失性，以归纳法得出：
\begin{enumerate}
\item $H^p(X, \mathcal{F}/\mathfrak m^n \mathcal{F}
\otimes_{\mathcal{O}_X} \mathcal{L}^{\otimes d}) = 0$
对所有 $p > 0$、$d \geq d_0$ 及 $n \geq 0$ 成立；以及
\item $H^0(X, \mathcal{F}/\mathfrak m^n \mathcal{F}
\otimes_{\mathcal{O}_X} \mathcal{L}^{\otimes d}) \to
H^0(X_y, \mathcal{F}_y \otimes_{\mathcal{O}_{X_y}} \mathcal{L}_y^{\otimes d})$
对所有 $d \geq d_0$ 及 $n \geq 1$ 都是满射。
\end{enumerate}
由形式函数定理（《空间的上同调》定理
\ref{spaces-cohomology-theorem-formal-functions}），可知
$H^p(X, \mathcal{F} \otimes_{\mathcal{O}_X} \mathcal{L}^{\otimes d})$
的 $\mathfrak m$-进完备化对所有 $d \geq d_0$ 及 $p > 0$ 都为零。
由《空间的上同调》引理
\ref{spaces-cohomology-lemma-proper-over-affine-cohomology-finite}，
$H^p(X, \mathcal{F} \otimes_{\mathcal{O}_X} \mathcal{L}^{\otimes d})$ 是有限
$A$-模，故由 Nakayama 引理（《代数》引理 \ref{algebra-lemma-NAK}），
$H^p(X, \mathcal{F} \otimes_{\mathcal{O}_X} \mathcal{L}^{\otimes d})$
对所有 $d \geq d_0$ 及 $p > 0$ 都为零。对 $p = 0$，由《空间的上同调》引理
\ref{spaces-cohomology-lemma-ML-cohomology-powers-ideal} 第 (3) 部分可知
$H^0(X, \mathcal{F} \otimes_{\mathcal{O}_X} \mathcal{L}^{\otimes d}) \to
H^0(X_y, \mathcal{F}_y \otimes_{\mathcal{O}_{X_y}} \mathcal{L}_y^{\otimes d})$
是满射，这就给出引理的最后一个断言。
\end{proof}

\begin{lemma}
\label{spaces-divisors-lemma-ample-in-neighbourhood}
（更一般的版本见《空间上的下降》引理
\ref{spaces-descent-lemma-ample-in-neighbourhood}。）设 $Y$ 为诺特概形，
$X$ 为 $Y$ 上的代数空间，并且结构态射 $f : X \to Y$ 固有。设
$\mathcal{L}$ 为可逆 $\mathcal{O}_X$-模。设 $y \in Y$ 为一点，使得 $X_y$
是概形，且 $\mathcal{L}_y$ 在 $X_y$ 上丰沛。则存在 $y$ 的开邻域
$V \subset Y$，使得 $\mathcal{L}|_{f^{-1}(V)}$ 在 $f^{-1}(V)/V$ 上丰沛
（如定义 \ref{spaces-divisors-definition-relatively-ample}）。
\end{lemma}

\begin{proof}
对 $\mathcal{F} = \mathcal{O}_X$，选取引理 \ref{spaces-divisors-lemma-ample-on-fibre} 中的
$d_0$。选取 $d \geq d_0$，使得存在 $r \geq 0$ 及截面
$s_{y, 0}, \ldots, s_{y, r} \in H^0(X_y, \mathcal{L}_y^{\otimes d})$
，它们定义闭浸入
$$
\varphi_y =
\varphi_{\mathcal{L}_y^{\otimes d}, (s_{y, 0}, \ldots, s_{y, r})} :
X_y \to \mathbf{P}^r_{\kappa(y)}.
$$
这由《态射》引理
\ref{morphisms-lemma-finite-type-over-affine-ample-very-ample} 保证；此外，为看出
$\varphi_y$ 是闭浸入，我们还使用《态射》引理
\ref{morphisms-lemma-image-proper-scheme-closed}，并使用《构造》第
\ref{constructions-section-projective-space} 节关于以可逆层和截面描述到射影空间
之态射的内容。由 $d_0$ 的选取，在用 $y$ 的一个开邻域替换 $Y$ 后，可以选取
$s_0, \ldots, s_r \in H^0(X, \mathcal{L}^{\otimes d})$
，它们映到 $s_{y, 0}, \ldots, s_{y, r}$。令 $X_{s_i} \subset X$ 为 $s_i$
生成 $\mathcal{L}^{\otimes d}$ 所在的开子空间。由于诸 $s_{y, i}$ 生成
$\mathcal{L}_y^{\otimes d}$，可见
$|X_y| \subset U = \bigcup |X_{s_i}|$。由于 $X \to Y$ 是闭的，存在开邻域
$y \in V \subset Y$，使得 $|f|^{-1}(V) \subset U$。用 $V$ 替换 $Y$ 后，
可以假设诸 $s_i$ 生成 $\mathcal{L}^{\otimes d}$。于是得到态射
$$
\varphi = \varphi_{\mathcal{L}^{\otimes d}, (s_0, \ldots, s_r)} :
X \longrightarrow \mathbf{P}^r_Y
$$
，满足 $\mathcal{L}^{\otimes d} \cong \varphi^*\mathcal{O}_{\mathbf{P}^r_Y}(1)$，
其到 $y$ 的基变换给出
$\varphi_y$（严格来说，我们需要写出一个证明，说明《构造》第
\ref{constructions-section-projective-space} 节给出的到射影空间之态射的构造也能
用于描述从代数空间到射影空间的态射；我们略去细节）。

\medskip\noindent
我们将用一个小技巧完成证明；“正确的”证明是直接说明，在基变换到 $y$ 的一个
开邻域后，$\varphi$ 是闭浸入。具体而言，由《空间的上同调》引理
\ref{spaces-cohomology-lemma-proper-finite-fibre-finite-in-neighbourhood} 可见，
% P09-ERR-0221
$\varphi$ 在 $\mathbf{P}^r_Y \to Y$ 位于 $y$ 上方的纤维
$\mathbf{P}^r_{\kappa(y)}$ 的一个开邻域上是有限的。利用
$\mathbf{P}^r_Y \to Y$ 是闭的，在缩小 $Y$ 后，可以假设 $\varphi$ 有限。
特别地，$X$ 是概形。于是，由非常一般的《态射》引理
\ref{morphisms-lemma-pullback-ample-tensor-relatively-ample}，
$\mathcal{L}^{\otimes d} \cong \varphi^*\mathcal{O}_{\mathbf{P}^r_Y}(1)$ 是丰沛的。
\end{proof}








\section{相对 Proj 的闭子空间}
\label{spaces-divisors-section-closed-in-relative-proj}

\noindent
这里给出一些关于相对 Proj 的闭子空间的辅助引理。本节是《除子》第
\ref{divisors-section-closed-in-relative-proj} 节的对应版本。

\begin{lemma}
\label{spaces-divisors-lemma-closed-subscheme-proj}
设 $S$ 为概形，$X$ 为 $S$ 上的代数空间，$\mathcal{A}$ 为拟凝聚分次
$\mathcal{O}_X$-代数。设
$\pi : P = \underline{\text{Proj}}_X(\mathcal{A}) \to X$ 为 $\mathcal{A}$ 的
相对 Proj，$i : Z \to P$ 为闭子空间。以 $\mathcal{I} \subset \mathcal{A}$
表示典范映射
$$
\mathcal{A}
\longrightarrow
\bigoplus\nolimits_{d \geq 0} \pi_*\left((i_*\mathcal{O}_Z)(d)\right)
$$
的核。若 $\pi$ 拟紧，则有同构
$Z = \underline{\text{Proj}}_X(\mathcal{A}/\mathcal{I})$。
\end{lemma}

\begin{proof}
由引理 \ref{spaces-divisors-lemma-relative-proj-separated}，态射 $\pi$ 分离。由于 $\pi$ 拟紧，
$\pi_*$ 把拟凝聚模变为拟凝聚模，见《空间的态射》引理
\ref{spaces-morphisms-lemma-pushforward}。因此，$\mathcal{I}$ 是拟凝聚
$\mathcal{O}_X$-模。特别地，$\mathcal{B} = \mathcal{A}/\mathcal{I}$ 是拟凝聚分次
$\mathcal{O}_X$-代数。函子性态射
$Z' = \underline{\text{Proj}}_X(\mathcal{B}) \to
\underline{\text{Proj}}_X(\mathcal{A})$ 处处有定义，并且是闭浸入，见引理
\ref{spaces-divisors-lemma-surjective-graded-rings-map-relative-proj}。因此，只需证明
$Z = Z'$（作为 $P$ 的闭子空间）。

\medskip\noindent
在此之后，该问题在基上是平展局部的，故通过平展局部化归约到概形情形
（《除子》引理 \ref{divisors-lemma-closed-subscheme-proj}）。
\end{proof}

\noindent
当闭子空间局部由有限多个方程截出时，可以用 $\mathcal{A}$ 的有限型理想层定义它。

\begin{lemma}
\label{spaces-divisors-lemma-closed-subscheme-proj-finite}
设 $S$ 为概形，$X$ 为 $S$ 上拟紧且拟分离的代数空间，$\mathcal{A}$ 为拟凝聚分次
$\mathcal{O}_X$-代数。设
$\pi : P = \underline{\text{Proj}}_X(\mathcal{A}) \to X$ 为 $\mathcal{A}$ 的
相对 Proj，$i : Z \to P$ 为闭子概形。若 $\pi$ 拟紧，且
% P09-ERR-0222
$i$ 有限表示，则存在 $d > 0$ 以及有限型拟凝聚 $\mathcal{O}_X$-子模
$\mathcal{F} \subset \mathcal{A}_d$，使得
$Z = \underline{\text{Proj}}_X(\mathcal{A}/\mathcal{F}\mathcal{A})$.
\end{lemma}

\begin{proof}
读者可以重做概形情形所用的论证。不过，我们将用一个技巧说明该引理由概形情形
得出。设 $\mathcal{I} \subset \mathcal{A}$ 为引理
\ref{spaces-divisors-lemma-closed-subscheme-proj} 中截出 $Z$ 的拟凝聚分次理想。选取仿射概形
$U$ 及满平展态射 $U \to X$，见《空间的性质》引理
\ref{spaces-properties-lemma-quasi-compact-affine-cover}。由概形情形（《除子》引理
\ref{divisors-lemma-closed-subscheme-proj-finite}），存在 $d > 0$ 及有限型拟凝聚
$\mathcal{O}_U$-子模
$\mathcal{F}' \subset \mathcal{I}_d|_U \subset \mathcal{A}_d|_U$
，使得 $Z \times_X U$ 等于
$\underline{\text{Proj}}_U(\mathcal{A}|_U/\mathcal{F}'\mathcal{A}|_U)$.
由《空间的极限》引理 \ref{spaces-limits-lemma-directed-colimit-finite-type}，可找到
有限型拟凝聚子模 $\mathcal{F} \subset \mathcal{I}_d$，使得
$\mathcal{F}' \subset \mathcal{F}|_U$。令
$Z' = \underline{\text{Proj}}_X(\mathcal{A}/\mathcal{F}\mathcal{A})$.
则 $Z' \to P$ 是闭浸入
% P09-ERR-0223
（引理 \ref{spaces-divisors-lemma-surjective-generated-degree-1-map-relative-proj}），且因为
$\mathcal{F}\mathcal{A} \subset \mathcal{I}$，有 $Z \subset Z'$。另一方面，由
$\mathcal{F}$ 的选取，$Z' \times_X U \subset Z \times_X U$。故如愿有
$Z = Z'$。
\end{proof}

\begin{lemma}
\label{spaces-divisors-lemma-closed-subscheme-proj-finite-type}
设 $S$ 为概形，$X$ 为 $S$ 上拟紧且拟分离的代数空间，$\mathcal{A}$ 为拟凝聚分次
$\mathcal{O}_X$-代数。设
$\pi : P = \underline{\text{Proj}}_X(\mathcal{A}) \to X$ 为 $\mathcal{A}$ 的
相对 Proj。设
% P09-ERR-0224
$i : Z \to X$ 为闭子空间，$U \subset X$ 为开集。假设：
\begin{enumerate}
\item $\pi$ 拟紧；
% P09-ERR-0222
\item $i$ 有限表示；
\item $|U| \cap |\pi|(|i|(|Z|)) = \emptyset$,
\item $U$ 拟紧；
\item 对所有 $n$，$\mathcal{A}_n$ 都是有限型 $\mathcal{O}_X$-模。
\end{enumerate}
则存在 $d > 0$ 及有限型拟凝聚 $\mathcal{O}_X$-子模
$\mathcal{F} \subset \mathcal{A}_d$，满足 (a)
$Z = \underline{\text{Proj}}_X(\mathcal{A}/\mathcal{F}\mathcal{A})$
，且 (b) $\mathcal{A}_d/\mathcal{F}$ 的支撑与 $U$ 不交。
\end{lemma}

\begin{proof}
我们采用与引理 \ref{spaces-divisors-lemma-closed-subscheme-proj-finite} 的证明相同的技巧，归约到
概形情形。设 $\mathcal{I} \subset \mathcal{A}$ 为引理
\ref{spaces-divisors-lemma-closed-subscheme-proj} 中截出 $Z$ 的拟凝聚分次理想。选取仿射概形
$W$ 及满平展态射 $W \to X$，见《空间的性质》引理
\ref{spaces-properties-lemma-quasi-compact-affine-cover}。由概形情形（《除子》引理
\ref{divisors-lemma-closed-subscheme-proj-finite-type}），存在 $d > 0$ 及有限型
拟凝聚 $\mathcal{O}_W$-子模
$\mathcal{F}' \subset \mathcal{I}_d|_W \subset \mathcal{A}_d|_W$
，使得 (a) $Z \times_X W$ 等于
$\underline{\text{Proj}}_W(\mathcal{A}|_W/\mathcal{F}'\mathcal{A}|_W)$
，且 (b) $\mathcal{A}_d|_W/\mathcal{F}'$ 的支撑与 $U \times_X W$ 不交。
由《空间的极限》引理 \ref{spaces-limits-lemma-directed-colimit-finite-type}，可找到
有限型拟凝聚子模 $\mathcal{F} \subset \mathcal{I}_d$，使得
$\mathcal{F}' \subset \mathcal{F}|_W$。令
$Z' = \underline{\text{Proj}}_X(\mathcal{A}/\mathcal{F}\mathcal{A})$.
则 $Z' \to P$ 是闭浸入（引理
% P09-ERR-0223
\ref{spaces-divisors-lemma-surjective-generated-degree-1-map-relative-proj}），且因为
$\mathcal{F}\mathcal{A} \subset \mathcal{I}$，有 $Z \subset Z'$。另一方面，由
$\mathcal{F}$ 的选取，$Z' \times_X W \subset Z \times_X W$。故 $Z = Z'$。
最后，$\mathcal{A}_d|_W/\mathcal{F}|_W$ 是
$\mathcal{A}_d|_W/\mathcal{F}'$ 的商，而后者的支撑包含于
% P09-ERR-0225
$W \setminus U \times_X W$，故可见 $\mathcal{A}_d/\mathcal{F}$ 的支撑包含于
$X \setminus U$。引理得证。
\end{proof}

\begin{lemma}
\label{spaces-divisors-lemma-conormal-sheaf-section-projective-bundle}
设 $S$ 为概形，$X$ 为 $S$ 上的代数空间，$\mathcal{E}$ 为拟凝聚
$\mathcal{O}_X$-模。存在双射
$$
\left\{
\begin{matrix}
\text{截面 }\sigma\text{，其所属的} \\
\text{态射为 } \mathbf{P}(\mathcal{E}) \to X
\end{matrix}
\right\}
\leftrightarrow
\left\{
\begin{matrix}
\text{满射 }\mathcal{E} \to \mathcal{L}\text{，其中} \\
\mathcal{L}\text{ 是可逆 }\mathcal{O}_X\text{-模}
\end{matrix}
\right\}
$$
在此情形下，$\sigma$ 是闭浸入，并且有典范同构
$$
\Ker(\mathcal{E} \to \mathcal{L})
\otimes_{\mathcal{O}_X} \mathcal{L}^{\otimes -1}
\longrightarrow
\mathcal{C}_{\sigma(X)/\mathbf{P}(\mathcal{E})}
$$
该双射和同构都与基变换相容。
\end{lemma}

\begin{proof}
由于这些构造与基变换相容，只需在 $X$ 上平展局部地验证该断言。因此，可以假设
$X$ 是概形，此时结果即《除子》引理
\ref{divisors-lemma-conormal-sheaf-section-projective-bundle}。
\end{proof}






\section{爆破}
\label{spaces-divisors-section-blowing-up}

\noindent
爆破是代数几何中的重要工具。

\begin{definition}
\label{spaces-divisors-definition-blow-up}
设 $S$ 为概形，$X$ 为 $S$ 上的代数空间。设
$\mathcal{I} \subset \mathcal{O}_X$ 为拟凝聚理想层，并设 $Z \subset X$ 为
$\mathcal{I}$ 所对应的闭子空间（《空间的态射》引理
\ref{spaces-morphisms-lemma-closed-immersion-ideals}）。沿 $Z$ 的
{\it $X$ 的爆破}，或在理想层 $\mathcal{I}$ 处的 {\it $X$ 的爆破}，是态射
$$
b :
\underline{\text{Proj}}_X
\left(\bigoplus\nolimits_{n \geq 0} \mathcal{I}^n\right)
\longrightarrow
X
$$
。爆破的 {\it 例外除子} 是逆像 $b^{-1}(Z)$。有时称 $Z$ 为爆破的
{\it 中心}。
\end{definition}

\noindent
稍后会看到，例外除子是有效 Cartier 除子。此外，爆破可以刻画为 $X$ 上使 $Z$
的逆像成为有效 Cartier 除子的“最小”代数空间。

\medskip\noindent
若 $b : X' \to X$ 是 $X$ 在 $Z$ 处的爆破，则常以 $\mathcal{O}_{X'}(n)$ 表示
结构层的扭转。注意，它们是可逆 $\mathcal{O}_{X'}$-模，并且
$\mathcal{O}_{X'}(n) = \mathcal{O}_{X'}(1)^{\otimes n}$
，因为 $X'$ 是一个由 $1$ 次部分生成的拟凝聚分次 $\mathcal{O}_X$-代数的相对
Proj，见引理 \ref{spaces-divisors-lemma-relative-proj-generated-in-degree-1}。

\begin{lemma}
\label{spaces-divisors-lemma-blowing-up-affine}
设 $S$ 为概形，$X$ 为 $S$ 上的代数空间，
$\mathcal{I} \subset \mathcal{O}_X$ 为拟凝聚理想层。设 $U = \Spec(A)$ 为
平展于 $X$ 的仿射概形，并设 $I \subset A$ 为 $\mathcal{I}|_U$ 所对应的理想。
若 $X' \to X$ 是 $X$ 在 $\mathcal{I}$ 处的爆破，则有典范同构
$$
U \times_X X' = \text{Proj}(\bigoplus\nolimits_{d \geq 0} I^d)
$$
，这是 $U$ 上概形的同构；右端是 $A$ 中 $I$ 的 Rees 代数的齐次谱。此外，
$U \times_X X'$ 有一个由仿射爆破代数 $A[\frac{I}{a}]$ 的谱组成的仿射开覆盖。
\end{lemma}

\begin{proof}
注意，限制 $\mathcal{I}|_U$ 等于 $\mathcal{I}$ 沿态射 $U \to X$ 的拉回，见
《空间的性质》第 \ref{spaces-properties-section-modules} 节。因此，将引理
\ref{spaces-divisors-lemma-relative-proj} 与《除子》引理
\ref{divisors-lemma-blowing-up-affine} 结合即得本引理。
\end{proof}

\begin{lemma}
\label{spaces-divisors-lemma-flat-base-change-blowing-up}
设 $S$ 为概形，$X_1 \to X_2$ 为 $S$ 上代数空间的平坦态射，
$Z_2 \subset X_2$ 为闭子空间。设 $Z_1$ 为 $Z_2$ 在 $X_1$ 中的逆像，并设
$X'_i$ 为 $X_i$ 在 $Z_i$ 处的爆破。则存在笛卡尔图
$$
\xymatrix{
X_1' \ar[r] \ar[d] & X_2' \ar[d] \\
X_1 \ar[r] & X_2
}
$$
，其中各项都是 $S$ 上的代数空间。
\end{lemma}

\begin{proof}
设 $\mathcal{I}_2$ 为 $Z_2$ 在 $X_2$ 中的理想层，以 $g : X_1 \to X_2$ 表示
给定态射。则 $Z_1$ 的理想层 $\mathcal{I}_1$ 是映射
$g^*\mathcal{I}_2 \to \mathcal{O}_{X_1}$
的像（见《空间的态射》定义
\ref{spaces-morphisms-definition-inverse-image-closed-subspace} 及定义后的讨论）。
由引理 \ref{spaces-divisors-lemma-relative-proj-base-change} 可见，
$X_1 \times_{X_2} X_2'$ 是 $\bigoplus_{n \geq 0} g^*\mathcal{I}_2^n$ 的
相对 Proj。由于 $g$ 平坦，映射
$g^*\mathcal{I}_2^n \to \mathcal{O}_{X_1}$ 是单射，其像为
$\mathcal{I}_1^n$。故 $X_1 \times_{X_2} X_2' = X_1'$。
\end{proof}

\begin{lemma}
\label{spaces-divisors-lemma-blowing-up-gives-effective-Cartier-divisor}
设 $S$ 为概形，$X$ 为 $S$ 上的代数空间，$Z \subset X$ 为闭子空间。
$X$ 在 $Z$ 处的爆破 $b : X' \to X$ 具有下列性质：
\begin{enumerate}
\item $b|_{b^{-1}(X \setminus Z)} : b^{-1}(X \setminus Z) \to X \setminus Z$
是同构；
\item 例外除子 $E = b^{-1}(Z)$ 是 $X'$ 上的有效 Cartier 除子；
\item 有典范同构
$\mathcal{O}_{X'}(-1) = \mathcal{O}_{X'}(E)$
\end{enumerate}
\end{lemma}

\begin{proof}
设 $U$ 为概形，$U \to X$ 为满平展态射。由于爆破与平坦基变换可交换（引理
\ref{spaces-divisors-lemma-flat-base-change-blowing-up}），可以在基变换到 $U$ 后证明各断言。
这把问题归约到概形情形，此时结果即《除子》引理
\ref{divisors-lemma-blowing-up-gives-effective-Cartier-divisor}。
\end{proof}

\begin{lemma}[爆破的泛性质]
\label{spaces-divisors-lemma-universal-property-blowing-up}
% P09-ERR-0226
\begin{slogan}
爆破一个闭子集，使之成为 Cartier 的。
\end{slogan}
设 $S$ 为概形，$X$ 为 $S$ 上的代数空间，$Z \subset X$ 为闭子空间。设
$\mathcal{C}$ 为 $(\textit{Spaces}/X)$ 的满子范畴，其对象是满足如下条件的
$Y \to X$：$Z$ 的逆像是 $Y$ 上的有效 Cartier 除子。则 $X$ 在 $Z$ 处的爆破
$b : X' \to X$ 是 $\mathcal{C}$ 的终对象。
\end{lemma}

\begin{proof}
由引理 \ref{spaces-divisors-lemma-blowing-up-gives-effective-Cartier-divisor} 可见，
$b : X' \to X$ 是 $\mathcal{C}$ 的对象。设 $f : Y \to X$ 为
$\mathcal{C}$ 的对象。我们须证明存在唯一的 $X$ 上态射 $Y \to X'$。令
$D = f^{-1}(Z)$。设 $\mathcal{I} \subset \mathcal{O}_X$ 为 $Z$ 的理想层，
$\mathcal{I}_D$ 为 $D$ 的理想层。则
$f^*\mathcal{I} \to \mathcal{I}_D$ 是到可逆 $\mathcal{O}_Y$-模的满射。
它延拓为分次 $\mathcal{O}_Y$-代数映射
$\psi : \bigoplus f^*\mathcal{I}^d \to \bigoplus \mathcal{I}_D^d$
。（注意，因为 $D$ 是有效 Cartier 除子，有
$\mathcal{I}_D^d = \mathcal{I}_D^{\otimes d}$。）由引理
% P09-ERR-0227
\ref{spaces-divisors-lemma-relative-proj-generated-in-degree-1}，三元组
$(f : Y \to X, \mathcal{I}_D, \psi)$ 定义了 $X$ 上态射 $Y \to X'$。其限制
$$
Y \setminus D \longrightarrow X' \setminus b^{-1}(Z) = X \setminus Z
$$
是唯一的。由引理 \ref{spaces-divisors-lemma-complement-effective-Cartier-divisor}，开子空间
$Y \setminus D$ 在 $Y$ 中概形论稠密。因此，由《空间的态射》引理
\ref{spaces-morphisms-lemma-equality-of-morphisms}，态射 $Y \to X'$ 唯一
（而且由引理 \ref{spaces-divisors-lemma-relative-proj-separated}，$b$ 分离）。
\end{proof}

\begin{lemma}
\label{spaces-divisors-lemma-blow-up-effective-Cartier-divisor}
设 $S$ 为概形，$X$ 为 $S$ 上的代数空间，$Z \subset X$ 为有效 Cartier
除子。$X$ 在 $Z$ 处的爆破是 $X$ 的恒等态射。
\end{lemma}

\begin{proof}
这立即由爆破的泛性质（引理 \ref{spaces-divisors-lemma-universal-property-blowing-up}）得出。
\end{proof}

\begin{lemma}
\label{spaces-divisors-lemma-blow-up-reduced-space}
设 $S$ 为概形，$X$ 为 $S$ 上的代数空间，
$\mathcal{I} \subset \mathcal{O}_X$ 为拟凝聚理想层。若 $X$ 约化，则 $X$
在 $\mathcal{I}$ 处的爆破 $X'$ 约化。
\end{lemma}

\begin{proof}
设 $U$ 为概形，$U \to X$ 为满平展态射。由于爆破与平坦基变换可交换（引理
\ref{spaces-divisors-lemma-flat-base-change-blowing-up}），可以在基变换到 $U$ 后证明该断言。
这把问题归约到概形情形，此时结果即《除子》引理
\ref{divisors-lemma-blow-up-reduced-scheme}。
\end{proof}

\begin{lemma}
\label{spaces-divisors-lemma-blowup-finite-nr-irreducibles}
设 $S$ 为概形，$X$ 为 $S$ 上的代数空间，$b : X' \to X$ 为 $X$ 在某闭子空间
处的爆破。若 $X$ 满足《空间的态射》引理
\ref{spaces-morphisms-lemma-prepare-normalization} 的等价条件，则 $X'$ 也满足。
\end{lemma}

\begin{proof}
这立即由陈述中所引用的引理、引理 \ref{spaces-divisors-lemma-blowing-up-affine} 中爆破的平展局部
描述以及《除子》引理
\ref{divisors-lemma-blow-up-and-irreducible-components} 得出。
\end{proof}

\begin{lemma}
\label{spaces-divisors-lemma-blow-up-pullback-effective-Cartier}
设 $S$ 为概形，$X$ 为 $S$ 上的代数空间，$b : X' \to X$ 为 $X$ 在某闭子空间
处的爆破。对 $X$ 上任意有效 Cartier 除子 $D$，拉回 $b^{-1}D$ 有定义
（见定义 \ref{spaces-divisors-definition-pullback-effective-Cartier-divisor}）。
\end{lemma}

\begin{proof}
由引理 \ref{spaces-divisors-lemma-blowing-up-affine} 及
\ref{spaces-divisors-lemma-characterize-effective-Cartier-divisor}，这归约为下列代数事实：
设 $A$ 为环，$I \subset A$ 为理想，$a \in I$，而 $x \in A$ 为非零因子。
则 $x$ 在 $A[\frac{I}{a}]$ 中的像是非零因子。事实上，假设
$x (y/a^n) = 0$ 在 $A[\frac{I}{a}]$ 中成立。则对某个 $m$，在 $A$ 中有
$a^mxy = 0$。由于 $x$ 是非零因子，故 $a^my = 0$。于是如所需，$y/a^n$
在 $A[\frac{I}{a}]$ 中为零。
\end{proof}

\begin{lemma}
\label{spaces-divisors-lemma-blowing-up-two-ideals}
设 $S$ 为概形，$X$ 为 $S$ 上的代数空间。设
% P09-ERR-0228
$\mathcal{I} \subset \mathcal{O}_X$ 及 $\mathcal{J}$ 为拟凝聚理想层。设
$b : X' \to X$ 为 $X$ 在 $\mathcal{I}$ 处的爆破，$b' : X'' \to X'$ 为
$X'$ 在 $b^{-1}\mathcal{J} \mathcal{O}_{X'}$ 处的爆破。则 $X'' \to X$ 典范
同构于 $X$ 在 $\mathcal{I}\mathcal{J}$ 处的爆破。
\end{lemma}

\begin{proof}
设 $E \subset X'$ 为 $b$ 的例外除子，由引理
\ref{spaces-divisors-lemma-blowing-up-gives-effective-Cartier-divisor}，它是有效 Cartier 除子。
由引理 \ref{spaces-divisors-lemma-blow-up-pullback-effective-Cartier}，$(b')^{-1}E$ 是
$X''$ 上的有效 Cartier 除子。设 $E' \subset X''$ 为 $b'$ 的例外除子（它也是
有效 Cartier 除子）。考虑有效 Cartier 除子
$E'' = E' + (b')^{-1}E$。按构造，$E''$ 的理想为
$(b \circ b')^{-1}\mathcal{I} (b \circ b')^{-1}\mathcal{J} \mathcal{O}_{X''}$.
因此，由引理 \ref{spaces-divisors-lemma-universal-property-blowing-up}，存在从 $X''$ 到
$X$ 在 $\mathcal{I}\mathcal{J}$ 处的爆破 $c : Y \to X$ 的典范态射。反过来，
由于 $\mathcal{I}\mathcal{J}$ 拉回为可逆理想，可见
$c^{-1}\mathcal{I}\mathcal{O}_Y$ 定义有效 Cartier 除子，见引理
\ref{spaces-divisors-lemma-sum-closed-subschemes-effective-Cartier}。因此，由引理
% P09-ERR-0229
\ref{spaces-divisors-lemma-universal-property-blowing-up}，存在 $X$ 上态射 $c' : Y \to X'$。
于是 $(c')^{-1}b^{-1}\mathcal{J}\mathcal{O}_Y = c^{-1}\mathcal{J}\mathcal{O}_Y$，
它也定义有效 Cartier 除子。因此，存在
$X'$ 上态射 $c'' : Y \to X''$。我们略去验证此态射与先前构造的态射
$X'' \to Y$ 互逆。
\end{proof}

\begin{lemma}
\label{spaces-divisors-lemma-blowing-up-projective}
设 $S$ 为概形，$X$ 为 $S$ 上的代数空间，
$\mathcal{I} \subset \mathcal{O}_X$ 为拟凝聚理想层。设 $b : X' \to X$
为 $X$ 在理想层 $\mathcal{I}$ 处的爆破。若 $\mathcal{I}$ 是有限型的，则
$b : X' \to X$ 是固有态射。
\end{lemma}

\begin{proof}
设 $U$ 为概形，$U \to X$ 为满平展态射。由于爆破与平坦基变换可交换（引理
\ref{spaces-divisors-lemma-flat-base-change-blowing-up}），可以在基变换到 $U$ 后证明该断言
（见《空间的态射》引理 \ref{spaces-morphisms-lemma-proper-local}）。这把问题归约到
概形情形。此时，由《除子》引理 \ref{divisors-lemma-blowing-up-projective}，
态射 $b$ 是射影的，故由《态射》引理
\ref{morphisms-lemma-locally-projective-proper}，它是固有的。
\end{proof}

\begin{lemma}
\label{spaces-divisors-lemma-composition-finite-type-blowups}
设 $S$ 为概形，$X$ 为 $S$ 上的代数空间。假设 $X$ 拟紧且拟分离。设
$Z \subset X$ 为有限表示闭子空间，$b : X' \to X$ 为中心为 $Z$ 的爆破。
设 $Z' \subset X'$ 为有限表示闭子空间，$X'' \to X'$ 为中心为 $Z'$ 的爆破。
则存在有限表示闭子空间 $Y \subset X$，使得：
\begin{enumerate}
\item $|Y| = |Z| \cup |b|(|Z'|)$；
\item 复合态射 $X'' \to X$ 同构于 $X$ 在 $Y$ 处的爆破。
\end{enumerate}
\end{lemma}

\begin{proof}
条件 $Z \to X$ 有限表示意味着 $Z$ 由有限型拟凝聚理想层
$\mathcal{I} \subset \mathcal{O}_X$ 截出，见《空间的态射》引理
\ref{spaces-morphisms-lemma-closed-immersion-finite-presentation}。写
$\mathcal{A} = \bigoplus_{n \geq 0} \mathcal{I}^n$，从而
$X' = \underline{\text{Proj}}(\mathcal{A})$。注意，由《空间的极限》引理
\ref{spaces-limits-lemma-quasi-coherent-finite-type-ideals}，$X \setminus Z$ 是
$X$ 的拟紧开子空间。由于 $b^{-1}(X \setminus Z) \to X \setminus Z$ 是同构
（引理 \ref{spaces-divisors-lemma-blowing-up-gives-effective-Cartier-divisor}），同一结果表明
$b^{-1}(X \setminus Z) \setminus Z'$ 是 $X'$ 的拟紧开子空间。因此，
$U = X \setminus (Z \cup b(Z'))$ 是 $X$ 的拟紧开子空间。由引理
\ref{spaces-divisors-lemma-closed-subscheme-proj-finite-type}，存在 $d > 0$ 及有限型
$\mathcal{O}_X$-子模 $\mathcal{F} \subset \mathcal{I}^d$，使得
$Z' = \underline{\text{Proj}}(\mathcal{A}/\mathcal{F}\mathcal{A})$，并且
$\mathcal{I}^d/\mathcal{F}$ 的支撑包含于 $X \setminus U$。

\medskip\noindent
由于 $\mathcal{F} \subset \mathcal{I}^d$ 是 $\mathcal{O}_X$-子模，可以把
$\mathcal{F} \subset \mathcal{I}^d \subset \mathcal{O}_X$ 看作 $X$ 上的有限型
拟凝聚理想层。为避免混淆，将其记为
$\mathcal{J} \subset \mathcal{O}_X$。由于 $\mathcal{I}^d / \mathcal{J}$ 及
$\mathcal{O}/\mathcal{I}^d$ 的支撑都包含于 $|X| \setminus |U|$，可见
$|V(\mathcal{J})|$ 包含于 $|X| \setminus |U|$。反过来，由于
$\mathcal{J} \subset \mathcal{I}^d$，可见
$|Z| \subset |V(\mathcal{J})|$。在
$X \setminus Z \cong X' \setminus b^{-1}(Z)$ 上，理想层 $\mathcal{J}$ 截出
$Z'$（见下方公式）。因此，$|V(\mathcal{J})|$ 等于
$|Z| \cup |b|(|Z'|)$，从而还有
$|V(\mathcal{I}\mathcal{J})| = |Z| \cup |b|(|Z'|)$。此外，
$\mathcal{I}\mathcal{J}$ 是两个有限型理想之积，故也是有限型理想。我们断言
$X'' \to X$ 同构于 $X$ 在 $\mathcal{I}\mathcal{J}$ 处的爆破；令
$Y = V(\mathcal{I}\mathcal{J})$，这便完成引理的证明。

\medskip\noindent
首先回顾，$X$ 在 $\mathcal{I}\mathcal{J}$ 处的爆破等同于 $X'$ 在
$b^{-1}\mathcal{J} \mathcal{O}_{X'}$ 处的爆破，见引理
\ref{spaces-divisors-lemma-blowing-up-two-ideals}。因此，只需证明 $X'$ 在
$b^{-1}\mathcal{J} \mathcal{O}_{X'}$ 处的爆破与 $X'$ 在 $Z'$ 处的爆破一致。
我们将证明
$$
b^{-1}\mathcal{J} \mathcal{O}_{X'} = \mathcal{I}_E^d \mathcal{I}_{Z'}
$$
% P09-ERR-0230
，这是 $X''$ 上理想层的等式。这将证明所需，因为 $\mathcal{I}_E^d$ 截出有效
Cartier 除子 $dE$，而我们可以使用引理
\ref{spaces-divisors-lemma-blow-up-effective-Cartier-divisor} 和
\ref{spaces-divisors-lemma-blowing-up-two-ideals}。

\medskip\noindent
为验证所示理想的等式，可以局部进行。采用引理
\ref{spaces-divisors-lemma-blowing-up-affine} 中的记号 $A$、$I$、$a \in I$，可见
% P09-ERR-0231
$\mathcal{F}$ 对应于一个 $R$-子模 $M \subset I^d$，后者同构映到理想
$J \subset R$。条件
$Z' = \underline{\text{Proj}}(\mathcal{A}/\mathcal{F}\mathcal{A})$
意味着 $Z' \cap \Spec(A[\frac{I}{a}])$ 由元素 $m/a^d$（$m \in M$）生成的
理想截出。设元素 $m \in M$ 对应于函数 $f \in J$。那么在仿射爆破代数
$A' = A[\frac{I}{a}]$ 中，有 $f = (a^dm)/a^d = a^d (m/a^d)$。
故等式成立。
\end{proof}









\section{严格变换}
\label{spaces-divisors-section-strict-transform}

\noindent
本节是《除子》第 \ref{divisors-section-strict-transform} 节的对应版本。设 $S$
为概形，$B$ 为 $S$ 上的代数空间，且 $Z \subset B$ 为闭子空间。设
$b : B' \to B$ 为 $B$ 在 $Z$ 处的爆破，并以 $E \subset B'$ 表示例外除子
% P09-ERR-0232
$E = b^{-1}Z$。下文常考虑 $B$ 上的代数空间 $X$，并作笛卡尔图
$$
\xymatrix{
\text{pr}_{B'}^{-1}E \ar[r] \ar[d] &
X \times_B B' \ar[r]_-{\text{pr}_X} \ar[d]_{\text{pr}_{B'}} &
X \ar[d]^f \\
E \ar[r] & B' \ar[r] & B
}
$$
。由于 $E$ 是有效 Cartier 除子（引理
\ref{spaces-divisors-lemma-blowing-up-gives-effective-Cartier-divisor}），可见
$\text{pr}_{B'}^{-1}E \subset X \times_B B'$ 是局部主的（引理
\ref{spaces-divisors-lemma-pullback-locally-principal}）。因此，
$\text{pr}_{B'}^{-1}E$ 在 $X \times_B B'$ 中之补集的包含态射是仿射的，特别地
是拟紧的（引理 \ref{spaces-divisors-lemma-complement-locally-principal-closed-subscheme}）。
因此，对拟凝聚 $\mathcal{O}_{X \times_B B'}$-模 $\mathcal{G}$，支撑于
$|\text{pr}_{B'}^{-1}E|$ 的截面子层是拟凝聚子模，见《空间的极限》定义
\ref{spaces-limits-definition-subsheaf-sections-supported-on-closed}。若
$\mathcal{G}$ 是拟凝聚代数层，例如
$\mathcal{G} = \mathcal{O}_{X \times_B B'}$，则此子层是 $\mathcal{G}$ 的理想。

\begin{definition}
\label{spaces-divisors-definition-strict-transform}
% P09-ERR-0233
沿用上述 $Z \subset B$ 与 $f : X \to B$。
\begin{enumerate}
\item 给定拟凝聚 $\mathcal{O}_X$-模 $\mathcal{F}$，$\mathcal{F}$ 关于 $B$
在 $Z$ 处爆破的{\it 严格变换}，是 $\text{pr}_X^*\mathcal{F}$ 模去支撑于
$|\text{pr}_{B'}^{-1}E|$ 的截面子模所得的商 $\mathcal{F}'$。
\item $X$ 的{\it 严格变换}是闭子空间
$X' \subset X \times_B B'$，它由 $\mathcal{O}_{X \times_B B'}$ 中支撑于
$|\text{pr}_{B'}^{-1}E|$ 的截面所成的拟凝聚理想截出。
\end{enumerate}
\end{definition}

\noindent
注意，沿一次爆破取严格变换依赖于该爆破所使用的闭子空间（而不仅依赖于态射
$B' \to B$）。

\begin{lemma}[平展局部化与严格变换]
\label{spaces-divisors-lemma-strict-transform-local}
处于定义 \ref{spaces-divisors-definition-strict-transform} 的情形。设
$$
\xymatrix{
U \ar[r] \ar[d] & X \ar[d] \\
V \ar[r] & B
}
$$
为态射的交换图，其中 $U$ 与 $V$ 为概形，水平箭头为平展态射。设
$V' \to V$ 为 $V$ 在 $Z \times_B V$ 处的爆破。则
\begin{enumerate}
\item $V' = V \times_B B'$，且态射
$V' \to B'$ 与 $U \times_V V' \to X \times_B B'$ 都是平展的；
\item $U$ 关于 $V' \to V$ 的严格变换 $U'$ 等于 $X' \times_X U$，其中
$X'$ 是 $X$ 关于 $B' \to B$ 的严格变换；并且
\item 对拟凝聚 $\mathcal{O}_X$-模 $\mathcal{F}$，严格变换
$\mathcal{F}'$ 在 $U \times_V V'$ 上的限制是 $\mathcal{F}|_U$ 关于
$V' \to V$ 的严格变换。
\end{enumerate}
\end{lemma}

\begin{proof}
(1) 由爆破与平坦基变换可交换（引理
\ref{spaces-divisors-lemma-flat-base-change-blowing-up}）、平展态射平坦以及平展态射的基变换仍平展
这些事实得出。(3) 随即由下述事实得出：取支撑于闭集的截面层与沿平展态射的拉回
可交换，见《空间的极限》引理
% P09-ERR-0234
\ref{spaces-limits-lemma-sections-supported-on-closed-subset}。(2) 由把 (3) 应用于
$\mathcal{F} = \mathcal{O}_X$ 得出。
\end{proof}

\begin{lemma}
\label{spaces-divisors-lemma-strict-transform}
处于定义 \ref{spaces-divisors-definition-strict-transform} 的情形。
\begin{enumerate}
\item $X$ 的严格变换 $X'$ 是 $X$ 在闭子空间 $f^{-1}Z$ 处的爆破；这里该闭子空间
位于 $X$ 中。
\item 对拟凝聚 $\mathcal{O}_X$-模 $\mathcal{F}$，严格变换
$\mathcal{F}'$ 典范同构于 $\mathcal{F}$ 关于爆破 $X' \to X$ 的严格变换沿
$X' \to X \times_B B'$ 的正像。
\end{enumerate}
\end{lemma}

\begin{proof}
设 $X'' \to X$ 为 $X$ 在 $f^{-1}Z$ 处的爆破。由爆破的泛性质（引理
\ref{spaces-divisors-lemma-universal-property-blowing-up}），存在交换图
$$
\xymatrix{
X'' \ar[r] \ar[d] & X \ar[d] \\
B' \ar[r] & B
}
$$
，从而得到态射 $i : X'' \to X \times_B B'$。引理的第一个断言是 $i$ 为像等于
$X'$ 的闭浸入。第二个断言是 $\mathcal{F}' = i_*\mathcal{F}''$，其中
$\mathcal{F}''$ 是 $\mathcal{F}$ 关于爆破 $X'' \to X$ 的严格变换。可以在 $X$
上平展局部地检验这些断言，因而归约到概形情形（《除子》引理
\ref{divisors-lemma-strict-transform}）。略去若干细节。
\end{proof}

\begin{lemma}
\label{spaces-divisors-lemma-strict-transform-flat}
处于定义 \ref{spaces-divisors-definition-strict-transform} 的情形。
\begin{enumerate}
\item 若 $X$ 在所有位于 $Z$ 上方的点处都平坦于 $B$，则 $X$ 的严格变换等于
基变换 $X \times_B B'$。
\item 设 $\mathcal{F}$ 为拟凝聚 $\mathcal{O}_X$-模。若 $\mathcal{F}$ 在所有
位于 $Z$ 上方的点处都平坦于 $B$，则 $\mathcal{F}$ 的严格变换
$\mathcal{F}'$ 等于拉回 $\text{pr}_X^*\mathcal{F}$。
\end{enumerate}
\end{lemma}

\begin{proof}
略。提示：由概形情形（《除子》引理
\ref{divisors-lemma-strict-transform-flat}）经平展局部化（引理
\ref{spaces-divisors-lemma-strict-transform-local}）得出。
\end{proof}

\begin{lemma}
\label{spaces-divisors-lemma-strict-transform-affine}
设 $S$ 为概形，$B$ 为 $S$ 上的代数空间，$Z \subset B$ 为闭子空间。设
$b : B' \to B$ 为 $B$ 在 $Z$ 处的爆破，$g : X \to Y$ 为 $B$ 上代数空间的
仿射态射。设 $\mathcal{F}$ 为 $X$ 上的拟凝聚层。设
$g' : X \times_B B' \to Y \times_B B'$ 为 $g$ 的基变换，$\mathcal{F}'$ 为
$\mathcal{F}$ 关于 $b$ 的严格变换。则 $g'_*\mathcal{F}'$ 是
$g_*\mathcal{F}$ 的严格变换。
\end{lemma}

\begin{proof}
略。提示：由概形情形（《除子》引理
\ref{divisors-lemma-strict-transform-affine}）经平展局部化（引理
\ref{spaces-divisors-lemma-strict-transform-local}）得出。
\end{proof}

\begin{lemma}
\label{spaces-divisors-lemma-strict-transform-different-centers}
设 $S$ 为概形，$B$ 为 $S$ 上的代数空间，$Z \subset B$ 为闭子空间。设
$D \subset B$ 为有效 Cartier 除子，$Z' \subset B$ 为由 $Z$ 与 $D$ 的理想层
之积截出的闭子空间。设 $B' \to B$ 为 $B$ 在 $Z$ 处的爆破。
\begin{enumerate}
\item $B$ 在 $Z'$ 处的爆破同构于 $B' \to B$。
\item 设 $f : X \to B$ 为代数空间的态射，$\mathcal{F}$ 为拟凝聚
$\mathcal{O}_X$-模。若 $\mathcal{F}$ 中支撑于 $|f^{-1}D|$ 的截面子层为零，
则 $\mathcal{F}$ 关于在 $Z$ 处爆破的严格变换，与 $\mathcal{F}$ 关于 $B$
在 $Z'$ 处爆破的严格变换一致。
\end{enumerate}
\end{lemma}

\begin{proof}
略。提示：由概形情形（《除子》引理
\ref{divisors-lemma-strict-transform-different-centers}）经平展局部化（引理
\ref{spaces-divisors-lemma-strict-transform-local}）得出。
\end{proof}

\begin{lemma}
\label{spaces-divisors-lemma-strict-transform-composition-blowups}
设 $S$ 为概形，$B$ 为 $S$ 上的代数空间，$Z \subset B$ 为闭子空间。设
$b : B' \to B$ 为中心为 $Z$ 的爆破。设 $Z' \subset B'$ 为闭子空间，
$B'' \to B'$ 为中心为 $Z'$ 的爆破。设
% P09-ERR-0235
$Y \subset B$ 为闭子概形，并且 $|Y| = |Z| \cup |b|(|Z'|)$，复合态射
$B'' \to B$ 同构于 $B$ 在 $Y$ 处的爆破。在此情形，对任意 $B$ 上概形 $X$
及 $\mathcal{F} \in \QCoh(\mathcal{O}_X)$，有
\begin{enumerate}
\item $\mathcal{F}$ 关于 $B$ 在 $Y$ 处爆破的严格变换，等于先取
$\mathcal{F}$ 关于 $B$ 在 $Z$ 处的爆破 $B' \to B$ 的严格变换，再取所得层关于
在 $Z'$ 处的爆破 $B'' \to B'$ 的严格变换；并且
\item $X$ 关于 $B$ 在 $Y$ 处爆破的严格变换，等于先取 $X$ 关于 $B$ 在 $Z$
处的爆破 $B' \to B$ 的严格变换，再取所得空间关于在 $Z'$ 处的爆破
$B'' \to B'$ 的严格变换。
\end{enumerate}
\end{lemma}

\begin{proof}
略。提示：由概形情形（《除子》引理
\ref{divisors-lemma-strict-transform-composition-blowups}）经平展局部化（引理
\ref{spaces-divisors-lemma-strict-transform-local}）得出。
\end{proof}

\begin{lemma}
\label{spaces-divisors-lemma-strict-transform-universally-injective}
处于定义 \ref{spaces-divisors-definition-strict-transform} 的情形。假设
$$
0 \to \mathcal{F}_1 \to \mathcal{F}_2 \to \mathcal{F}_3 \to 0
$$
是 $X$ 上拟凝聚层的正合列，并且在任意基变换 $T \to B$ 后仍正合。则关于任意
爆破 $B' \to B$ 的严格变换
% P09-ERR-0236
$\mathcal{F}_i'$ 也组成短正合列
$0 \to \mathcal{F}'_1 \to \mathcal{F}'_2 \to \mathcal{F}'_3 \to 0$。
\end{lemma}

\begin{proof}
略。提示：由概形情形（《除子》引理
\ref{divisors-lemma-strict-transform-universally-injective}）经平展局部化（引理
\ref{spaces-divisors-lemma-strict-transform-local}）得出。
\end{proof}

\begin{lemma}
\label{spaces-divisors-lemma-strict-transform-blowup-fitting-ideal}
设 $S$ 为概形，$B$ 为 $S$ 上的代数空间，$\mathcal{F}$ 为有限型拟凝聚
$\mathcal{O}_B$-模。设
% P09-ERR-0237
$Z_k \subset S$ 为由 $\text{Fit}_k(\mathcal{F})$ 截出的闭子概形，见第
\ref{spaces-divisors-section-fitting-ideals} 节。设 $B' \to B$ 为 $B$ 在 $Z_k$ 处的爆破，
$\mathcal{F}'$ 为 $\mathcal{F}$ 的严格变换。则 $\mathcal{F}'$ 局部可由
$\leq k$ 个截面生成。
\end{lemma}

\begin{proof}
略。由概形情形（《除子》引理
\ref{divisors-lemma-strict-transform-blowup-fitting-ideal}）经平展局部化（引理
\ref{spaces-divisors-lemma-strict-transform-local}）得出。
\end{proof}

\begin{lemma}
\label{spaces-divisors-lemma-strict-transform-blowup-fitting-ideal-locally-free}
设 $S$ 为概形，$B$ 为 $S$ 上的代数空间，$\mathcal{F}$ 为有限型拟凝聚
$\mathcal{O}_B$-模。设
% P09-ERR-0237
$Z_k \subset S$ 为由 $\text{Fit}_k(\mathcal{F})$ 截出的闭子概形，见第
\ref{spaces-divisors-section-fitting-ideals} 节。假设 $\mathcal{F}$ 在 $B \setminus Z_k$ 上
局部自由且秩为 $k$。设 $B' \to B$ 为 $B$ 在 $Z_k$ 处的爆破，$\mathcal{F}'$
为 $\mathcal{F}$ 的严格变换。则 $\mathcal{F}'$ 局部自由且秩为 $k$。
\end{lemma}

\begin{proof}
略。由概形情形（《除子》引理
\ref{divisors-lemma-strict-transform-blowup-fitting-ideal-locally-free}）
经平展局部化（引理 \ref{spaces-divisors-lemma-strict-transform-local}）得出。
\end{proof}












\section{容许爆破}
\label{spaces-divisors-section-admissible-blowups}

\noindent
为更好地控制所用的爆破，引入以下标准术语。

\begin{definition}
\label{spaces-divisors-definition-admissible-blowup}
设 $S$ 为概形，$X$ 为 $S$ 上的代数空间，$U \subset X$ 为开子空间。若存在
有限表示闭浸入 $Z \to X$，其中 $Z$ 与 $U$ 不交，并且 $X'$ 同构于 $X$ 在 $Z$ 处的爆破，
则称态射 $X' \to X$ 为{\it $U$-容许爆破}。
\end{definition}

\noindent
回顾：$Z \to X$ 有限表示，当且仅当理想层
$\mathcal{I}_Z \subset \mathcal{O}_X$ 为有限型，见《空间的态射》引理
\ref{spaces-morphisms-lemma-closed-immersion-finite-presentation}。特别地，
$U$-容许爆破是固有态射，见引理 \ref{spaces-divisors-lemma-blowing-up-projective}。注意，可能有多个
中心给出同一态射。因此，要求仅仅是存在某个与 $U$ 不交且产生 $X'$ 的中心。最后，
由于态射 $b : X' \to X$ 在 $U$ 上是同构（见引理
\ref{spaces-divisors-lemma-blowing-up-gives-effective-Cartier-divisor}），我们常滥用记号，也把 $U$
看作 $X'$ 的开子空间。

\begin{lemma}
\label{spaces-divisors-lemma-composition-admissible-blowups}
设 $S$ 为概形，$X$ 为 $S$ 上拟紧拟分离的代数空间，$U \subset X$ 为拟紧开子空间。
设 $b : X' \to X$ 为 $U$-容许爆破，$X'' \to X'$ 为 $U$-容许爆破。则复合态射
$X'' \to X$ 是 $U$-容许爆破。
\end{lemma}

\begin{proof}
立即由更精确的引理 \ref{spaces-divisors-lemma-composition-finite-type-blowups} 得出。
\end{proof}

\begin{lemma}
\label{spaces-divisors-lemma-extend-admissible-blowups}
设 $S$ 为概形。设
% P09-ERR-0238
$X$ 为拟紧拟分离代数空间，$U, V \subset X$ 为拟紧开子空间。设
$b : V' \to V$ 为 $U \cap V$-容许爆破。则存在 $U$-容许爆破 $X' \to X$，
其在 $V$ 上的限制为 $V'$。
\end{lemma}

\begin{proof}
设 $\mathcal{I} \subset \mathcal{O}_V$ 为有限型拟凝聚理想层，使得
$V(\mathcal{I})$ 与 $U \cap V$ 不交，且 $V'$ 同构于 $V$ 在 $\mathcal{I}$ 处的
爆破。设 $\mathcal{I}' \subset \mathcal{O}_{U \cup V}$ 为拟凝聚理想层，其在 $U$
上的限制为 $\mathcal{O}_U$，在 $V$ 上的限制为 $\mathcal{I}$。由《空间的极限》
引理 \ref{spaces-limits-lemma-extend}，存在有限型拟凝聚理想层
$\mathcal{J} \subset \mathcal{O}_X$，其在 $U \cup V$ 上的限制为
$\mathcal{I}'$。引理得证。
\end{proof}

\begin{lemma}
\label{spaces-divisors-lemma-dominate-admissible-blowups}
设 $S$ 为概形，$X$ 为 $S$ 上拟紧拟分离的代数空间，$U \subset X$ 为拟紧开子空间。
设
% P09-ERR-0239
$b_i : X_i \to X$（$i = 1, \ldots, n$）为 $U$-容许爆破。则存在
$U$-容许爆破 $b : X' \to X$，使得 (a) 对 $i = 1, \ldots, n$，$b$ 分解为
$X' \to X_i \to X$；且 (b) 每个态射 $X' \to X_i$ 都是 $U$-容许爆破。
\end{lemma}

\begin{proof}
设 $\mathcal{I}_i \subset \mathcal{O}_X$ 为有限型拟凝聚理想层，使得
$V(\mathcal{I}_i)$ 与 $U$ 不交，且 $X_i$ 同构于 $X$ 在
$\mathcal{I}_i$ 处的爆破。令
$\mathcal{I} = \mathcal{I}_1 \cdot \ldots \cdot \mathcal{I}_n$，并令 $X'$ 为
$X$ 在 $\mathcal{I}$ 处的爆破。则由引理 \ref{spaces-divisors-lemma-blowing-up-two-ideals}，
$X' \to X$ 经由 $b_i$ 分解。
\end{proof}

\begin{lemma}
\label{spaces-divisors-lemma-separate-disjoint-opens-by-blowing-up}
设 $S$ 为概形，$X$ 为 $S$ 上拟紧拟分离的代数空间，$U, V$ 为 $X$ 的两个不交
拟紧开子空间。则存在 $U \cup V$-容许爆破 $b : X' \to X$，使得 $X'$ 是开子空间
的不交并 $X' = X'_1 \amalg X'_2$，并且 $b^{-1}(U) \subset X'_1$、
$b^{-1}(V) \subset X'_2$。
\end{lemma}

\begin{proof}
选取有限型拟凝聚理想层 $\mathcal{I}$ 与 $\mathcal{J}$，分别使得
$X \setminus U = V(\mathcal{I})$ 与 $X \setminus V = V(\mathcal{J})$，见
《空间的极限》引理
\ref{spaces-limits-lemma-quasi-coherent-finite-type-ideals}。于是
$|V(\mathcal{I}\mathcal{J})| = |X|$，故 $\mathcal{I}\mathcal{J}$ 是局部幂零
理想层。由于 $\mathcal{I}$ 与 $\mathcal{J}$ 是有限型的，且 $X$ 拟紧，存在
$n > 0$ 使得 $\mathcal{I}^n \mathcal{J}^n = 0$。可以并且确实用
$\mathcal{I}^n$ 代替 $\mathcal{I}$，用 $\mathcal{J}^n$ 代替 $\mathcal{J}$。
于是 $\mathcal{I} \mathcal{J} = 0$。设 $b : X' \to X$ 为在
$\mathcal{I} + \mathcal{J}$ 处的爆破。由于
% P09-ERR-0240
$|V(\mathcal{I} + \mathcal{J})| = |X| \setminus |U| \cup |V|$，它是
$U \cup V$-容许的。我们将证明，如引理陈述中那样，$X'$ 是开子空间的不交并
$X' = X'_1 \amalg X'_2$。

\medskip\noindent
由于 $|V(\mathcal{I} + \mathcal{J})|$ 是 $|U \cup V|$ 的补集，可知
$V \cup U$ 在 $X'$ 中概形论稠密，见引理
\ref{spaces-divisors-lemma-blowing-up-gives-effective-Cartier-divisor} 及
\ref{spaces-divisors-lemma-complement-effective-Cartier-divisor}。因此，若存在这样的分解
$X' = X'_1 \amalg X'_2$ 为开闭子空间，则 $X'_1$ 是 $U$ 在 $X'$ 中的概形论闭包，
类似地，$X'_2$ 是 $V$ 在 $X'$ 中的概形论闭包。由于 $U \to X'$ 与
$V \to X'$ 拟紧，取概形论闭包与平展局部化可交换（《空间的态射》引理
\ref{spaces-morphisms-lemma-quasi-compact-scheme-theoretic-image}）。因此，为验证
$X'_1$ 与 $X'_2$ 的存在性，可以在 $X$ 上平展局部地进行。这把问题归约到概形
情形；该情形已在《除子》引理
\ref{divisors-lemma-separate-disjoint-opens-by-blowing-up} 的证明中处理。
\end{proof}
